\documentclass[11pt, a4paper]{article}

\usepackage[english]{estilo-mate}
\tema{Linear Algebra}

\title{\textbf{Linear Algebra Notes}}
\author{Lucas Allara}
\date{}

\begin{document}

\maketitle
\newpage

\tableofcontents
\newpage

\section{Vector Spaces and Subspaces}

\subsection{Vector Spaces}

\begin{definicion}{$\mathbb{K}$-Vector Space}{espacio_vectorial}
Let $\mathbb{K}$ be a field. A nonempty set $V$ is said to be a $\mathbb{K}$-vector space (or vector space over $\mathbb{K}$) if it is equipped with two operations:
\begin{enumerate}
    \item[i)] Addition: $+: V \times V \to V$
    \item[ii)] Scalar multiplication: $\cdot: \mathbb{K} \times V \to V$
\end{enumerate}
such that, for all $u, v, w \in V$ and $\alpha, \beta \in \mathbb{K}$, the following conditions hold:

Properties of addition (Abelian Group Structure):
\begin{itemize}
    \item Commutativity: $u+v = v+u$.
    \item Associativity: $(u+v)+w = u+(v+w)$.
    \item Existence of a neutral element: there is $0_V \in V$ such that $v + 0_V = v$.
    \item Existence of an opposite: for each $v$ there is $-v$ such that $v + (-v) = 0_V$.
\end{itemize}

Properties of scalar multiplication:
\begin{itemize}
    \item Associativity: $\alpha(\beta v) = (\alpha\beta)v$.
    \item Identity of the field: $1 \cdot v = v$ (where $1$ is the unit of $\mathbb{K}$).
    \item Distributivity with respect to vector addition: $\alpha(u+v) = \alpha u + \alpha v$.
    \item Distributivity with respect to scalar addition: $(\alpha+\beta)v = \alpha v + \beta v$.
\end{itemize}
\end{definicion}

\obs By a usual abuse of notation, from now on $V$ is said to be a $\mathbb{K}$-vector space, referring to the triple $(V, +, \cdot)$ of Definition \ref{def:espacio_vectorial}. The space is identified with its underlying set and the operations of addition and scalar multiplication are understood, unless the context requires specifying them.

\subsection{Vectors and Operations}

\begin{definicion}{Vector in the General Sense}{vector_general}
Let $V$ be a $\mathbb{K}$-vector space. Any element $v \in V$ is called a \textbf{vector}.
\end{definicion}

\obs What is essential about a vector space is not the nature of its elements (polynomials, matrices, functions), but their behavior under the addition and scalar multiplication defined in Definition \ref{def:espacio_vectorial}.

From the perspective of linear algebra, if two spaces share the same operational structure, they are essentially the same object. Distinguishing a $2\times 2$ matrix from a vector of $\mathbb{R}^4$ based only on how they are written would be as artificial as claiming that the same word changes meaning depending on whether it is written in italics or in print. The representation changes, but the entity is the same.

\begin{ejemplo}{Matrices}{ej_matrices}
Let $M_{2\times 2}$ be the set of $2 \times 2$ matrices. A vector in this space is a matrix:
\[
A = \begin{pmatrix} a & b \\ c & d \end{pmatrix}.
\]
Here vector addition is matrix addition and scalar multiplication is performed entry by entry.
\end{ejemplo}

\begin{ejemplo}{The Space $\mathbb{R}^n$}{ej_rn_general}
Let $\mathbb{R}^n$ be the set of all ordered lists of $n$ real numbers.
A vector in this space has the form:
\[
\vec{v} = (v_1, v_2, \dots, v_n), \quad \text{where } v_i \in \mathbb{R}.
\]
Here the operations are defined coordinate by coordinate:
\begin{itemize}
    \item Addition: if $\vec{u} = (u_1, \dots, u_n)$ and $\vec{v} = (v_1, \dots, v_n)$,
    \[
    \vec{u} + \vec{v} = (u_1 + v_1, \ u_2 + v_2, \ \dots, \ u_n + v_n).
    \]
    \item Scalar multiplication: if $\alpha \in \mathbb{R}$,
    \[
    \alpha \vec{v} = (\alpha v_1, \ \alpha v_2, \ \dots, \ \alpha v_n).
    \]
\end{itemize}
The zero vector in this space is the list of zeros: $\vec{0} = (0, 0, \dots, 0)$.
\end{ejemplo}

\begin{ejemplo}{Polynomials}{ej_polinomios}
Let $P_n$ be the set of polynomials of degree less than or equal to $n$ with real coefficients.
A vector in this space is a polynomial of the form:
\[
p(x) = a_0 + a_1 x + a_2 x^2 + \dots + a_n x^n.
\]
\end{ejemplo}

\subsection{Vector Subspaces}

\begin{definicion}{Vector Subspace}{subespacio}
Let $V$ be a $\mathbb{K}$-vector space. A \textbf{vector subspace} is a subset $W \subseteq V$ such that:
\begin{enumerate}
    \item $0 \in W$,
    \item If $u,v \in W$, then $u+v \in W$,
    \item If $u \in W$ and $\alpha \in \mathbb{K}$, then $\alpha u \in W$.
\end{enumerate}
\end{definicion}

\subsection{Linear Combinations and Generated Subspace}

\begin{definicion}{Linear Combination}{comb_lineal}
Let $v_1,v_2,\dots,v_k \in V$. A \textbf{linear combination} of these vectors is a vector of the form
\[
\alpha_1 v_1 + \alpha_2 v_2 + \cdots + \alpha_k v_k,
\quad \alpha_i \in \mathbb{K}.
\]
\end{definicion}

\begin{definicion}{Generated Subspace}{span}
The set of all linear combinations of $v_1,v_2,\dots,v_k$ is called the \textbf{span} or \textbf{generated subspace} and is denoted
\[
\operatorname{span}\{v_1,v_2,\dots,v_k\}
= \{ \alpha_1 v_1 + \cdots + \alpha_k v_k \mid \alpha_i \in \mathbb{K} \}.
\]
The vectors $v_1,\dots,v_k$ are called \textbf{generators} of this subspace.
\end{definicion}

\begin{proposicion}{The Span is a Subspace}{span_es_subespacio}
Let $V$ be a $\mathbb{K}$-vector space and $v_1,\dots,v_k \in V$. Then $\operatorname{span}\{v_1,\dots,v_k\}$ is a subspace of $V$.
\end{proposicion}

\begin{proof}
The three subspace conditions (Definition \ref{def:subespacio}) are verified:
\begin{enumerate}
    \item the zero vector belongs to the span, since with $\alpha_i = 0$ one obtains $0 = 0v_1 + \cdots + 0v_k$;
    \item given $u = \sum_{i=1}^k \alpha_i v_i$ and $w = \sum_{i=1}^k \beta_i v_i$ in the span,
    \[
    u+w = \sum_{i=1}^k (\alpha_i+\beta_i) v_i
    \]
    is another linear combination, so that $u+w \in \operatorname{span}\{v_1,\dots,v_k\}$;
    \item given $u = \sum_{i=1}^k \alpha_i v_i$ in the span and $\lambda \in \mathbb{K}$,
    \[
    \lambda u = \sum_{i=1}^k (\lambda \alpha_i) v_i
    \]
    is also a linear combination, hence $\lambda u \in \operatorname{span}\{v_1,\dots,v_k\}$.
\end{enumerate}
Therefore the span is a subspace of $V$.
\end{proof}

\subsection{Representation of Subspaces}

\begin{enumerate}
    \item \textit{Parametric form.} $S$ is described by how it is built (through the span,
    Definition \ref{def:span}): a set of vectors is given and $S$ is the set of all their
    linear combinations.
    \[ S = \operatorname{span}\{u_1, \dots, u_k\} = \{ \alpha_1 u_1 + \dots + \alpha_k u_k \mid \alpha_i \in \mathbb{K} \}. \]
    With this form it is direct to exhibit points of the subspace (by giving values to the $\alpha_i$), but it is hard to decide whether a given point belongs to it.

    \item \textit{Implicit form.} $S$ is described by the conditions it must satisfy: a
    set of restrictions (homogeneous linear equations) is given and $S$ is the set where all of them
    are satisfied.
    \[ S = \{ x \in V \mid Ax = 0 \}. \]
    With this form it is direct to decide whether a point belongs to the subspace (by checking the restrictions), but it is hard to exhibit its points.
\end{enumerate}

\subsection{Intersection and Sum of Subspaces}

\begin{definicion}{Intersection of Subspaces}{interseccion}
Let $S$ and $T$ be subspaces of a $\mathbb{K}$-vector space $V$. The intersection of $S$ and $T$ is defined as:
\[
S \cap T = \{ v \in V \mid v \in S \text{ and } v \in T \}.
\]
\end{definicion}

\begin{proposicion}{Intersection of Subspaces}{prop_interseccion}
The intersection $S \cap T$ is a subspace of $V$, it is contained in both, and it is the largest subspace with this property.
\end{proposicion}

\begin{proof}
\textit{$S \cap T$ is a subspace} (Definition \ref{def:subespacio}):
\begin{itemize}
    \item since $S$ and $T$ are subspaces, $0 \in S$ and $0 \in T$, hence $0 \in S \cap T$;
    \item if $u, v \in S \cap T$, then $u, v \in S$ and $u, v \in T$; since both are subspaces, $u+v \in S$ and $u+v \in T$, so that $u+v \in S \cap T$;
    \item if $u \in S \cap T$ and $\alpha \in \mathbb{K}$, then $u \in S$ and $u \in T$, and by closure $\alpha u \in S$ and $\alpha u \in T$, hence $\alpha u \in S \cap T$.
\end{itemize}

\textit{It is contained in $S$ and in $T$:} by definition of intersection, if $x \in S \cap T$ then $x \in S$ and $x \in T$, that is, $S \cap T \subseteq S$ and $S \cap T \subseteq T$.

\textit{It is the largest subspace contained in $S$ and $T$:} let $W$ be a subspace with $W \subseteq S$ and $W \subseteq T$; it must be seen that $W \subseteq S \cap T$. Given $w \in W$, since $W \subseteq S$ one has $w \in S$, and since $W \subseteq T$, $w \in T$; being in both, $w \in S \cap T$. Hence $W \subseteq S \cap T$.
\end{proof}

\begin{definicion}{Sum of Subspaces}{suma_subespacios}
Let $S$ and $T$ be subspaces of a $\mathbb{K}$-vector space $V$. The sum of $S$ and $T$ is defined as the smallest subspace of $V$ containing $S$ and $T$:
\[
S + T = \bigcap \{ W \subseteq V \mid W \text{ is a subspace of $V$ and } S \subseteq W, \ T \subseteq W \}.
\]
\end{definicion}

\begin{proposicion}{Equivalence of the Sum}{suma_equiv}
The following definition:
\[
S+T = \{ s+t \mid s \in S, t \in T \}
\]
is equivalent to the previous definition.
\end{proposicion}

\begin{proof}
Let $W = \{ s+t \mid s \in S, t \in T \}$.

\textit{$W$ is a subspace} (Definition \ref{def:subespacio}):
\begin{itemize}
    \item with $s=0 \in S$ and $t=0 \in T$, $0 = 0+0 \in W$;
    \item if $w_1 = s_1+t_1$ and $w_2 = s_2+t_2$ are in $W$ (with $s_i \in S$, $t_i \in T$), then
    \[
    w_1 + w_2 = (s_1+s_2) + (t_1+t_2) \in W,
    \]
    since $s_1+s_2 \in S$ and $t_1+t_2 \in T$;
    \item for $\alpha \in \mathbb{K}$, $\alpha w = (\alpha s) + (\alpha t) \in W$, because $\alpha s \in S$ and $\alpha t \in T$.
\end{itemize}

\textit{It contains $S$ and $T$:} for $s \in S$, $s+0 \in W$, so that $S \subseteq W$; for $t \in T$, $0+t \in W$, so that $T \subseteq W$.

\textit{It is the smallest one containing $S$ and $T$:} let $W'$ be a subspace with $S \subseteq W'$ and $T \subseteq W'$. For $w = s+t \in W$ with $s \in S$, $t \in T$, one has $s, t \in W'$, and since $W'$ is a subspace, $s+t \in W'$. Hence $W \subseteq W'$.
\end{proof}

\subsection{Direct Sum and Unique Writing}

\begin{definicion}{Uniqueness of the Writing}{unicidad_general}
Let $v$ be a vector. The writing of $v$ is \textbf{unique} if any two representations of $v$ coincide component by component.
\end{definicion}

\begin{ejemplo}{Application to Linear Combinations}{}
That $v$ is written uniquely as a linear combination of $\{u_1, \dots, u_n\}$ means that if
\[
v = \sum \alpha_i u_i \quad \text{and} \quad v = \sum \beta_i u_i,
\]
then necessarily $\alpha_i = \beta_i$ for every $i$.
\end{ejemplo}

\begin{ejemplo}{Application to Sums of Subspaces}{unicidad_suma}
That a vector $v \in S+T$ is written uniquely as the sum of an element of $S$ and one of $T$ (Definition \ref{def:unicidad_general}) means that if two decompositions are considered:
\[
v = s + t \quad \text{and} \quad v = s' + t',
\]
with $s, s' \in S$ and $t, t' \in T$, the uniqueness condition necessarily implies that:
\[
s = s' \quad \text{and} \quad t = t'.
\]
\end{ejemplo}

\begin{definicion}{Direct Sum}{suma_directa}
If every vector $v \in S+T$ has a unique writing, $S+T$ is the \textbf{direct sum} of $S$ and $T$, and one writes
\[
S \oplus T = S+T.
\]
\end{definicion}

\begin{proposicion}{Characterization of the Direct Sum}{suma_directa_carac}
Let $S$ and $T$ be subspaces of $V$. The sum $S+T$ is direct if and only if
\[
S \cap T = \{0\}.
\]
\end{proposicion}

\begin{proof}
Both implications are proved.

\textit{($\Rightarrow$).} Suppose $S \oplus T$ and let $v \in S \cap T$. Then $v \in S$ and $v \in T$, and
\[
v = v + 0 = 0 + v,
\]
with $v, 0 \in S$ and $v, 0 \in T$. Since the writing is unique (Definition \ref{def:suma_directa}), $v = 0$. Hence $S \cap T = \{0\}$.

\textit{($\Leftarrow$).} Suppose $S \cap T = \{0\}$ and let $v \in S+T$ with two writings
\[
v = s_1 + t_1 = s_2 + t_2, \quad s_i \in S, \ t_i \in T.
\]
Subtracting, $0 = (s_1 - s_2) + (t_1 - t_2)$, with $s_1 - s_2 \in S$ and $t_1 - t_2 \in T$. Since $S \cap T = \{0\}$,
\[
s_1 - s_2 = 0, \quad t_1 - t_2 = 0,
\]
that is, $s_1 = s_2$ and $t_1 = t_2$. Therefore the writing of $v$ is unique and $S+T$ is a direct sum.
\end{proof}

\subsection{Inner Product}

\begin{definicion}{Inner Product}{producto_interno_general}
Let $V$ be a $\mathbb{K}$-vector space (where $\mathbb{K}$ is usually $\mathbb{R}$ or $\mathbb{C}$). An \textbf{inner product} on $V$ is a function
\[
\langle \cdot , \cdot \rangle : V \times V \to \mathbb{K}
\]
that satisfies the following properties for all $x, y, z \in V$ and $\alpha \in \mathbb{K}$:

\begin{enumerate}
    \item Conjugate (or Hermitian) Symmetry:
    \[
    \langle x, y \rangle = \overline{\langle y, x \rangle}.
    \]
    \item Linearity in the first variable:
    \[
    \langle x + y, z \rangle = \langle x, z \rangle + \langle y, z \rangle,
    \quad
    \langle \alpha x, y \rangle = \alpha \langle x, y \rangle.
    \]
    \item Positive definiteness:
    \[
    \langle x, x \rangle \in \mathbb{R}_{\geq 0} \quad (\text{it is a nonnegative real number}),
    \]
    and moreover $\langle x, x \rangle = 0 \iff x = 0$.
\end{enumerate}
\end{definicion}

\obs If $\mathbb{K} = \mathbb{C}$, property 2 implies that in the second variable the scalars come out conjugated:
\[
\langle x, \alpha y \rangle = \overline{\langle \alpha y, x \rangle} = \overline{\alpha \langle y, x \rangle} = \overline{\alpha} \overline{\langle y, x \rangle} = \overline{\alpha} \langle x, y \rangle.
\]
If $\mathbb{K} = \mathbb{R}$, conjugation is the identity ($\bar{a} = a$), and symmetry and bilinearity are recovered.

\obs The usual inner product on $\mathbb{R}^n$ can be seen as a matrix multiplication: representing $\vec{x}$ as a row vector and $\vec{y}$ as a column vector,
\[
\langle \vec{x}, \vec{y} \rangle = 
\begin{bmatrix} x_1 & x_2 & \cdots & x_n \end{bmatrix}
\begin{bmatrix} y_1 \\ y_2 \\ \vdots \\ y_n \end{bmatrix}.
\]

\subsection{Norm and Cauchy-Schwarz}

\begin{definicion}{Induced Norm}{norma}
Let $(V,\langle\cdot,\cdot\rangle)$ be an inner product space. The \textbf{induced norm} of a vector $x \in V$ is:
\[
\|x\| := \sqrt{\langle x,x\rangle}.
\]
\end{definicion}

\begin{proposicion}{Cauchy-Schwarz Inequality}{cauchy_schwarz}
In every inner product space $(V,\langle\cdot,\cdot\rangle)$
\[
\;|\langle x,y\rangle| \le \|x\|\,\|y\|\; \qquad\forall x,y\in V.
\]
\end{proposicion}

\begin{proof}
The case $\mathbb{K} = \mathbb{R}$ is proved; the complex case is reduced to it in the remark
that follows. The idea is to use that the quadratic $\varphi(t) = \langle tx+y, tx+y\rangle$ is nonnegative. Let
$x, y \in V$; for every $t \in \mathbb{R}$ one has $\varphi(t) \ge 0$, and by bilinearity
\[
\varphi(t) = t^2\langle x,x\rangle + 2t\langle x,y\rangle + \langle y,y\rangle,
\]
a quadratic $a t^2 + b t + c$ with $a = \|x\|^2$, $b = 2\langle x,y\rangle$, $c = \|y\|^2$. Since
$\varphi(t)\ge 0$, its discriminant satisfies
\[
\Delta = b^2 - 4ac = 4\Big(\langle x,y\rangle^2 - \|x\|^2\|y\|^2\Big) \le 0,
\]
whence $\langle x,y\rangle^2 \le \|x\|^2\|y\|^2$ and, taking square roots, $|\langle x,y\rangle| \le \|x\|\,\|y\|$.
\end{proof}

\obs
If the field is $\mathbb{C}$, the polynomial cannot be formed directly with $\|x+ty\|^2$ because the middle term would be complex and $\varphi$ would not be a real quadratic in $t$. The problem is transformed to the real case as follows.

Suppose $\langle x, y \rangle \neq 0$ (if it is $0$, the inequality holds directly). This complex number is written in polar form:
\[
\langle x, y \rangle = |\langle x, y \rangle| e^{i\theta}.
\]
Define an auxiliary vector $\tilde{y}$ by rotating $y$ to cancel that phase:
\[
\tilde{y} = e^{i\theta} y.
\]
Two things are observed about $\tilde{y}$:
\begin{enumerate}
    \item It has the same norm as $y$: $\|\tilde{y}\| = |e^{i\theta}| \|y\| = 1 \cdot \|y\| = \|y\|$.
    \item Its inner product with $x$ is now a real number:
    \[
    \langle x, \tilde{y} \rangle = \langle x, e^{i\theta} y \rangle = \overline{e^{i\theta}} \langle x, y \rangle = e^{-i\theta} \left( |\langle x, y \rangle| e^{i\theta} \right) = |\langle x, y \rangle|.
    \]
\end{enumerate}

Now consider the quadratic polynomial in the real variable $t \in \mathbb{R}$:
\[
\varphi(t) = \| x + t\tilde{y} \|^2 \ge 0.
\]
Expanding, using that $\langle x, \tilde{y} \rangle$ is real,
\begin{align*}
\| x + t\tilde{y} \|^2 &= \langle x + t\tilde{y}, x + t\tilde{y} \rangle \\
&= \langle x, x \rangle + t \langle x, \tilde{y} \rangle + t \langle \tilde{y}, x \rangle + t^2 \langle \tilde{y}, \tilde{y} \rangle \\
&= \|x\|^2 + t |\langle x, y \rangle| + t \overline{|\langle x, y \rangle|} + t^2 \|\tilde{y}\|^2 \\
&= \|x\|^2 + 2t |\langle x, y \rangle| + t^2 \|y\|^2.
\end{align*}
From here the proof continues as in the real case, and since $\langle x, \tilde{y} \rangle = |\langle x, y \rangle|$ and $\|\tilde{y}\| = \|y\|$, the conclusion is the desired inequality.

\begin{proposicion}{Properties of the Induced Norm}{prop_norma}
The induced norm satisfies the following properties for all $x, y \in V$ and $\alpha \in \mathbb{K}$:
\begin{enumerate}
    \item Positivity: $\|x\| \ge 0$, and $\|x\|=0 \iff x=0$.
    \item Homogeneity: $\|\alpha x\| = |\alpha|\,\|x\|$.
    \item Triangle Inequality: $\|x+y\| \le \|x\| + \|y\|$.
\end{enumerate}
That is, the induced norm is a norm.
\end{proposicion}

\begin{proof}
The first two are verified from the definition of inner product (Definition \ref{def:producto_interno_general}); the last one follows from the Cauchy-Schwarz inequality (Proposition \ref{prop:cauchy_schwarz}).
\end{proof}

\begin{proposicion}{Equality Condition in Cauchy-Schwarz}{igualdad_cs}
Let $x, y$ be vectors in an inner product space. The equality
\[
|\langle x, y \rangle| = \|x\| \|y\|
\]
holds if and only if one is a scalar multiple of the other.
\end{proposicion}

\begin{proof}
Both implications are proved.

\textit{($\Rightarrow$).} Suppose that $|\langle x, y \rangle| = \|x\| \|y\|$. If $x = 0$, then
$x$ is a multiple of $y$ and there is nothing to prove; suppose $x \neq 0$ and consider the vector $z$:
\[
z = y - \frac{\langle y, x \rangle}{\|x\|^2} x.
\]
Its squared norm, expanding $\langle z, z \rangle$, is:
\[
\|z\|^2 = \left\langle y - \frac{\langle y, x \rangle}{\|x\|^2} x, \ y - \frac{\langle y, x \rangle}{\|x\|^2} x \right\rangle.
\]
On expanding (in the complex case, $\langle \alpha u, v \rangle = \alpha \langle u, v \rangle$ and $\langle u, \alpha v \rangle = \bar{\alpha} \langle u, v \rangle$):
\[
\|z\|^2 = \|y\|^2 - \frac{\overline{\langle y, x \rangle}}{\|x\|^2}\langle y, x \rangle - \frac{\langle y, x \rangle}{\|x\|^2}\langle x, y \rangle + \left| \frac{\langle y, x \rangle}{\|x\|^2} \right|^2 \|x\|^2.
\]
Since $\overline{\langle y, x \rangle} = \langle x, y \rangle$, the middle terms combine:
\[
\|z\|^2 = \|y\|^2 - \frac{|\langle x, y \rangle|^2}{\|x\|^2} - \frac{|\langle x, y \rangle|^2}{\|x\|^2} + \frac{|\langle x, y \rangle|^2}{\|x\|^4} \|x\|^2.
\]
Simplifying:
\[
\|z\|^2 = \|y\|^2 - 2\frac{|\langle x, y \rangle|^2}{\|x\|^2} + \frac{|\langle x, y \rangle|^2}{\|x\|^2} = \|y\|^2 - \frac{|\langle x, y \rangle|^2}{\|x\|^2}.
\]
By hypothesis $|\langle x, y \rangle|^2 = \|x\|^2 \|y\|^2$, whence:
\[
\|z\|^2 = \|y\|^2 - \frac{\|x\|^2 \|y\|^2}{\|x\|^2} = \|y\|^2 - \|y\|^2 = 0.
\]
Since $\|z\|^2 = 0$ implies $z = 0$ (positivity of the norm, Proposition \ref{prop:prop_norma}), one obtains:
\[
y - \frac{\langle y, x \rangle}{\|x\|^2} x = 0 \implies y = \left( \frac{\langle y, x \rangle}{\|x\|^2} \right) x.
\]
Therefore $y$ is a scalar multiple of $x$.

\textit{($\Leftarrow$).} There is $\lambda \in \mathbb{K}$ with $y = \lambda x$. Substituting on both sides:
\begin{itemize}
    \item Left-hand side:
    \[
    |\langle x, y \rangle| = |\langle x, \lambda x \rangle| = |\bar{\lambda} \langle x, x \rangle| = |\lambda| \|x\|^2.
    \]
    \item Right-hand side:
    \[
    \|x\| \|y\| = \|x\| \|\lambda x\| = \|x\| (|\lambda| \|x\|) = |\lambda| \|x\|^2.
    \]
\end{itemize}
Both sides coincide.
\end{proof}

\subsection{Distance and Angles}

The existence of an inner product endows the space with a geometric structure. From it, metric concepts such as the length of vectors (norm) and the distance between them are induced.

There is a hierarchy: an inner product induces a norm, and a norm induces a distance; the converse path does not always hold (there are distances that do not come from norms, and norms that do not come from inner products).

\begin{definicion}{Induced Distance}{distancia}
Let $(V,\|.\|)$ be a normed space. The \textbf{induced distance} between two vectors $x, y \in V$ is:
\[
d(x,y) := \|x-y\|.
\]
\end{definicion}

\begin{proposicion}{Properties of the Induced Distance}{prop_distancia}
The function $d$ defines a metric on $V$, that is:
\begin{enumerate}
    \item $d(x,y) \ge 0$ and $d(x,y)=0 \iff x=y$.
    \item Symmetry: $d(x,y) = d(y,x)$.
    \item Triangle Inequality: $d(x,z) \le d(x,y) + d(y,z)$.
\end{enumerate}
\end{proposicion}

\begin{proof}
The three properties are verified from those of the norm (Proposition \ref{prop:prop_norma}).
\end{proof}

\begin{definicion}{Angle between Vectors}{angulo}
Let $x, y \in V$ be nonzero vectors. The \textbf{angle} $\theta$ between them is the unique real number such that:
\[
\cos\theta = \frac{\langle x,y\rangle}{\|x\|\cdot\|y\|}, \qquad 0 \le \theta \le \pi.
\]
\end{definicion}

\obs This definition makes sense thanks to the Cauchy-Schwarz inequality (Proposition \ref{prop:cauchy_schwarz}).
One has $|\langle x, y \rangle| \le \|x\| \|y\|$; dividing by the product of the norms,
\[
\frac{|\langle x, y \rangle|}{\|x\| \|y\|} \le 1 \implies -1 \le \frac{\langle x, y \rangle}{\|x\| \|y\|} \le 1.
\]
Since the value always falls in the interval $[-1, 1]$, it is within the domain of the arccosine function, which guarantees that the angle $\theta$ is well defined.

\newpage

\section{Orthogonal Spaces, Linear Independence, Basis and Dimension}

\subsection{Orthogonality and Orthogonal Complement}

\begin{definicion}{Orthogonality}{ortogonalidad}
A vector $x$ is \textbf{orthogonal} to $y$ ($x \perp y$) if $\langle x, y \rangle = 0$.
\end{definicion}

\begin{definicion}{Orthogonal Subspaces}{sub_ortogonales}
Let $V$ be a $\mathbb{K}$-vector space with an inner product and let $S$ and $T$ be subspaces of $V$.
$S$ and $T$ are \textbf{orthogonal} ($S \perp T$) if:
\[
\langle x, y \rangle = 0, \quad \forall x \in S, \ \forall y \in T.
\]
\end{definicion}

\begin{proposicion}{Trivial Intersection}{int_trivial}
If $S \perp T$, then $S \cap T = \{0\}$.
\end{proposicion}

\begin{proof}
Let $z \in S \cap T$. Since $z \in S$, $z \in T$ and $S \perp T$, one has $\langle z, z \rangle = 0$, hence $z = 0$.
\end{proof}

\begin{definicion}{Orthogonal Complement}{complemento_ortogonal}
Let $V$ be a $\mathbb{K}$-vector space with an inner product and let $S$ be a subspace of $V$. The \textbf{orthogonal complement} of $S$, denoted $S^\perp$, is:
\[
S^\perp = \{ x \in V \mid \langle x, s \rangle = 0, \ \forall s \in S \}.
\]
\end{definicion}

\begin{proposicion}{Properties of the Orthogonal Complement}{prop_complemento}
Let $V$ be a $\mathbb{K}$-vector space with an inner product and let $S$ be a subspace of $V$. The orthogonal complement $S^\perp$ is a subspace of $V$, it is orthogonal to $S$, and it is the largest subspace with this property.
\end{proposicion}

\begin{proof}
\textit{$S^\perp$ is a subspace:}
\begin{itemize}
    \item the zero vector belongs to $S^\perp$, since $\langle 0, s \rangle = 0$ for every $s \in S$;
    \item if $x, y \in S^\perp$, for every $s \in S$ one has $\langle x+y, s \rangle = \langle x, s \rangle + \langle y, s \rangle = 0$, hence $x+y \in S^\perp$;
    \item if $x \in S^\perp$ and $\alpha \in \mathbb{K}$, $\langle \alpha x, s \rangle = \alpha \langle x, s \rangle = 0$, hence $\alpha x \in S^\perp$.
\end{itemize}

\textit{$S^\perp \perp S$:} by definition of orthogonal complement (Definition \ref{def:complemento_ortogonal}), $\langle x, s \rangle = 0$ for every $x \in S^\perp$ and every $s \in S$.

\textit{It is the largest with that property:} let $W$ be a subspace with $W \perp S$. Given $w \in W$, since $W \perp S$ one has $\langle w, s \rangle = 0$ for every $s \in S$, so that $w \in S^\perp$. Hence $W \subseteq S^\perp$.
\end{proof}

\subsection{Projection onto Subspaces and Complementarity}

\begin{definicion}{Finite Dimension}{dim_finita}
A $\mathbb{K}$-vector space $V$ is said to be \textbf{finite-dimensional} (or finitely generated) if it has a finite generating set, that is, if there is a finite subset $G \subset V$ such that:
\[ \operatorname{span}(G) = V \]
Otherwise, it is said to be \textbf{infinite-dimensional}.
\end{definicion}

The following theorem is a classical result of analysis and is stated without proof.

\begin{teorema}{Best Approximation in Subspaces}{mejor_aprox}
Let $V$ be a finite-dimensional $\mathbb{K}$-vector space with an inner product and let $S$ be a subspace of $V$.
For every vector $v \in V$, there is a unique element $p \in S$ such that:
\[ \|v - p\| \leq \|v - s\| \quad \text{for every } s \in S. \]
That is, $p$ is the vector of $S$ closest to $v$.
\end{teorema}

\begin{definicion}{Orthogonal Projection}{proyeccion_ortogonal}
The vector $p \in S$ of Theorem \ref{teo:mejor_aprox} is called the \textbf{orthogonal projection} of $v$ onto $S$ and is denoted $P_S(v)$.
\end{definicion}

\begin{teorema}{Orthogonal Decomposition}{descomposicion_ortogonal}
Let $V$ be a finite-dimensional $\mathbb{K}$-vector space with an inner product and $S$ a subspace of $V$. Then:
\[
V = S \oplus S^\perp
\]
This means that every vector $v \in V$ can be written uniquely as the sum of two orthogonal components:
\[ v = p + q, \quad \text{where } p \in S \text{ and } q \in S^\perp. \]
Specifically, $p = P_S(v)$ and $q = v - P_S(v)$.
\end{teorema}

\begin{proof}
Let $p \in S$ be the element that minimizes the distance to $v$ (Theorem \ref{teo:mejor_aprox}) and $q = v - p$. It must be proved that $q \in S^\perp$, that is, $\langle q, y \rangle = 0$ for every $y \in S$.

Fix $y \in S$. For every $\alpha \in \mathbb{R}$, the vector $p + \alpha y$ belongs to $S$, and the function
\[
f(\alpha) = \|v - (p + \alpha y)\|^2 = \|q - \alpha y\|^2
\]
has, by the minimality of $p$ (Theorem \ref{teo:mejor_aprox}), an absolute minimum at $\alpha = 0$. Expanding,
\[
f(\alpha) = \langle q - \alpha y, q - \alpha y \rangle = \|q\|^2 - 2\alpha \operatorname{Re}\langle q, y \rangle + \alpha^2 \|y\|^2.
\]
Since $f$ is differentiable and has a minimum at $0$, $f'(0) = -2 \operatorname{Re}\langle q, y \rangle = 0$, so that $\operatorname{Re}\langle q, y \rangle = 0$. In the real case this gives $\langle q, y \rangle = 0$; in the complex case, the same argument applied to $iy$ (following Remark) gives $\langle q, y \rangle = 0$. Hence $q \perp S$ and $v = p + q$.

Finally, since $S$ and $S^\perp$ are orthogonal, $S \cap S^\perp = \{0\}$ (Proposition \ref{prop:int_trivial}), so that the sum is direct and the writing is unique.
\end{proof}

\obs In the previous proof, with $\alpha \in \mathbb{R}$ it is proved that the real part of the inner product vanishes. For the imaginary part, in a complex space:
\begin{itemize}
    \item since $S$ is a subspace, $y \in S \implies iy \in S$;
    \item applying the previous result to $iy$, $\operatorname{Re}\langle q, iy \rangle = 0$;
    \item since the scalars in the second variable come out conjugated, $\langle q, iy \rangle = -i \langle q, y \rangle$;
    \item then $0 = \operatorname{Re}(-i\langle q, y \rangle) = \operatorname{Im}(\langle q, y \rangle)$.
\end{itemize}
Since the real and imaginary parts are zero, the inner product is zero.

\begin{corolario}{Double Orthogonal}{doble_orto}
$(S^\perp)^\perp = S$.
\end{corolario}

\begin{proof}
Double inclusion is proved.

\textit{($\supseteq$).} Let $s \in S$. For every $z \in S^\perp$, $\langle s, z \rangle = 0$, hence $s \in (S^\perp)^\perp$.

\textit{($\subseteq$).} Let $x \in (S^\perp)^\perp$ and write $x = s + z$ with $s \in S$, $z \in S^\perp$. Then $\langle x, z \rangle = \langle s, z \rangle + \langle z, z \rangle$; since $x \in (S^\perp)^\perp$ and $z \in S^\perp$, $\langle x, z \rangle = 0$, and since $s \in S$, $\langle s, z \rangle = 0$. This leaves $0 = \|z\|^2$, so that $z = 0$ and $x = s \in S$.
\end{proof}

\begin{proposicion}{Inclusion Relations and Operations}{relaciones_orto}
Let $V$ be a finite-dimensional $\mathbb{K}$-vector space with an inner product and let $S$ and $T$ be subspaces of $V$. Then:
\begin{enumerate}
    \item $S \subset T \iff T^\perp \subset S^\perp$.
    \item $(S+T)^\perp = S^\perp \cap T^\perp$.
    \item $(S \cap T)^\perp = S^\perp + T^\perp$.
\end{enumerate}
\end{proposicion}

\begin{proof}
\textit{(1)} Both implications are proved.

\textit{($\Rightarrow$).} Let $x \in T^\perp$, that is, $\langle x, t \rangle = 0$ for every $t \in T$. Since $S \subset T$, in particular $\langle x, s \rangle = 0$ for every $s \in S$, so that $x \in S^\perp$. Hence $T^\perp \subset S^\perp$.

\textit{($\Leftarrow$).} Applying the forward direction to $T^\perp \subset S^\perp$ one obtains $(S^\perp)^\perp \subseteq (T^\perp)^\perp$, and by the double orthogonal complement (Corollary \ref{cor:doble_orto}), $S \subset T$.

\textit{(2)} Double inclusion. \textit{($\subseteq$)} if $x \in (S+T)^\perp$, then $\langle x, s+t \rangle = 0$; with $t=0$, $\langle x, s \rangle = 0$, so that $x \in S^\perp$, and with $s=0$, $x \in T^\perp$, hence $x \in S^\perp \cap T^\perp$. \textit{($\supseteq$)} if $x \in S^\perp \cap T^\perp$, then $\langle x, s \rangle = 0$ and $\langle x, t \rangle = 0$, so that $\langle x, s+t \rangle = 0$ and $x \in (S+T)^\perp$.

\textit{(3)} By (2), $(S^\perp + T^\perp)^\perp = S^{\perp\perp} \cap T^{\perp\perp} = S \cap T$; taking orthogonals on both sides, $S^\perp + T^\perp = (S \cap T)^\perp$.
\end{proof}

\subsection{Linear Independence}

\begin{proposicion}{Equivalent Forms of Linear Independence}{equiv_li}
Let $V$ be a $\mathbb{K}$-vector space and let $v_1, v_2, \dots, v_n$ be vectors in $V$. The following statements are equivalent:
\begin{enumerate}[label=(\roman*)] % roman labels so that they match the references (i)/(ii)/(iii) in the proof
    \item Each vector $x \in \operatorname{span}\{v_1, \dots, v_n\}$ is written uniquely as a linear combination of the $v_i$.
    \item No $v_j$ is a linear combination of the other vectors of the set.
    \item The only zero linear combination is the trivial one:
    \[
    \sum_{i=1}^{n} \alpha_i v_i = 0 \implies \alpha_1 = \alpha_2 = \dots = \alpha_n = 0.
    \]
\end{enumerate}
\end{proposicion}

\begin{proof}
The chain (i) $\implies$ (ii) $\implies$ (iii) $\implies$ (i) is proved.

\textit{(i) $\implies$ (ii).} By contradiction: if some $v_j$ were a linear combination of the others, $v_j = \sum_{i \neq j} \beta_i v_i$, there would be two different writings of $v_j$ (that one and $1 \cdot v_j$), against the uniqueness of (i).

\textit{(ii) $\implies$ (iii).} By contraposition: if there were a zero linear combination $\sum_{i} \alpha_i v_i = 0$ with some $\alpha_j \neq 0$, solving,
\[
v_j = \sum_{i \neq j} \left( -\frac{\alpha_i}{\alpha_j} \right) v_i,
\]
so that $v_j$ would be a combination of the others, against (ii).

\textit{(iii) $\implies$ (i).} Let $x$ be in the span with two writings $x = \sum_{i} \alpha_i v_i = \sum_{i} \beta_i v_i$. Subtracting,
\[
0 = \sum_{i=1}^{n} (\alpha_i - \beta_i) v_i,
\]
and by (iii) all the coefficients are zero, $\alpha_i = \beta_i$ for every $i$. Hence the writing is unique.
\end{proof}

\begin{definicion}{Linearly Independent Set}{def_li}
Let $V$ be a $\mathbb{K}$-vector space and let $v_1, v_2, \dots, v_n$ be vectors in $V$.
\begin{itemize}
    \item The set is \textbf{linearly independent} if it satisfies any of the equivalent conditions of Proposition \ref{prop:equiv_li}.
    \item Otherwise, it is \textbf{linearly dependent}.
\end{itemize}
\end{definicion}

\subsection{Basis and Dimension}

\begin{definicion}{Basis of a Space (Hamel)}{def_base}
Let $V$ be a $\mathbb{K}$-vector space. A set of vectors $\mathcal{B} \subset V$ is a \textbf{basis} of $V$ if it satisfies two conditions:
\begin{enumerate}
    \item Every vector $v \in V$ can be expressed as a finite linear combination of elements of $\mathcal{B}$.
    That is, there are $b_1, \dots, b_k \in \mathcal{B}$ and scalars $\alpha_1, \dots, \alpha_k$ such that:
    \[ v = \sum_{i=1}^k \alpha_i b_i. \]
    \item It is linearly independent.
\end{enumerate}
\end{definicion}

\obs The finiteness condition on the sum is essential: infinite sums of vectors are not defined, so that, although a basis may have infinitely many elements, each vector is built with finitely many of them.

\begin{proposicion}{Reduction to a Basis}{reduccion_base}
Let $A = \{v_1, \dots, v_k\}$ be a set of vectors that generates a $\mathbb{K}$-vector space $V$. If $A$ is linearly dependent, then there is a proper subset of $A$ that still generates $V$.
\end{proposicion}

\begin{proof}
Since $A$ is linearly dependent, by Proposition \ref{prop:equiv_li} some $v_j \in A$ is a linear combination of the others:
\[
v_j = \sum_{i \neq j} \alpha_i v_i.
\]
It must be seen that $A' = A \setminus \{v_j\}$ still generates $V$. The inclusion $\operatorname{span}(A') \subseteq \operatorname{span}(A)$ holds because $A' \subset A$. For the other one, let $x \in \operatorname{span}(A) = V$,
\[
x = \beta_j v_j + \sum_{i \neq j} \beta_i v_i;
\]
substituting $v_j$,
\[
x = \beta_j \left( \sum_{i \neq j} \alpha_i v_i \right) + \sum_{i \neq j} \beta_i v_i,
\]
a linear combination of the $v_i$ with $i \neq j$, so that $x \in \operatorname{span}(A')$. Hence $v_j$ can be removed without altering the span.
\end{proof}

\begin{proposicion}{Extension to a Basis}{extension_base}
Let $A = \{v_1, \dots, v_k\}$ be a linearly independent set contained in a $\mathbb{K}$-vector space $V$. If $A$ does not generate $V$, then a vector $v_{k+1} \in V$ can be added such that the new set $\{v_1, \dots, v_k, v_{k+1}\}$ is still linearly independent.
\end{proposicion}

\begin{proof}
Since $\operatorname{span}(A) \neq V$, there is $v_{k+1} \in V$ with $v_{k+1} \notin \operatorname{span}\{v_1, \dots, v_k\}$; it must be seen that $\{v_1, \dots, v_k, v_{k+1}\}$ is linearly independent.

By contradiction, suppose that it is linearly dependent: there are scalars $\alpha_1, \dots, \alpha_{k+1}$, not all zero, with
\[
\alpha_1 v_1 + \dots + \alpha_k v_k + \alpha_{k+1} v_{k+1} = 0.
\]
\begin{itemize}
    \item If $\alpha_{k+1} = 0$, the equation reduces to $\alpha_1 v_1 + \dots + \alpha_k v_k = 0$; since $\{v_1, \dots, v_k\}$ is linearly independent, all the $\alpha_i$ are $0$, against not all of them being zero.
    \item If $\alpha_{k+1} \neq 0$, solving,
    \[
    v_{k+1} = \sum_{i=1}^k \left( -\frac{\alpha_i}{\alpha_{k+1}} \right) v_i \in \operatorname{span}\{v_1, \dots, v_k\},
    \]
    against the choice of $v_{k+1}$.
\end{itemize}
In both cases there is a contradiction, so that the set is linearly independent.
\end{proof}

\begin{teorema}{Existence of a Basis (Finite Case)}{existencia_base}
Every finite-dimensional $\mathbb{K}$-vector space $V \neq \{0\}$ has a basis.
\end{teorema}

\begin{proof}
By the definition of finite dimension (Definition \ref{def:dim_finita}), $V$ admits a finite generating set $G = \{v_1, \dots, v_k\}$. $G$ is reduced to a basis:
\begin{itemize}
    \item if $G$ is linearly independent, since it also generates $V$, it is already a basis;
    \item if $G$ is linearly dependent, by Proposition \ref{prop:reduccion_base} a vector can be removed, obtaining $G'$ with $\operatorname{span}(G') = V$ and one element fewer.
\end{itemize}
Since $G$ is finite, the elimination process stops in at most $k-1$ steps. The final set $\mathcal{B}$ generates $V$ and is linearly independent, hence it is a basis of $V$.
\end{proof}

\obs The zero subspace $S = \{0\}$ is the only subspace that does not have a basis.

\begin{teorema}{Relation between Generators and Linear Independence}{teorema_intercambio}
Let $V$ be a $\mathbb{K}$-vector space generated by a set of vectors $\{u_1, \dots, u_r\}$ (that is, $V = \operatorname{span}\{u_1, \dots, u_r\}$).
If $\{v_1, \dots, v_k\}$ is a set of linearly independent vectors contained in $V$, then:
\[
k \le r.
\]
\end{teorema}

\begin{proof}
By contradiction, suppose $k > r$. Since $\{u_1, \dots, u_r\}$ generates $V$ and each $v_j \in V$, there are scalars $a_{ij}$ with
\[
v_j = \sum_{i=1}^{r} a_{ij} u_i.
\]
A zero linear combination $\sum_{j=1}^{k} x_j v_j = 0$, after substituting, is
\[
\sum_{i=1}^{r} \left( \sum_{j=1}^{k} a_{ij} x_j \right) u_i = 0,
\]
and it suffices to make each coefficient vanish, that is, to solve the homogeneous system
\[
\begin{cases}
a_{11}x_1 + a_{12}x_2 + \dots + a_{1k}x_k = 0 \\
a_{21}x_1 + a_{22}x_2 + \dots + a_{2k}x_k = 0 \\
\vdots \\
a_{r1}x_1 + a_{r2}x_2 + \dots + a_{rk}x_k = 0
\end{cases}
\]
of $r$ equations and $k$ unknowns. Since $k > r$, it has a nontrivial solution, with some $x_j \neq 0$, so that $\sum_{j} x_j v_j = 0$ with coefficients not all zero and $\{v_1, \dots, v_k\}$ would be linearly dependent, against the hypothesis. Hence $k \le r$.
\end{proof}

\begin{corolario}{Invariance of the Number of Elements of a Basis}{invariancia_base}
All the bases of the same finite-dimensional $\mathbb{K}$-vector space $V$ have the same number of elements.
\end{corolario}

\begin{proof}
Let $B_1 = \{v_1, \dots, v_n\}$ and $B_2 = \{w_1, \dots, w_m\}$ be two bases of $V$. Considering $B_1$ as a generating set and $B_2$ as a linearly independent set, Theorem \ref{teo:teorema_intercambio} gives $m \le n$. Reversing the roles ($B_2$ generating, $B_1$ independent), the same theorem gives $n \le m$. Hence $n = m$, that is, both bases have the same number of elements.
\end{proof}

\begin{definicion}{Dimension of a Vector Space}{definicion_dimension}
Let $V$ be a finite-dimensional $\mathbb{K}$-vector space.
The \textbf{dimension} of $V$, denoted $\dim(V)$, is the number of vectors of any basis of $V$.

This number is well defined (it is unique for each space) thanks to Corollary \ref{cor:invariancia_base}, which guarantees that all the bases of the same vector space have the same number of elements.
\end{definicion}

\obs For $V = \{0\}$ the convention is that its basis is the empty set $\emptyset$; since it has $0$ elements, $\dim(\{0\}) = 0$.

\begin{proposicion}{Relation between Subspaces and Dimension}{relacion_subespacios}
Let $V$ be a $\mathbb{K}$-vector space. Let $S$ and $T$ be finite-dimensional subspaces of $V$ such that $T \subseteq S$. Then:
\begin{enumerate}
    \item $\dim(T) \le \dim(S)$.
    \item If $\dim(T) = \dim(S)$, then $T = S$.
\end{enumerate}
\end{proposicion}

\begin{proof}
\textit{(1)} Let $\mathcal{B}_T = \{v_1, \dots, v_k\}$ be a basis of $T$, with $k = \dim(T)$. Since $T \subseteq S$, its vectors are in $S$ and are linearly independent. Let $\mathcal{B}_S = \{u_1, \dots, u_r\}$ be a basis of $S$, with $r = \dim(S)$. By Theorem \ref{teo:teorema_intercambio} (the $v_i$ independent, the $u_i$ generators), $k \le r$, that is, $\dim(T) \le \dim(S)$.

\textit{(2)} Let $\dim(T) = \dim(S) = n$ and, by contradiction, suppose $S \neq T$. Since $T \subseteq S$, there is $x \in S$ with $x \notin T$. If $\mathcal{B}_T = \{v_1, \dots, v_n\}$ is a basis of $T$, since $x \notin T = \operatorname{span}(\mathcal{B}_T)$, the set $\{v_1, \dots, v_n, x\}$ is linearly independent in $S$, with $n+1$ vectors. This is absurd, since $\dim(S) = n$ and, by Corollary \ref{cor:invariancia_base}, there are no more than $n$ independent vectors in $S$. Hence $S = T$.
\end{proof}

\begin{teorema}{Basis Extension Theorem}{extension_base}
Let $V$ be a finite-dimensional $\mathbb{K}$-vector space and let $A$ be a set of linearly independent vectors in $V$.
If $A$ is not a basis (that is, it does not generate $V$), then it is possible to add vectors to $A$ until it becomes a basis of $V$.
\end{teorema}

\begin{proof}
Let $A_0 = A$ and consider the iterative process: if $\operatorname{span}(A_k) = V$, then $A_k$ is a basis and the process ends; if $\operatorname{span}(A_k) \subsetneq V$, there is $v_{k+1} \in V \setminus \operatorname{span}(A_k)$, and adding to a linearly independent set a vector outside its span preserves independence, so that $A_{k+1} = A_k \cup \{v_{k+1}\}$ is still linearly independent.

Since $V$ is finite-dimensional, it admits a generating set with $n$ elements, and by Corollary \ref{cor:invariancia_base} there are no linearly independent sets with more than $n$ vectors; the process stops at some step $m$ (with $|A_m| \le n$). When it stops, $\operatorname{span}(A_m) = V$, and since $A_m$ is linearly independent, it is the desired basis extending $A$.
\end{proof}

\begin{teorema}{Dimension of the Sum Theorem}{dim_suma}
Let $S$ and $T$ be finite-dimensional subspaces of a $\mathbb{K}$-vector space $V$. Then:
\[
\dim(S+T) = \dim(S) + \dim(T) - \dim(S \cap T).
\]
\end{teorema}

\begin{proof}
Let $\{v_1, \dots, v_k\}$ be a basis of $S \cap T$, so that $\dim(S \cap T) = k$. Since it is linearly independent and contained in $S$, it is extended to a basis of $S$ by adding $s_1, \dots, s_p$ (Theorem \ref{teo:extension_base}):
\[
B_S = \{v_1, \dots, v_k, s_1, \dots, s_p\}.
\]
Then $\dim(S) = k + p$. Analogously, it is extended to a basis of $T$ by adding $t_1, \dots, t_q$:
\[
B_T = \{v_1, \dots, v_k, t_1, \dots, t_q\}.
\]
Then $\dim(T) = k + q$. It is proved that
\[
B = \{v_1, \dots, v_k, s_1, \dots, s_p, t_1, \dots, t_q\}
\]
is a basis of $S+T$.

\textit{It generates $S+T$:} let $x \in S+T$, that is, $x = s + t$ with $s \in S$, $t \in T$. Since $B_S$ is a basis of $S$ and $B_T$ of $T$,
\[
s = \sum_{i=1}^{k} \alpha_i v_i + \sum_{j=1}^{p} \beta_j s_j, \qquad
t = \sum_{i=1}^{k} \delta_i v_i + \sum_{l=1}^{q} \gamma_l t_l,
\]
and adding,
\[
x = \sum_{i=1}^{k} (\alpha_i + \delta_i) v_i + \sum_{j=1}^{p} \beta_j s_j + \sum_{l=1}^{q} \gamma_l t_l,
\]
a linear combination of the elements of $B$. Hence $B$ generates $S+T$.

\textit{Linear independence:} consider a zero linear combination
\[
\sum_{i=1}^k \alpha_i v_i + \sum_{j=1}^p \beta_j s_j + \sum_{l=1}^q \gamma_l t_l = 0.
\]
Isolating the part with the vectors exclusive to $T$,
\[
\sum_{l=1}^q \gamma_l t_l = - \sum_{i=1}^k \alpha_i v_i - \sum_{j=1}^p \beta_j s_j =: y.
\]
The left-hand side is in $T$ and the right-hand side in $S$, so that $y \in S \cap T$, and it is written as a combination of its basis:
\[
y = \sum_{i=1}^k \delta_i v_i.
\]
Equating with the left-hand side,
\[
\sum_{l=1}^q \gamma_l t_l - \sum_{i=1}^k \delta_i v_i = 0,
\]
a zero linear combination of the vectors of $B_T$, which are linearly independent; in particular $\gamma_l = 0$ for every $l$. With this, the original equation reduces to
\[
\sum_{i=1}^k \alpha_i v_i + \sum_{j=1}^p \beta_j s_j = 0,
\]
a zero combination of the vectors of $B_S$, which are linearly independent, so that $\alpha_i = 0$ and $\beta_j = 0$. Since all the scalars are zero, $B$ is linearly independent.

Consequently, $B$ is a basis of $S+T$ with $k + p + q$ elements, and
\[
\dim(S) + \dim(T) - \dim(S \cap T) = (k+p) + (k+q) - k = k + p + q = \dim(S+T). \qedhere
\]
\end{proof}

\subsection{Orthogonal and Orthonormal Bases}

\begin{definicion}{Orthogonal/Orthonormal Sets and Bases}{base_orto}
Let $(V, \langle \cdot, \cdot \rangle)$ be a $\mathbb{K}$-vector space with an inner product and let $\mathcal{B} \subset V$ be a subset of vectors.

\begin{enumerate}
    \item \textbf{Orthogonal set:} $\mathcal{B}$ is an orthogonal set if all its vectors are nonzero and pairwise orthogonal:
    \[
    \forall u, v \in \mathcal{B}, \quad u \neq v \implies \langle u, v \rangle = 0.
    \]

    \item \textbf{Orthonormal set:} $\mathcal{B}$ is an orthonormal set if it is orthogonal and, moreover, each vector has unit norm.

    \item \textbf{Orthogonal/orthonormal basis:} if $\mathcal{B}$ is a basis of $V$ and an orthogonal/orthonormal set, it is called an orthogonal/orthonormal basis.
\end{enumerate}
\end{definicion}

\begin{proposicion}{Independence of Orthogonal Sets}{indep_ortogonal}
If $\{v_1, \dots, v_k\}$ is an orthogonal set of nonzero vectors, then the set is linearly independent.
\end{proposicion}

\begin{proof}
Consider a zero linear combination $\sum_{i=1}^{k} \alpha_i v_i = 0$; it is proved that all the $\alpha_i$ are zero. Given an index $j$, taking the inner product against $v_j$ and using linearity,
\[
\sum_{i=1}^{k} \alpha_i \langle v_i, v_j \rangle = \langle 0, v_j \rangle = 0.
\]
By orthogonality, $\langle v_i, v_j \rangle = 0$ for $i \neq j$, so that only
\[
\alpha_j \langle v_j, v_j \rangle = \alpha_j \|v_j\|^2 = 0
\]
survives. Since $v_j \neq 0$, $\|v_j\|^2 \neq 0$, hence $\alpha_j = 0$. Since this holds for every $j$, the set is linearly independent.
\end{proof}

\begin{teorema}{Projection Formula}{proyeccion_general}
Let $(V, \langle \cdot, \cdot \rangle)$ be a $\mathbb{K}$-vector space with an inner product and let $S$ be a finite-dimensional subspace.
Let $\mathcal{B} = \{v_1, \dots, v_k\}$ be an orthogonal basis of $S$.

Then, for every vector $v \in V$, the orthogonal projection of $v$ onto $S$ is given by:
\[
P_S(v) = \sum_{i=1}^{k} \frac{\langle v, v_i \rangle}{\|v_i\|^2} v_i.
\]
\end{teorema}

\begin{proof}
Since $S$ is finite-dimensional, $V = S \oplus S^\perp$ (Theorem \ref{teo:descomposicion_ortogonal}), so that every $v \in V$ is written uniquely as $v = s + t$ with $s \in S$ and $t \in S^\perp$, where $s = P_S(v)$. Since $\mathcal{B} = \{v_1, \dots, v_k\}$ is a basis of $S$,
\[
s = \sum_{j=1}^{k} c_j v_j,
\]
and it remains to determine the $c_j$. Taking the inner product of $v$ with $v_i$ and using linearity,
\[
\langle v, v_i \rangle = \langle s, v_i \rangle + \langle t, v_i \rangle = \langle s, v_i \rangle,
\]
since $t \in S^\perp$ and $v_i \in S$ give $\langle t, v_i \rangle = 0$. Substituting $s$ and using the orthogonality of the basis ($\langle v_j, v_i \rangle = 0$ for $j \neq i$),
\[
\langle v, v_i \rangle = \sum_{j=1}^{k} c_j \langle v_j, v_i \rangle = c_i \|v_i\|^2,
\]
whence $c_i = \langle v, v_i \rangle / \|v_i\|^2$ (the $v_i$ are nonzero). Therefore,
\[
P_S(v) = \sum_{i=1}^{k} \frac{\langle v, v_i \rangle}{\|v_i\|^2} v_i. \qedhere
\]
\end{proof}

\begin{corolario}{Orthonormal Case}{caso_ortonormal}
If the basis $\mathcal{B} = \{u_1, \dots, u_k\}$ is orthonormal, the formula simplifies:
\[
P_S(v) = \sum_{i=1}^{k} \langle v, u_i \rangle u_i.
\]
\end{corolario}

\begin{teorema}{Gram-Schmidt Orthogonalization Process}{gram_schmidt}
Let $(V, \langle \cdot, \cdot \rangle)$ be a $\mathbb{K}$-vector space with an inner product and let $S$ be a finite-dimensional subspace.
Given an arbitrary basis $\mathcal{B} = \{u_1, \dots, u_n\}$ of $S$, there is an orthogonal basis $\mathcal{O} = \{v_1, \dots, v_n\}$ of $S$ that satisfies the property of preservation of the generated subspaces:
\[
\operatorname{span}\{u_1, \dots, u_k\} = \operatorname{span}\{v_1, \dots, v_k\}, \quad \forall k \in \{1, \dots, n\}.
\]
\end{teorema}

\begin{proof}
Define the vectors recursively:
\begin{align*}
v_1 &= u_1 \\
v_2 &= u_2 - \frac{\langle u_2, v_1 \rangle}{\|v_1\|^2} v_1 \\
&\vdots \\
v_k &= u_k - \sum_{j=1}^{k-1} \frac{\langle u_k, v_j \rangle}{\|v_j\|^2} v_j.
\end{align*}

\textit{The $v_k$ are nonzero.} If $v_k = 0$, the definition of $v_k$ would give
$u_k \in \operatorname{span}\{v_1, \dots, v_{k-1}\}$, which by the preservation of the spans
(proved below, with index $k-1$) is $\operatorname{span}\{u_1, \dots, u_{k-1}\}$; this
contradicts the linear independence of $\mathcal{B}$. In particular, the divisions by
$\|v_j\|^2$ are well defined.

\textit{$\{v_1, \dots, v_n\}$ is orthogonal.} By induction on $k$. The case $\{v_1\}$ holds vacuously. Suppose $\{v_1, \dots, v_{k-1}\}$ orthogonal; it must be seen that $\langle v_k, v_i \rangle = 0$ for $i < k$. By the definition of $v_k$ and linearity,
\[
\langle v_k, v_i \rangle = \langle u_k, v_i \rangle - \sum_{j=1}^{k-1} \frac{\langle u_k, v_j \rangle}{\|v_j\|^2} \langle v_j, v_i \rangle.
\]
By the inductive hypothesis, $\langle v_j, v_i \rangle = 0$ except for $j = i$, so that
\[
\langle v_k, v_i \rangle = \langle u_k, v_i \rangle - \frac{\langle u_k, v_i \rangle}{\|v_i\|^2} \|v_i\|^2 = 0.
\]
Hence $\{v_1, \dots, v_k\}$ is orthogonal.

\textit{$\operatorname{span}\{u_1, \dots, u_p\} = \operatorname{span}\{v_1, \dots, v_p\}$ for every $p$.} By induction on $p$. The case $p = 1$ holds because $v_1 = u_1$. Suppose the equality for $p-1$. From the construction,
\[
v_p = u_p - \sum_{j=1}^{p-1} c_j v_j \implies u_p = v_p + \sum_{j=1}^{p-1} c_j v_j.
\]
\textit{($\subseteq$).} $u_p \in \operatorname{span}\{v_1, \dots, v_p\}$ by the previous equality, and $u_1, \dots, u_{p-1} \in \operatorname{span}\{v_1, \dots, v_{p-1}\} \subseteq \operatorname{span}\{v_1, \dots, v_p\}$ by the inductive hypothesis. \textit{($\supseteq$).} $v_p = u_p - w$ with $w \in \operatorname{span}\{v_1, \dots, v_{p-1}\} = \operatorname{span}\{u_1, \dots, u_{p-1}\}$, so that $v_p \in \operatorname{span}\{u_1, \dots, u_p\}$. By double inclusion, the subspaces coincide.

With $p = n$, $\mathcal{O}$ generates $S$, and being an orthogonal set of nonzero vectors it is linearly independent (Proposition \ref{prop:indep_ortogonal}); hence $\mathcal{O}$ is an orthogonal basis of $S$.
\end{proof}

\obs The construction of the vectors $v_k$ is called the \textbf{Gram-Schmidt orthogonalization process}.

\subsection{Coordinates and Matrix Representation}

\begin{definicion}{Coordinate Vector}{vector_coordenadas}
Let $V$ be a vector space of finite dimension $n$ and let $\mathcal{B} = \{v_1, v_2, \dots, v_n\}$ be an ordered basis of $V$.
For any vector $v \in V$, there are unique scalars $\alpha_1, \dots, \alpha_n$ such that:
\[ v = \alpha_1 v_1 + \alpha_2 v_2 + \dots + \alpha_n v_n \]
The \textbf{coordinate vector} of $v$ in the basis $\mathcal{B}$, denoted $[v]_\mathcal{B}$, is the column vector of these scalars:
\[ [v]_\mathcal{B} = \begin{pmatrix} \alpha_1 \\ \alpha_2 \\ \vdots \\ \alpha_n \end{pmatrix} \in \mathbb{K}^n \]
\end{definicion}

\obs One must distinguish the vector $v$ (an abstract object) from $[v]_\mathcal{B}$: when the basis $\mathcal{B}$ changes, the vector $v$ remains the same, but its coordinates $[v]_\mathcal{B}$ change.

\subsection{Existence of Bases in General Spaces and Curiosities}

Reading this section is optional. It is suggested to become familiar beforehand with Zorn's Lemma and the Baire Category Theorem (in one of its versions).

\begin{teorema}{Existence of a Basis (General Case)}{existencia_base_zorn}
Every vector space $V$ has a basis.
\end{teorema}

\begin{proof}
Let $\mathcal{X}$ be the set of linearly independent subsets of $V$, ordered by inclusion: $A \le B \iff A \subseteq B$. It is nonempty, since $\emptyset \in \mathcal{X}$.

Let $\mathcal{Y}$ be a chain in $\mathcal{X}$; it must be seen that it has an upper bound in $\mathcal{X}$. Let $C = \bigcup_{B \in \mathcal{Y}} B$; it suffices to see that $C$ is linearly independent. By contradiction, suppose that it is dependent: there are $v_1, \dots, v_n \in C$ and scalars $\lambda_i$ not all zero with $\sum_{i=1}^n \lambda_i v_i = 0_V$. Each $v_i$ is in some $B_i \in \mathcal{Y}$; since $\{B_1, \dots, B_n\}$ is finite and $\mathcal{Y}$ is a chain, it has a maximum $B_{\max}$, and all the $v_i \in B_{\max}$. But then $B_{\max}$ would contain a dependent set, against $B_{\max} \in \mathcal{X}$. Hence $C$ is linearly independent and, since $B \subseteq C$ for every $B \in \mathcal{Y}$, it is an upper bound of the chain.

By Zorn's lemma, $\mathcal{X}$ has a maximal element $G$, which is a basis of $V$:
\begin{itemize}
    \item \textit{Independence:} $G \in \mathcal{X}$ is linearly independent.
    \item \textit{Generation:} if $G$ did not generate $V$, there would be $v \notin \operatorname{span}(G)$, and $G' = G \cup \{v\}$ would be linearly independent (if it were dependent, $v$ would be a linear combination of $G$), with $G \subsetneq G'$ in $\mathcal{X}$, against the maximality of $G$.
\end{itemize}
Hence $G$ generates $V$ and is a basis.
\end{proof}

\begin{teorema}{Cardinality of the Basis}{base_frechet}
Let $V$ be an infinite-dimensional complete normed vector space (Banach space).
Then any Hamel basis of $V$ is uncountable.
\end{teorema}

\begin{proof}
By contradiction, suppose that $V$ admits a countable Hamel basis $\mathcal{B} = \{e_1, e_2, \dots\}$, and let
\[
F_n = \operatorname{span}\{e_1, \dots, e_n\}.
\]
These subspaces satisfy:
\begin{enumerate}
    \item \textit{They are closed:} every finite-dimensional subspace is closed.
    \item \textit{They have empty interior:} since $V$ is infinite-dimensional, $F_n \subsetneq V$ is proper; a proper subspace with nonempty interior would contain a neighborhood of the origin and would generate the whole space, against $F_n \neq V$. Hence $\operatorname{int}(F_n) = \emptyset$.
\end{enumerate}
Since $\mathcal{B}$ is a basis, every $x \in V$ is a finite linear combination of its elements, so that
\[
V = \bigcup_{n=1}^{\infty} F_n,
\]
a countable union of closed sets with empty interior (nowhere dense). This contradicts the Baire Category Theorem, according to which a complete metric space is not of the first category. Hence the Hamel basis cannot be countable.
\end{proof}

\obs This result must not be confused with bases of Hilbert spaces or Fourier series.
\begin{itemize}
    \item A Hamel basis (the algebraic one) requires finite linear combinations. In an infinite-dimensional Banach space, these bases are necessarily uncountable.
    \item A Schauder basis (the topological one) allows infinite linear combinations.
\end{itemize}
The previous theorem says that an infinite-dimensional Banach space cannot be generated using only finite sums of countably many vectors.

\obs On the relation between separability and Schauder bases:
\begin{enumerate}
    \item If a Banach space has a Schauder basis, then it is necessarily separable.
    \item The converse is false: not every separable Banach space admits a Schauder basis.
\end{enumerate}
This was a famous open problem for 40 years, solved negatively by Per Enflo in 1973, who constructed a separable Banach space that admits no Schauder basis.

\newpage

\section{Linear Transformations}

\subsection{Definition and First Properties}

\begin{definicion}{Linear Transformation}{def_trans_lineal}
Let $V$ and $W$ be two $\mathbb{K}$-vector spaces.
A function $T: V \to W$ is said to be a linear transformation if it satisfies the following two conditions:

\begin{enumerate}
    \item Additivity:
    For every pair of vectors $u, v \in V$:
    \[ T(u + v) = T(u) + T(v) \]
    \item Homogeneity:
    For every scalar $\lambda \in \mathbb{K}$ and every vector $v \in V$:
    \[ T(\lambda v) = \lambda T(v) \]
\end{enumerate}
\end{definicion}

\begin{proposicion}{Properties}{prop_estructurales_tl}
Let $T: V \to W$ be a linear transformation. The following properties hold:

\begin{enumerate}
    \item The origin is a fixed point: $T(0_V) = 0_W$.
    \item If the values that $T$ takes on a basis $\{v_1, \dots, v_n\}$ of $V$ are known, then $T(v)$ is determined for any vector $v \in V$.
    \item If $S$ is a vector subspace of $V$, then its image $T(S)$ is a vector subspace of $W$.
    \item Given two linear transformations $T: U \to V$ and $L: V \to W$, their composition $L \circ T: U \to W$ is also a linear transformation.
    \item If $W'$ is a subspace of $W$, then the preimage of $W'$, defined as:
    \[ T^{-1}(W') = \{ v \in V \mid T(v) \in W' \} \]
    is a vector subspace of $V$.
\end{enumerate}
\end{proposicion}

\begin{proof}
\textit{(1)} From $0_V + 0_V = 0_V$ and additivity, $T(0_V) = T(0_V) + T(0_V)$; subtracting $T(0_V)$, one gets $0_W = T(0_V)$.

\textit{(2)} Let $\mathcal{B} = \{v_1, \dots, v_n\}$ be a basis of $V$. Each $v \in V$ is written uniquely as $v = \alpha_1 v_1 + \dots + \alpha_n v_n$, and by linearity
\[
T(v) = \alpha_1 T(v_1) + \dots + \alpha_n T(v_n),
\]
so that $T(v)$ is determined by the $\alpha_i$ and the images $T(v_i)$.

\textit{(3)} Let $S$ be a subspace of $V$; it is verified that $T(S)$ is a subspace of $W$. Since $0_V \in S$ and $T(0_V) = 0_W$, $0_W \in T(S)$. If $y_1 = T(s_1)$ and $y_2 = T(s_2)$ with $s_1, s_2 \in S$, then $y_1 + y_2 = T(s_1 + s_2)$ with $s_1 + s_2 \in S$; and for $\lambda \in \mathbb{K}$, $\lambda y_1 = T(\lambda s_1)$ with $\lambda s_1 \in S$. Hence $T(S)$ is a subspace.

\textit{(4)} For $u, v \in U$ and $\lambda \in \mathbb{K}$, $(L \circ T)(u + v) = L(T(u) + T(v)) = (L \circ T)(u) + (L \circ T)(v)$ and $(L \circ T)(\lambda u) = \lambda (L \circ T)(u)$, so that $L \circ T$ is linear.

\textit{(5)} Let $Z = T^{-1}(W')$. Since $T(0_V) = 0_W \in W'$, $0_V \in Z$. If $u, v \in Z$, then $T(u), T(v) \in W'$ and $T(u+v) = T(u) + T(v) \in W'$, so that $u + v \in Z$. If $v \in Z$ and $\lambda \in \mathbb{K}$, $T(\lambda v) = \lambda T(v) \in W'$, so that $\lambda v \in Z$. Hence $T^{-1}(W')$ is a subspace of $V$.
\end{proof}

\subsection{The Matrix Representation Theorem}

As a visual complement to the relation between a linear transformation and the matrix product, \href{https://youtu.be/3gxjK9zkKqc?si=_vvIvbiUcNyruPvx}{this explanatory video} can be consulted.

\begin{definicion}{Canonical Basis of $\mathbb{K}^n$}{base_canonica_kn}
The \textbf{canonical basis} of $\mathbb{K}^n$ is the set $\{e_1, \dots, e_n\}$ where $e_j$ has a $1$ in position $j$ and $0$ in the others. Every $x \in \mathbb{K}^n$ is written uniquely as $x = \sum_{j=1}^n x_j e_j$.
\end{definicion}

\begin{teorema}{Matrix Representation}{teo_matriz_asociada}
Let $V$ and $W$ be spaces of dimensions $n$ and $m$, with bases $\mathcal{B}_V$ and $\mathcal{B}_W$. Let $T: V \to W$ be linear.
There is a unique matrix $A \in \mathbb{K}^{m \times n}$, denoted $[T]_{\mathcal{B}_V}^{\mathcal{B}_W}$, such that for every vector $v \in V$:
\[
[T(v)]_{\mathcal{B}_W} = [T]_{\mathcal{B}_V}^{\mathcal{B}_W} \cdot [v]_{\mathcal{B}_V}
\]
The columns of that matrix are built by computing the coordinates of the images of the basis of the domain:
\[ \text{Col}_j(A) = [T(v_j)]_{\mathcal{B}_W} \]
\end{teorema}

\begin{proof}
For existence, $A$ is built with the column formula of the statement. Let $v \in V$ with coordinates $x = [v]_{\mathcal{B}_V}$, that is, $v = x_1 v_1 + \dots + x_n v_n$. In $\mathbb{K}^n$, $x = x_1 e_1 + \dots + x_n e_n$, and by the distributivity of the matrix product,
\[ A x = x_1 (A e_1) + \dots + x_n (A e_n) = x_1 [T(v_1)]_{\mathcal{B}_W} + \dots + x_n [T(v_n)]_{\mathcal{B}_W}, \]
since $A e_j$ is column $j$ of $A$, equal to $[T(v_j)]_{\mathcal{B}_W}$ by construction. By the linearity of taking coordinates and that of $T$,
\[ A x = \left[ x_1 T(v_1) + \dots + x_n T(v_n) \right]_{\mathcal{B}_W} = \left[ T( x_1 v_1 + \dots + x_n v_n ) \right]_{\mathcal{B}_W} = [T(v)]_{\mathcal{B}_W}. \]
Hence the constructed matrix satisfies the required property.

For uniqueness, let $A$ and $A'$ be two matrices satisfying the property. Taking $v = v_j$ (whose coordinates are $[v_j]_{\mathcal{B}_V} = e_j$), one obtains $A e_j = [T(v_j)]_{\mathcal{B}_W} = A' e_j$ for every $j$: the columns coincide and $A = A'$.
\end{proof}

\subsection{Composition and Matrix Product}

\obs \textit{Order of multiplication.} Given $T: U \to V$ and $L: V \to W$, the matrix of the composition $L \circ T$ is the product of the associated matrices:
\[ [L \circ T]_{\mathcal{B}_U}^{\mathcal{B}_W} = [L]_{\mathcal{B}_V}^{\mathcal{B}_W} \cdot [T]_{\mathcal{B}_U}^{\mathcal{B}_V} \]
Note that the order is reversed with respect to the usual left-to-right writing, but it coincides with the order of evaluation of functions: $(L \circ T)(x) = L(T(x))$.

\subsection{The Matrix Product as Composition}

The rule for multiplying matrices (row by column) is obtained by requiring that the product of matrices represent the composition of linear transformations.

Consider two consecutive linear transformations
\[ T: U \to V \quad \text{and} \quad S: V \to W, \]
with associated matrices $A$ (for $T$) and $B$ (for $S$). Given $x \in U$, applying $T$ yields $y = Ax$, and applying $S$ yields $z = By$. The matrix $C$ representing the composition $S \circ T$ is sought, that is, such that $z = Cx$. Substituting,
\[ z = B(Ax). \]

At the level of components, the $i$-th component of $z$ is $z_i = \sum_{k} B_{ik} y_k$, and the $k$-th component of $y$ is $y_k = \sum_{j} A_{kj} x_j$. Substituting,
\[ z_i = \sum_{k} B_{ik} \left( \sum_{j} A_{kj} x_j \right) = \sum_{j} \left( \sum_{k} B_{ik} A_{kj} \right) x_j. \]
For this to take the form $z_i = \sum_{j} C_{ij} x_j$, the $(i, j)$ entry of $C$ must be the term in parentheses.

\obs
The formula for the matrix product
\[ (BA)_{ij} = \sum_{k} B_{ik} A_{kj} \]
is the only definition that guarantees that multiplying matrices is equivalent to composing the corresponding linear transformations.

\subsubsection{Associativity as a Consequence}

The arithmetic proof that the matrix product is associative ($(AB)C = A(BC)$) is laborious, since it involves triple sums. The conceptual proof is direct when using that multiplying matrices is equivalent to composing linear transformations.

\begin{proposicion}{Associativity of the Matrix Product}{asoc_matrices}
The matrix product is associative. That is, for matrices $A, B, C$ of compatible sizes:
\[ (AB)C = A(BC) \]
\end{proposicion}

\begin{proof}
% Correction: with T, S, R represented by A, B, C, the product AB represents T∘S
% (the order is reversed), not S∘T as stated before.
Let $T, S, R$ be the linear transformations represented by $A, B, C$. Since the matrix product represents composition in the same order ($AB$ represents $T \circ S$), the matrix $(AB)C$ represents the composition $(T \circ S) \circ R$ and the matrix $A(BC)$ represents $T \circ (S \circ R)$. Since composition of functions is associative, both coincide, and therefore $(AB)C = A(BC)$.
\end{proof}

\subsection{Kernel and Image}

\begin{definicion}{Kernel and Image}{def_nucleo_imagen}
Let $T: V \to W$ be a linear transformation.
\begin{enumerate}
    \item Kernel: $\ker(T) = \{v \in V \mid T(v) = 0_W\}$.
    \item Image: $\operatorname{Im}(T) = \{w \in W \mid \exists v \in V, T(v) = w\}$.
\end{enumerate}
\end{definicion}

\begin{proposicion}{Subspace Structure}{prop_subespacios_T}
The kernel $\ker(T)$ is a subspace of $V$ and the image $\operatorname{Im}(T)$ is a subspace of $W$.
\end{proposicion}

\begin{proof}
The kernel is the preimage of the subspace $\{0_W\}$, a subspace by Proposition \ref{prop:prop_estructurales_tl} (5). The image is $T(V)$, a subspace by Proposition \ref{prop:prop_estructurales_tl} (3).
\end{proof}

\begin{proposicion}{Monomorphisms and Kernel}{prop_monomorfismo}
$T$ is injective if and only if $\ker(T) = \{0_V\}$.
\end{proposicion}

\begin{proof}
Both implications are proved.

\textit{($\Rightarrow$).} Let $T$ be injective and $v \in \ker(T)$, that is, $T(v) = 0_W = T(0_V)$; by injectivity, $v = 0_V$. Hence $\ker(T) = \{0_V\}$.

\textit{($\Leftarrow$).} Let $\ker(T) = \{0_V\}$ and $u, v$ with $T(u) = T(v)$. By linearity, $T(u - v) = 0_W$, so that $u - v \in \ker(T) = \{0_V\}$, that is, $u = v$. Hence $T$ is injective.
\end{proof}

\begin{proposicion}{Characterization of Surjectivity}{prop_epi_img}
Let $T: V \to W$ be a linear transformation.
$T$ is surjective if and only if its image coincides with the whole codomain $W$.
\end{proposicion}

\begin{proof}
Both implications are proved.

\textit{($\Rightarrow$).} If $T$ is surjective, for every $w \in W$ there is $v$ with $T(v) = w$, so that $W \subseteq \operatorname{Im}(T)$; since always $\operatorname{Im}(T) \subseteq W$, one has $\operatorname{Im}(T) = W$.

\textit{($\Leftarrow$).} If $\operatorname{Im}(T) = W$, given $w \in W = \operatorname{Im}(T)$ there is $v$ with $T(v) = w$, so that $T$ is surjective.
\end{proof}

\subsection{Matrix Interpretation: Rows and Columns}

\begin{proposicion}{Image and Columns}{prop_img_cols}
Let $A$ be the matrix associated with $T$. Then the image of $T$ is the space generated by the columns of $A$:
\[ \operatorname{Im}(T) = \operatorname{Col}(A) \]
\end{proposicion}

\begin{proof}
Let $x \in \mathbb{K}^n$ and $C_1, \dots, C_n$ the columns of $A$. Then
\[
Ax = \begin{bmatrix} | & & | \\ C_1 & \dots & C_n \\ | & & | \end{bmatrix} \begin{bmatrix} x_1 \\ \vdots \\ x_n \end{bmatrix} = x_1 C_1 + x_2 C_2 + \dots + x_n C_n,
\]
a linear combination of the columns. Since every vector of the image is of the form $Ax$, the image is exactly the set of linear combinations of the columns, that is, $\operatorname{Col}(A)$.
\end{proof}

\begin{proposicion}{Kernel and Rows}{prop_ker_rows}
% TODO: for K = C the orthogonal complement argument requires the Hermitian product
% (conjugation); the statement and proof are written for the real case.
The kernel of $A$ is the orthogonal complement of the row space:
\[ \ker(A) = (\operatorname{Row}(A))^\perp \]
\end{proposicion}

\begin{proof}
Let $R_1, \dots, R_m$ be the rows of $A$. A vector $x$ belongs to the kernel if and only if $Ax = \mathbf{0}$, which, breaking the product down row by row through the standard inner product, is equivalent to
\[
Ax = \begin{bmatrix} \text{---} & R_1 & \text{---} \\ & \vdots & \\ \text{---} & R_m & \text{---} \end{bmatrix} \begin{bmatrix} x \end{bmatrix} = \begin{bmatrix} R_1 \cdot x \\ \vdots \\ R_m \cdot x \end{bmatrix} = \begin{bmatrix} 0 \\ \vdots \\ 0 \end{bmatrix},
\]
that is, to $R_i \cdot x = 0$ for every $i$. This means that $x$ is orthogonal to all the rows, and a vector orthogonal to a generating set is orthogonal to the whole subspace it generates. Hence $\ker(A) = (\operatorname{Row}(A))^\perp$.
\end{proof}

\subsection{Dimension Theorem and Rank}

\begin{teorema}{Dimension Theorem (Rank-Nullity)}{teo_dimension}
Let $T: V \to W$ be linear with $\dim(V) < \infty$. Then:
\[ \dim(V) = \dim(\ker(T)) + \dim(\operatorname{Im}(T)) \]
\end{teorema}

\begin{proof}
The strategy is to build a basis of $\operatorname{Im}(T)$ from a basis of the kernel extended to a basis of $V$. Let $n = \dim(V)$ and $k = \dim(\ker(T))$, and take a basis of the kernel $\mathcal{B}_K = \{u_1, \dots, u_k\}$. Being a linearly independent set in $V$, by the Basis Extension Theorem \ref{teo:extension_base} it is completed to a basis
\[ \mathcal{B}_V = \{u_1, \dots, u_k, v_1, \dots, v_r\}, \]
with $n = k + r$. It is proved that $B' = \{T(v_1), \dots, T(v_r)\}$ is a basis of $\operatorname{Im}(T)$, whence $\dim(\operatorname{Im}(T)) = r$.

\textit{$B'$ generates $\operatorname{Im}(T)$.} Let $w \in \operatorname{Im}(T)$, that is, $w = T(x)$ for some $x \in V$. Writing $x = \sum_{i=1}^k \alpha_i u_i + \sum_{j=1}^r \beta_j v_j$ and applying $T$, since $T(u_i) = 0_W$,
\[ w = \sum_{j=1}^r \beta_j T(v_j), \]
a linear combination of $B'$.

\textit{$B'$ is linearly independent.} Suppose $\sum_{j=1}^r c_j T(v_j) = 0_W$. By linearity, $T(\sum_j c_j v_j) = 0_W$, so that $\sum_j c_j v_j \in \ker(T)$ and it is written as $\sum_{j=1}^r c_j v_j = \sum_{i=1}^k d_i u_i$. Then
\[ \sum_{j=1}^r c_j v_j - \sum_{i=1}^k d_i u_i = 0_V \]
is a zero linear combination of the vectors of the basis $\mathcal{B}_V$, which forces all the scalars to be $0$; in particular, $c_j = 0$ for every $j$.

Hence $\dim(\operatorname{Im}(T)) = r$ and $\dim(V) = k + r = \dim(\ker(T)) + \dim(\operatorname{Im}(T))$.
\end{proof}

\begin{corolario}{Rank of a Matrix}{cor_rango_matriz}
Let $A \in \mathbb{K}^{m \times n}$. The number of linearly independent rows equals the number of linearly independent columns.
\[ \dim(\operatorname{Col}(A)) = \dim(\operatorname{Row}(A)) \]
\end{corolario}

\begin{proof}
Applying the Dimension Theorem \ref{teo:teo_dimension} to the transformation $T(x) = Ax$ and using that $\dim(\mathbb{K}^n) = n$,
\[ n = \dim(\ker(A)) + \dim(\operatorname{Col}(A)). \quad (*) \]
By Proposition \ref{prop:prop_ker_rows}, $\ker(A) = (\operatorname{Row}(A))^\perp$ inside $\mathbb{K}^n$. Since the sum of the dimensions of a subspace and its orthogonal complement is the dimension of the whole space,
\[ \dim(\operatorname{Row}(A)) + \dim(\ker(A)) = n. \quad (**) \]
Equating $(*)$ and $(**)$ and cancelling $\dim(\ker(A))$, one obtains $\dim(\operatorname{Col}(A)) = \dim(\operatorname{Row}(A))$.
\end{proof}

\subsection{Types of Morphisms}

\begin{definicion}{Homomorphism and Endomorphism}{homo_endo}
Let $V$ and $W$ be two $\mathbb{K}$-vector spaces. Every linear transformation $f: V \to W$ is called a homomorphism.
If the source and target spaces are the same ($V = W$), that is $f: V \to V$, $f$ is said to be an endomorphism.

\end{definicion}

\obs The concept of eigenvalue and eigenvector makes sense only for endomorphisms, since a vector $v$ is compared with its image $f(v)$ in the same space.

\begin{definicion}{Monomorphism, Epimorphism and Isomorphism}{mono_epi_iso}
Let $f: V \to W$ be a linear transformation.
\begin{enumerate}
    \item Monomorphism: if $f$ is injective, that is,
    \[ v_1 \neq v_2 \implies f(v_1) \neq f(v_2), \]
    which is equivalent to $\ker(f) = \{0\}$ (Proposition \ref{prop:prop_monomorfismo}).
    \item Epimorphism: if $f$ is surjective, that is, $\operatorname{Im}(f) = W$.
    \item Isomorphism: if $f$ is bijective (injective and surjective at the same time). In this case, $f$ is invertible and its inverse $f^{-1}: W \to V$ is also linear.
\end{enumerate}
\end{definicion}

\begin{definicion}{Automorphism}{automorfismo}
If $f: V \to V$ is an endomorphism that is moreover bijective (that is, an isomorphism of a space onto itself), it is called an automorphism.
\end{definicion}

\obs The set of all automorphisms of $V$, together with the operation of composition of functions, forms a group known as the General Linear Group of $V$, denoted by $GL(V)$. In matrix terms, it corresponds to the invertible square matrices.

\subsection{Equivalences for Endomorphisms}

\begin{corolario}{Equivalences}{cor_endomorfismos}
Let $T: V \to V$ be a linear transformation ($\dim(V)=n$). The following are equivalent:
\begin{enumerate}
    \item[(i)] $T$ is bijective.
    \item[(ii)] $T$ is injective.
    \item[(iii)] $T$ is surjective.
\end{enumerate}
\end{corolario}

\begin{proof}
(ii) $\Leftrightarrow$ (iii) is proved; together with the definitions, this gives the equivalence with (i).

\textit{(ii $\Rightarrow$ iii).} If $T$ is injective, $\dim(\ker(T)) = 0$, and by the Dimension Theorem \ref{teo:teo_dimension}, $\dim(\operatorname{Im}(T)) = n$. Since $\operatorname{Im}(T)$ is a subspace of $V$ with its same dimension, $\operatorname{Im}(T) = V$, so that $T$ is surjective.

\textit{(iii $\Rightarrow$ ii).} If $T$ is surjective, $\operatorname{Im}(T) = V$ and $\dim(\operatorname{Im}(T)) = n$. By the Dimension Theorem \ref{teo:teo_dimension}, $\dim(\ker(T)) = 0$, that is, $\ker(T) = \{0\}$, so that $T$ is injective.

Since injectivity and surjectivity are equivalent, either of them implies bijectivity (i).
\end{proof}

\begin{proposicion}{Matrix Version}{prop_matrices_invertibles}
Let $A \in \mathbb{K}^{n \times n}$. The following are equivalent:
\begin{enumerate}
    \item $A$ is invertible.
    \item The columns of $A$ are linearly independent.
    \item The columns of $A$ generate $\mathbb{K}^n$.
\end{enumerate}
\end{proposicion}

\begin{proof}
Let $T: \mathbb{K}^n \to \mathbb{K}^n$ be the transformation $T(x) = Ax$. The matrix $A$ is invertible if and only if $T$ is bijective. Since $Ax = x_1 C_1 + \dots + x_n C_n$ (a linear combination of the columns), $\ker(T) = \{0\}$ is equivalent to the columns being linearly independent, and $\operatorname{Im}(T) = \operatorname{Col}(A) = \mathbb{K}^n$ (Proposition \ref{prop:prop_img_cols}) is equivalent to the columns generating $\mathbb{K}^n$. By Corollary \ref{cor:cor_endomorfismos}, injectivity, surjectivity and bijectivity of $T$ are equivalent, whence the three conditions of the statement are as well.
\end{proof}

\subsection{Example: Orthogonal Projection}

\begin{proposicion}{Orthogonal Projection}{prop_proyeccion}
Let $V$ be an inner product space and $S$ a subspace such that $V = S \oplus S^\perp$.
Define $P_S(v) = s$, where $v = s + s^\perp$.
Then $P_S: V \to V$ is linear and its associated matrix $M$ satisfies $M^2 = M$ (idempotence).
\end{proposicion}

\begin{proof}
\textit{Linearity.} Let $u, v \in V$ with decompositions $u = u_S + u_\perp$ and $v = v_S + v_\perp$ with respect to $V = S \oplus S^\perp$, so that $P_S(u) = u_S$ and $P_S(v) = v_S$. Then
\[ u + v = (u_S + v_S) + (u_\perp + v_\perp), \]
with $u_S + v_S \in S$ and $u_\perp + v_\perp \in S^\perp$; by definition of the projection, $P_S(u+v) = u_S + v_S = P_S(u) + P_S(v)$. Analogously, for $\lambda \in \mathbb{K}$, $\lambda u = \lambda u_S + \lambda u_\perp$ with $\lambda u_S \in S$, hence $P_S(\lambda u) = \lambda u_S = \lambda P_S(u)$.

\textit{Idempotence of the operator.} Let $s = P_S(v) \in S$. Its decomposition is $s = s + 0$, so that $P_S(s) = s$, and then $P_S(P_S(v)) = s = P_S(v)$, that is, $P_S^2 = P_S$.

\textit{Idempotence of the matrix.} Let $M$ be the associated matrix. The matrix of $P_S \circ P_S$ is $M^2$, and since $P_S \circ P_S = P_S$, one has $M^2 = M$.
\end{proof}

\newpage

\section{The Dual Space}

\subsection{Definition and Motivation}

\begin{definicion}{Dual Space}{def_espacio_dual}
If $V$ is a vector space over a field $\mathbb{K}$, the \textbf{dual space}, denoted $V^*$, is the set of all linear transformations from $V$ to the field of scalars.
\[
V^* = \{ f: V \to \mathbb{K} \mid f \text{ is linear} \}.
\]
The elements of $V^*$ are called \textbf{linear functionals}.
\end{definicion}

The dual space is of interest for two reasons:

\begin{itemize}
    \item Object/measurement distinction: in the matrix product $x^T y$, the two vectors do not play the same role:
    \begin{itemize}
        \item the column vector $y \in V$ is the geometric object;
        \item the row vector $x^T \in V^*$ acts as an operator or measuring rule.
    \end{itemize}
    The dual space formalizes this difference: the measured object does not coincide with the rule with which it is measured.

    \item Correspondence (Riesz): in finite-dimensional inner product spaces (such as $\mathbb{R}^n$) there is a one-to-one correspondence between $V$ and $V^*$; for each linear functional $f \in V^*$ there is a unique vector $x \in V$ such that $f(y) = \langle x, y \rangle$.
\end{itemize}

\subsection{Dimension and Dual Basis}

\begin{definicion}{Dual Basis}{def_base_dual}
Let $V$ be a vector space of finite dimension $n$, and let $\mathcal{B} = \{v_1, v_2, \dots, v_n\}$ be a basis of $V$. 
The \textbf{dual basis}, denoted $\mathcal{B}^* = \{f_1, f_2, \dots, f_n\}$, is the set of linear functionals in $V^*$ defined by the condition:
\[
f_i(v_j) = \delta_{ij} = \begin{cases} 
1 & \text{if } i = j \\
0 & \text{if } i \neq j
\end{cases}
\]
\end{definicion}

\begin{teorema}{Existence and Uniqueness of the Dual Basis}{teo_base_dual}
If $V$ is a vector space of finite dimension $n$ with basis $\mathcal{B}$, then the set $\mathcal{B}^*$ defined above is indeed a basis of $V^*$ and, moreover, it is the unique basis that satisfies the Kronecker delta condition (the $\delta_{ij}$).
\end{teorema}

\begin{proof}
The strategy is to prove that $\mathcal{B}^*$ is linearly independent, that it generates $V^*$ and, finally, its uniqueness.

\textit{Independence.} Suppose $\sum_{i=1}^n c_i f_i = 0$. Evaluating at $v_j$,
\[
\left( \sum_{i=1}^n c_i f_i \right)(v_j) = \sum_{i=1}^n c_i \delta_{ij} = c_j = 0,
\]
so that $c_j = 0$ for every $j$.

\textit{Generation.} Given $\phi \in V^*$, let $\psi = \sum_{i=1}^n \phi(v_i) f_i$. For each $v_j$,
\[
\psi(v_j) = \sum_{i=1}^n \phi(v_i) \delta_{ij} = \phi(v_j),
\]
and since $\psi$ and $\phi$ coincide on the basis, $\phi = \psi$. Hence $\mathcal{B}^*$ is a basis of $V^*$, and $\dim(V^*) = n$.

\textit{Uniqueness.} If $\mathcal{G} = \{g_1, \dots, g_n\}$ also satisfies $g_i(v_j) = \delta_{ij}$, then $g_i$ and $f_i$ coincide on the whole basis $\mathcal{B}$, and since a linear transformation is determined by its values on a basis, $g_i = f_i$ for every $i$.
\end{proof}

\subsection{Importance of the Dual Basis}

\begin{proposicion}{Extraction of Coordinates}{prop_coords_dual}
Let $v \in V$. If $v$ is expressed in terms of the basis $\mathcal{B}$ as $v = \sum_{i=1}^n x_i v_i$, then the coefficients $x_i$ are given exactly by applying the dual basis:
\[
x_i = f_i(v)
\]
\end{proposicion}

\begin{proof}
Applying $f_i$ to $v = \sum_{j=1}^n x_j v_j$ and using linearity and the condition $f_i(v_j) = \delta_{ij}$,
\[
f_i(v) = \sum_{j=1}^n x_j f_i(v_j) = \sum_{j=1}^n x_j \delta_{ij} = x_i. \qedhere
\]
\end{proof}

This property allows writing any vector $v$ and any functional $\phi$ explicitly:

\begin{itemize}
    \item For vectors: $v = \sum_{i=1}^n f_i(v) v_i$
    \item For functionals: $\phi = \sum_{i=1}^n \phi(v_i) f_i$
\end{itemize}

\obs
In physics, this underlies the difference between contravariant components (the coordinates of the vector $x_i$) and covariant components (the components of the functional).

\begin{ejemplo}{The Case of $\mathbb{R}^3$}{ej_R3_dual}
Consider the standard basis $e_1, e_2, e_3$ of $\mathbb{R}^3$ (column vectors). The corresponding dual basis $e^*_1, e^*_2, e^*_3$ consists of the transposed row vectors:
\[
e^*_1(v) = \begin{bmatrix} 1 & 0 & 0 \end{bmatrix} v, \quad 
e^*_2(v) = \begin{bmatrix} 0 & 1 & 0 \end{bmatrix} v, \quad 
e^*_3(v) = \begin{bmatrix} 0 & 0 & 1 \end{bmatrix} v
\]
Their function is to extract the corresponding component of the vector on which they act.
\end{ejemplo}

The relation between the functionals of the dual space and the vectors of the space (via the inner product) admits a dynamic visualization in \href{https://www.youtube.com/watch?v=LyGKycYT2v0}{this explanatory video}.

\subsection{Recovery of the Primal Space}

\begin{proposicion}{Existence of a Pre-dual Basis}{prop_base_pre_dual}
Let $V$ be a $\mathbb{K}$-vector space of dimension $n$, and let $V^*$ be its dual space. 
Let $\mathcal{B}^* = \{\phi_1, \dots, \phi_n\}$ be a basis of $V^*$. 
Then there is a unique basis $\mathcal{B} = \{v_1, \dots, v_n\}$ of $V$ whose dual is exactly $\mathcal{B}^*$ (that is, $\phi_i(v_j) = \delta_{ij}$).
\end{proposicion}

\begin{proof}
\textit{Existence.} Let $\mathcal{W} = \{w_1, \dots, w_n\}$ be an auxiliary basis of $V$ and $M \in \mathbb{K}^{n \times n}$ the matrix $M_{ij} = \phi_i(w_j)$. Row $i$ of $M$ is the coordinate vector of $\phi_i$ in the dual basis $\mathcal{W}^*$ (since $\phi_i = \sum_j \phi_i(w_j)\, w_j^*$); since $\{\phi_1, \dots, \phi_n\}$ is linearly independent in $V^*$ and taking coordinates is an isomorphism, the rows of $M$ are linearly independent and $M$ is invertible. Let $A = (a_{ij}) = M^{-1}$, that is, $M \cdot A = I_n$, and define
\[
v_j = \sum_{k=1}^n a_{kj} w_k, \qquad 1 \le j \le n.
\]
Then
\[
\phi_i(v_j) = \sum_{k=1}^n a_{kj} \phi_i(w_k) = \sum_{k=1}^n a_{kj} M_{ik} = (M \cdot A)_{ij} = \delta_{ij}.
\]
Moreover, $\{v_1, \dots, v_n\}$ is a basis of $V$: since $\dim(V) = n$, linear independence suffices, and if $\sum_j \alpha_j v_j = 0$, applying $\phi_i$ one obtains $\alpha_i = 0$ for every $i$.

\textit{Uniqueness.} If $\mathcal{B}' = \{u_1, \dots, u_n\}$ is another basis with dual $\mathcal{B}^*$, its coordinates in the basis $\mathcal{B} = \{v_1, \dots, v_n\}$ are
\[
(u_i)_{\mathcal{B}} = (\phi_1(u_i), \dots, \phi_n(u_i)) = e_i = (v_i)_{\mathcal{B}},
\]
so that $u_i = v_i$ for every $i$.
\end{proof}

\begin{ejemplo}{Lagrange Interpolation}{ej_lagrange}
Let $V = \mathbb{R}_2[X]$ (polynomials of degree $\le 2$). Define the functionals $\epsilon_0, \epsilon_1, \epsilon_2 \in V^*$ as the evaluations at specific points:
\[
\epsilon_0(P) = P(0), \quad \epsilon_1(P) = P(1), \quad \epsilon_2(P) = P(2).
\]
It can be verified that $\{\epsilon_0, \epsilon_1, \epsilon_2\}$ is a basis of $V^*$. By the previous proposition, there must be a basis of polynomials $\{P_0, P_1, P_2\}$ such that $\epsilon_i(P_j) = \delta_{ij}$.

This implies, for example, that $P_0$ must satisfy:
\[
P_0(0)=1, \quad P_0(1)=0, \quad P_0(2)=0.
\]
The polynomial that satisfies this is $P_0(X) = \frac{1}{2}(X-1)(X-2)$.

To find a polynomial $P$ such that $P(0)=\alpha$, $P(1)=\beta$, $P(2)=\gamma$, the coordinates in this basis are used:
\[
P = \alpha P_0 + \beta P_1 + \gamma P_2,
\]
which is the Lagrange interpolation formula.
\end{ejemplo}

\subsection{Annihilator of a Subspace}

\begin{definicion}{Annihilator}{def_anulador}
Let $V$ be a vector space and $S$ a subspace of $V$. The \textbf{annihilator} of $S$, denoted $S^\circ$, is the set of functionals that vanish at every vector of $S$:
\[
S^\circ = \{ f \in V^* \mid f(s) = 0, \forall s \in S \} = \{ f \in V^* \mid S \subseteq \ker(f) \}.
\]
\end{definicion}

\begin{teorema}{Dimension of the Annihilator}{teo_dim_anulador}
Let $V$ have finite dimension $n$ and let $S$ be a subspace. Then:
\[
\dim(S^\circ) = n - \dim(S).
\]
\end{teorema}

\begin{proof}
Let $\{v_1, \dots, v_r\}$ be a basis of $S$, extended to a basis $\mathcal{B} = \{v_1, \dots, v_n\}$ of $V$, and let $\mathcal{B}^* = \{\phi_1, \dots, \phi_n\}$ be the dual basis. For $i > r$, $\phi_i$ vanishes at $v_1, \dots, v_r$, so that $\phi_i \in S^\circ$, and $\{\phi_{r+1}, \dots, \phi_n\}$ is linearly independent. Conversely, every $g \in S^\circ$ is written as $g = \sum_i \alpha_i \phi_i$, and since $g(v_k) = 0$ for $k \le r$, the coefficients $\alpha_1, \dots, \alpha_r$ are zero; hence $g$ depends only on $\phi_{r+1}, \dots, \phi_n$, which therefore generate $S^\circ$. Consequently, $\dim(S^\circ) = n - r = n - \dim(S)$.
\end{proof}

\begin{proposicion}{Recovery of the Subspace}{prop_recuperacion_sub}
The original subspace can be recovered from its annihilator:
\[
S = \{ x \in V \mid f(x) = 0, \forall f \in S^\circ \}.
\]
\end{proposicion}
\begin{proof}
Let $T = \{ x \in V \mid f(x) = 0, \ \forall f \in S^\circ \}$ be the set in the statement. By the definition of the annihilator, $S \subseteq T$. For the reverse inclusion, argue by contradiction: if there were $x \in T$ with $x \notin S$, completing a basis of $S$ with $x$ and taking the dual basis would give a functional $\phi$ that vanishes on $S$ (hence $\phi \in S^\circ$) but with $\phi(x) = 1$, which contradicts $x \in T$. Therefore $T = S$.
\end{proof}

\subsection{Application: The Gradient and Duality}

\begin{definicion}{The Differential (The Covector)}{def_diferencial}
Let $f: V \to \mathbb{R}$ be a function differentiable at a point $x \in V$. The \textbf{differential} of $f$ at $x$, denoted $df_x$, is the linear functional in $V^*$ that best approximates the change of the function:
\[
f(x + v) \approx f(x) + df_x(v)
\]
That is, $df_x \in V^*$. It takes a displacement vector $v$ and returns a number (the infinitesimal change).
\end{definicion}

Now, $df_x$ is a functional, not a direction of motion. To convert it into a vector an additional structure is required: an inner product.

\begin{teorema}{The Gradient (Riesz Representative)}{teo_gradiente_riesz}
Given a vector space $V$ with inner product $\langle \cdot, \cdot \rangle$, by the Riesz Representation Theorem, for the functional $df_x \in V^*$ there is a unique vector $g \in V$ such that:
\[
df_x(v) = \langle g, v \rangle \quad \forall v \in V
\]
This vector $g$ is called the \textbf{gradient} and is denoted $\nabla f(x)$.
\end{teorema}

\begin{proof}
The real finite-dimensional case is proved. Let $\{e_1, \dots, e_n\}$ be an orthonormal basis of $V$ and define $g = \sum_{i=1}^n df_x(e_i)\, e_i$. For $v = \sum_i v_i e_i$, by the linearity of $df_x$ and the orthonormality of the basis,
\[
\langle g, v \rangle = \sum_{i=1}^n df_x(e_i)\, v_i = df_x\!\left( \sum_{i=1}^n v_i e_i \right) = df_x(v).
\]
For uniqueness, if $\langle g, v \rangle = \langle g', v \rangle$ for every $v$, then $\langle g - g', v \rangle = 0$; taking $v = g - g'$ yields $\|g - g'\|^2 = 0$, that is, $g = g'$.
\end{proof}

\obs
The differential $df_x$ exists independently of the geometry. The gradient $\nabla f(x)$ depends on the chosen inner product. If the way of measuring distances changes, the gradient changes, even though the function is the same.

\subsection{Gradient Descent}

\begin{proposicion}{Direction of Steepest Descent}{prop_max_descenso}
The direction $v$ that minimizes $df_x(v)$ subject to $\|v\|=1$ is opposite to the gradient:
\[
v = -\frac{\nabla f(x)}{\|\nabla f(x)\|}
\]
\end{proposicion}

\begin{proof}
By the definition of the gradient, $df_x(v) = \langle \nabla f(x), v \rangle$. By the Cauchy-Schwarz inequality, for $\|v\| = 1$,
\[
|\langle \nabla f(x), v \rangle| \le \|\nabla f(x)\| \|v\| = \|\nabla f(x)\|,
\]
and the minimum (most negative) value is attained when the cosine of the angle between $v$ and $\nabla f(x)$ is $-1$, that is, when $v = -\nabla f(x)/\|\nabla f(x)\|$.
\end{proof}

\newpage

\section{Determinants}

\subsection{Notation}

Let $\mathcal{M}_{n}(\mathbb{K})$ be the space of square matrices of size $n \times n$ with coefficients in a field $\mathbb{K}$.
Given a matrix $A \in \mathcal{M}_{n}(\mathbb{K})$, its columns are denoted as vectors $A^1, A^2, \dots, A^n$ in $\mathbb{K}^n$.

The matrix $A$ is written explicitly as the ordered list of these columns:
\[
A = (A^1 \mid A^2 \mid \dots \mid A^n)
\]
In this way, a function $D: \mathcal{M}_{n}(\mathbb{K}) \to \mathbb{K}$ can be written as:
\[
D(A) = D(A^1 \mid A^2 \mid \dots \mid A^n)
\]

\subsection{Alternating Multilinear Functions}

\begin{definicion}{Alternating Multilinear Function}{multilineal_matriz}
A function $D: \mathcal{M}_{n}(\mathbb{K}) \to \mathbb{K}$ is called \textbf{alternating multilinear} if it satisfies the following conditions:

1. Multilinearity: the function is linear with respect to each column separately. When all columns except the $i$-th are fixed, the function behaves as a linear transformation in that variable.
For every $i \in \{1, \dots, n\}$:
\begin{itemize}
    \item Additivity:
    \[ D(\dots \mid A^i + {A^i}' \mid \dots) = D(\dots \mid A^i \mid \dots) + D(\dots \mid {A^i}' \mid \dots) \]
    \item Homogeneity:
    \[ D(\dots \mid \lambda A^i \mid \dots) = \lambda D(\dots \mid A^i \mid \dots) \]
\end{itemize}

2. Alternation: the function vanishes if the matrix has two equal columns, that is,
\[
A^i = A^j \text{ for some } i \neq j \implies D(A) = 0.
\]
\end{definicion}

\begin{proposicion}{Elementary Properties}{determinantes}
Let $D: \mathcal{M}_{n}(\mathbb{K}) \to \mathbb{K}$ be an alternating multilinear function. Then:

\begin{enumerate}
    \item
    If a matrix has a column of zeros ($A^i = 0_{\mathbb{K}^n}$), the function equals 0.
    \[ D(\dots \mid 0 \mid \dots) = 0. \]

    \item
    If two columns are swapped, the value of the function changes sign.
    \[ D(\dots \mid A^i \mid \dots \mid A^j \mid \dots) = - D(\dots \mid A^j \mid \dots \mid A^i \mid \dots). \]

    \item
    The value of the function does not change if a multiple of another, different column is added to a column.
    \[ D(\dots \mid A^i + \alpha A^j \mid \dots \mid A^j \mid \dots) = D(\dots \mid A^i \mid \dots \mid A^j \mid \dots). \]

    \item
    If a column is a linear combination of the others, the function equals 0.
    \[ A^i = \sum_{j \neq i} \alpha_j A^j \implies D(A) = 0. \]
\end{enumerate}
\end{proposicion}

\begin{proof}
\textit{(1)} If $A^i = 0_{\mathbb{K}^n}$, writing $0_{\mathbb{K}^n} = 0 \cdot 0_{\mathbb{K}^n}$ and using homogeneity in column $i$,
\[
D(\dots \mid 0_{\mathbb{K}^n} \mid \dots) = 0 \cdot D(\dots \mid 0_{\mathbb{K}^n} \mid \dots) = 0.
\]

\textit{(2)} Let $i \neq j$. By alternation (Definition \ref{def:multilineal_matriz}), a matrix with the vector $A^i + A^j$ in positions $i$ and $j$ has two equal columns, hence
\[
0 = D(\dots \mid \underbrace{A^i+A^j}_{\text{pos } i} \mid \dots \mid \underbrace{A^i+A^j}_{\text{pos } j} \mid \dots).
\]
Expanding by additivity in positions $i$ and $j$ yields four terms; the two with a repeated column ($A^i$ with $A^i$, $A^j$ with $A^j$) vanish, and what remains is
\[
0 = D(\dots \mid A^i \mid \dots \mid A^j \mid \dots) + D(\dots \mid A^j \mid \dots \mid A^i \mid \dots),
\]
whence $D(\dots \mid A^i \mid \dots \mid A^j \mid \dots) = - D(\dots \mid A^j \mid \dots \mid A^i \mid \dots)$.

\textit{(3)} By additivity in column $i$ and then homogeneity,
\[
D(\dots \mid A^i + \alpha A^j \mid \dots \mid A^j \mid \dots) = D(\dots \mid A^i \mid \dots \mid A^j \mid \dots) + \alpha\, D(\dots \mid A^j \mid \dots \mid A^j \mid \dots).
\]
The second term has $A^j$ in positions $i$ and $j$, so that by alternation it equals $0$, and the value does not change.

\textit{(4)} If $A^i = \sum_{k \neq i} c_k A^k$, by multilinearity
\[
D(A) = \sum_{\substack{k=1 \\ k \neq i}}^{n} c_k \cdot D(\dots \mid \underbrace{A^k}_{\text{pos } i} \mid \dots \mid \underbrace{A^k}_{\text{pos } k} \mid \dots).
\]
Each term has $A^k$ repeated in positions $i$ and $k$, so that by alternation it is zero and $D(A) = 0$.
\end{proof}

\subsection{Existence and Uniqueness of the Determinant}

The following lemma, which relates an alternating multilinear function defined on $\mathbb{K}^{(n+1) \times (n+1)}$ to others defined on $\mathbb{K}^{n \times n}$, is the basis of the recursive argument that allows proving the uniqueness of the determinant function.

\begin{lema}\label{lem:insercion}
% Correction: the statement said f_i(I_n) = α, but for i ≥ 2 the matrix built
% on I_n is a cyclic permutation of I_{n+1} and the correct value is (−1)^{i+1}·α
% (the previous version contradicted the final uniqueness formula).
Let $f: \mathbb{K}^{(n+1) \times (n+1)} \to \mathbb{K}$ be an alternating multilinear function such that $f(I_{n+1}) = \alpha$, and let $i$ with $1 \le i \le n+1$. Given $A \in \mathbb{K}^{n \times n}$, build $\tilde{A} \in \mathbb{K}^{(n+1) \times (n+1)}$ as follows:
\begin{itemize}
    \item column $1$ of $\tilde{A}$ is the canonical vector $e_i \in \mathbb{K}^{n+1}$;
    \item row $i$ of $\tilde{A}$ has zeros in all positions except $(i,1)$;
    \item the remaining coefficients correspond to the matrix $A$ keeping their relative order (the rows of $A$ occupy, in order, the rows of $\tilde{A}$ other than $i$, and its columns, the columns $2, \dots, n+1$).
\end{itemize}
Define $f_i: \mathbb{K}^{n \times n} \to \mathbb{K}$ by $f_i(A) = f(\tilde{A})$.
Then $f_i$ is an alternating multilinear function and $f_i(I_n) = (-1)^{i+1} \alpha$.
\end{lema}

\begin{proof}
The columns of $A$ correspond to columns $2, \dots, n+1$ of $\tilde{A}$: adding two columns of $A$ or multiplying one by a scalar produces exactly the same operation on the corresponding columns of $\tilde{A}$ (in row $i$ those columns equal $0$, so that they are not affected), and two equal columns of $A$ give two equal columns of $\tilde{A}$. Thus, the multilinearity and alternation of $f$ carry over to $f_i$.

For the value at the identity: if $A = I_n$, the columns of $\tilde{I_n}$ are, in order, $e_i$, $e_1$, $e_2$, \dots, $e_{i-1}$, $e_{i+1}$, \dots, $e_{n+1}$. Moving the first column to position $i$ through $i-1$ consecutive swaps yields $I_{n+1}$, and each swap changes the sign (Proposition \ref{prop:determinantes}), so that
\[
f_i(I_n) = f(\tilde{I_n}) = (-1)^{i-1} f(I_{n+1}) = (-1)^{i+1} \alpha. \qedhere
\]
\end{proof}

\begin{teorema}{Existence and Uniqueness}{existencia_unicidad}
Let $\alpha \in \mathbb{K}$. For each $n \in \mathbb{N}$, there is a unique alternating multilinear function $f: \mathbb{K}^{n \times n} \to \mathbb{K}$ such that $f(I_n) = \alpha$.
\end{teorema}

\begin{proof}
Given $A \in \mathbb{K}^{(n+1) \times (n+1)}$, denote by $A(i|j) \in \mathbb{K}^{n \times n}$ the matrix that results from deleting row $i$ and column $j$ of $A$. Existence and uniqueness are proved by induction on $n$.

\textit{Existence.} For $n=1$, $f(x) = \alpha \cdot x$ is alternating multilinear and satisfies $f(1) = \alpha$. Suppose that there is an alternating multilinear $g: \mathbb{K}^{n \times n} \to \mathbb{K}$ with $g(I_n) = \alpha$, and define $f: \mathbb{K}^{(n+1) \times (n+1)} \to \mathbb{K}$ by
\[
f(A) = \sum_{i=1}^{n+1} (-1)^{i+1} a_{1i} \cdot g(A(1|i)), \qquad A=(a_{jl})_{1 \le j,l \le n+1}.
\]
It is verified that $f$ is alternating multilinear and that $f(I_{n+1}) = \alpha$.

\textit{(i) Multilinearity.} Let $A = (A_1|\dots|A_k|\dots|A_{n+1})$, $A' = (A_1|\dots|A'_k|\dots|A_{n+1})$ and $\tilde{A} = (A_1|\dots|A_k+A'_k|\dots|A_{n+1})$. Then
\[
f(\tilde{A}) = \sum_{i \neq k}^{n+1} (-1)^{i+1} a_{1i} g(\tilde{A}(1|i)) + (-1)^{k+1} (a_{1k} + a'_{1k}) g(\tilde{A}(1|k)).
\]
If $i \neq k$, column $k$ of $\tilde{A}(1|i)$ is the sum of those of $A(1|i)$ and $A'(1|i)$, and by the multilinearity of $g$, $g(\tilde{A}(1|i)) = g(A(1|i)) + g(A'(1|i))$. If $i = k$, column $k$ is deleted and $g(\tilde{A}(1|k)) = g(A(1|k))$. Regrouping,
\[
f(\tilde{A}) = \left( \sum_{i=1}^{n+1} (-1)^{i+1} a_{1i} g(A(1|i)) \right) + \left( \sum_{i=1}^{n+1} (-1)^{i+1} a'_{1i} g(A'(1|i)) \right) = f(A) + f(A').
\]
The case of the product by a scalar $\lambda A_k$ is analogous.

\textit{(ii) Alternation.} Let $A = (A_1|\dots|A_j|\dots|A_k|\dots|A_{n+1})$ with column $k$ equal to column $j$ ($k > j$). Then
\[
f(A) = \sum_{i \neq k,j}^{n+1} (-1)^{i+1} a_{1i} g(A(1|i)) + (-1)^{j+1} a_{1j} g(A(1|j)) + (-1)^{k+1} a_{1k} g(A(1|k)).
\]
For each $i \neq j,k$, the matrix $A(1|i)$ has two equal columns, so that $g(A(1|i)) = 0$. Moreover, $A(1|j)$ is obtained from $A(1|k)$ through $k-1-j$ column swaps, so that $g(A(1|k)) = (-1)^{k-1-j} g(A(1|j))$. Hence
\begin{align*}
f(A) &= (-1)^{j+1} a_{1j} g(A(1|j)) + (-1)^{k+1} a_{1j} (-1)^{k-1-j} g(A(1|j)) \\
&= a_{1j} g(A(1|j)) \big( (-1)^{j+1} + (-1)^{2k-j} \big),
\end{align*}
and since $(-1)^{2k-j} = (-1)^{j} = -(-1)^{j+1}$, the sum is $0$.

\textit{(iii) Value at the identity.}
\[
f(I_{n+1}) = \sum_{i=1}^{n+1} (-1)^{i+1} (I_{n+1})_{1i} \cdot g(I_{n+1}(1|i)) = (-1)^{2} \cdot 1 \cdot g(I_n) = \alpha.
\]

\textit{Uniqueness.} For $n=1$, a function $f: \mathbb{K}^{1 \times 1} \to \mathbb{K}$ is alternating multilinear if and only if it is linear, and the value $f(1)$ determines it: $f(x) = f(1) \, x$. Suppose uniqueness on $\mathbb{K}^{n \times n}$ for each prescribed value at the identity (for each $\beta \in \mathbb{K}$ there is a unique alternating multilinear function that equals $\beta$ at $I_n$), let $g$ be the one that satisfies $g(I_n) = \alpha$, and let $f: \mathbb{K}^{(n+1) \times (n+1)} \to \mathbb{K}$ be alternating multilinear with $f(I_{n+1}) = \alpha$. By linearity in the first column, for $A = (a_{ij})$,
\[
f(A) = \sum_{i=1}^{n+1} a_{i1} f(e_i \mid A^2 \mid \dots \mid A^{n+1}),
\]
with $e_i$ the $i$-th canonical vector. In the term $f(e_i \mid A^2 \mid \dots \mid A^{n+1})$, subtracting from each column $j$ ($2 \le j \le n+1$) the first one multiplied by $a_{ij}$ (an operation that does not alter the value by Proposition \ref{prop:determinantes}), the resulting matrix has the vector $e_i$ in column $1$, zeros in row $i$ except at position $(i,1)$, and the submatrix $A(i|1)$ in the remaining rows and columns: it is exactly the matrix $\tilde{A}$ of Lemma \ref{lem:insercion} built on $A(i|1)$. Therefore,
\[
f(e_i \mid A^2 \mid \dots \mid A^{n+1}) = f_i(A(i|1)).
\]
By Lemma \ref{lem:insercion}, $f_i$ is alternating multilinear with $f_i(I_n) = (-1)^{i+1} \alpha$; the function $(-1)^{i+1} g$ is also alternating multilinear and takes that same value at $I_n$, so that by the inductive hypothesis $f_i = (-1)^{i+1} g$. Substituting,
\[
f(A) = \sum_{i=1}^{n+1} (-1)^{i+1} a_{i1} g(A(i|1)),
\]
so that $f$ is uniquely determined.
\end{proof}

\begin{definicion}{Determinant of Order $n$}{determinante}
The unique alternating multilinear function $f: \mathbb{K}^{n \times n} \to \mathbb{K}$ such that $f(I_n) = 1$ is called the \textbf{determinant} of order $n$.
It is usually denoted $\det: \mathbb{K}^{n \times n} \to \mathbb{K}$, and one writes $\det(A)$ or $|A|$.
\end{definicion}

\begin{corolario}{Expansion of the Determinant}{calculo_det}
Let $A \in \mathbb{K}^{n \times n}$. If for $r \in \mathbb{N}$, $\det: \mathbb{K}^{r \times r} \to \mathbb{K}$ denotes the determinant function of order $r$, then the following recursive formulas hold:

\[ \det(A) = \sum_{i=1}^{n} (-1)^{i+1} a_{i1} \det(A(i|1)) = \sum_{j=1}^{n} (-1)^{1+j} a_{1j} \det(A(1|j)) \]
\end{corolario}

\obs As a visual complement to the meaning of the determinant (the scale factor of an area or volume), \href{https://www.youtube.com/watch?v=Ip3X9LOh2dk}{this explanatory video} can be consulted.

\obs Any other alternating multilinear function $g$ on $\mathbb{K}^{n \times n}$ is a scalar multiple of the determinant.
Specifically, if $g$ is alternating multilinear and $g(I_n) = \alpha$, then by the uniqueness of Theorem \ref{teo:existencia_unicidad}:
\[ g(A) = \alpha \cdot \det(A) \]

\subsection{Properties}

\begin{proposicion}{Determinant of the Transpose}{transpuesta}
Let $A \in \mathbb{K}^{n \times n}$. Then:
\[ \det(A) = \det(A^t) \]
\end{proposicion}

\begin{proof}
By induction on $n$.

\textit{Base case ($n=1$).} If $A \in \mathbb{K}^{1 \times 1}$, then $A = A^t$ and the equality holds.

\textit{Inductive step.} Suppose the equality for $n$ and let $A \in \mathbb{K}^{(n+1) \times (n+1)}$. Let $B = A^t$, with $b_{ij} = a_{ji}$. Expanding $\det(B)$ along its first row,
\[
\det(B) = \sum_{i=1}^{n+1} (-1)^{1+i} b_{1i} \det(B(1|i)).
\]
Since $b_{1i} = a_{i1}$ and $B(1|i) = (A(i|1))^t$, by the inductive hypothesis $\det((A(i|1))^t) = \det(A(i|1))$, and then
\[
\det(A^t) = \sum_{i=1}^{n+1} (-1)^{1+i} a_{i1} \det(A(i|1)),
\]
which is the expansion of $\det(A)$ along the first column (Corollary \ref{cor:calculo_det}). Hence $\det(A^t) = \det(A)$.
\end{proof}

\obs As a consequence, the determinant is an alternating multilinear function of the rows. Everything proved for columns holds for rows.

\begin{proposicion}{Determinant of an Upper Triangular Matrix}{triangular}
Let $A = (a_{ij}) \in \mathbb{K}^{n \times n}$ be an upper triangular matrix ($a_{ij} = 0$ if $i > j$). Then:
\[ \det(A) = \prod_{i=1}^n a_{ii} \]
\end{proposicion}

\begin{proof}
By induction on $n$. For $n=1$, $\det(A) = a_{11}$. For the inductive step, let $A \in \mathbb{K}^{(n+1) \times (n+1)}$ be upper triangular. Expanding along the first column,
\[
\det(A) = \sum_{i=1}^{n+1} (-1)^{i+1} a_{i1} \det(A(i|1)),
\]
and since $a_{i1} = 0$ for $i \ge 2$, only the term $i=1$ remains:
\[
\det(A) = a_{11} \det(A(1|1)).
\]
The submatrix $A(1|1)$ is upper triangular of size $n \times n$ with diagonal $a_{22}, \dots, a_{n+1,n+1}$, so that by the inductive hypothesis $\det(A(1|1)) = \prod_{j=2}^{n+1} a_{jj}$, and
\[
\det(A) = a_{11} \prod_{j=2}^{n+1} a_{jj} = \prod_{j=1}^{n+1} a_{jj}.
\]
\end{proof}

\begin{proposicion}{General Laplace Expansion}{laplace_general}
Let $A \in \mathbb{K}^{n \times n}$.
\begin{enumerate}
    \item Along column $j$: $\det(A) = \sum_{i=1}^n (-1)^{i+j} a_{ij} \det(A(i|j))$
    \item Along row $i$: $\det(A) = \sum_{j=1}^n (-1)^{i+j} a_{ij} \det(A(i|j))$
\end{enumerate}
\end{proposicion}

\begin{proof}
The formula along columns is proved; the case along rows is analogous via the transpose. To expand along column $j$, that column is moved to position $1$ through $j-1$ successive swaps with the previous columns; if $A'$ is the reordered matrix,
\[ \det(A) = (-1)^{j-1} \det(A'). \]
Expanding $A'$ along its first column (the original column $j$ of $A$),
\[ \det(A') = \sum_{i=1}^n (-1)^{i+1} a'_{i1} \det(A'(i|1)), \]
and since $a'_{i1} = a_{ij}$ and $\det(A'(i|1)) = \det(A(i|j))$,
\[
\det(A) = (-1)^{j-1} \sum_{i=1}^n (-1)^{i+1} a_{ij} \det(A(i|j)) = \sum_{i=1}^n (-1)^{i+j} a_{ij} \det(A(i|j)).
\]
\end{proof}

\begin{proposicion}{Determinant of the Product}{producto}
Let $A, B \in \mathbb{K}^{n \times n}$. Then:
\[ \det(A \cdot B) = \det(A) \cdot \det(B) \]
\end{proposicion}

\begin{proof}
The strategy is to fix $A$ and verify that $f(X) = \det(A \cdot X)$ is alternating multilinear in the columns of $X$. Let $X_1, \dots, X_n$ be the columns of $X$; those of $AX$ are $AX_1, \dots, AX_n$.

\textit{Multilinearity.} For column $k$, by the distributivity of the matrix-vector product and the linearity of $\det$ in each column,
\begin{align*}
f(\dots \mid X_k + \lambda X'_k \mid \dots) &= \det(\dots \mid AX_k + \lambda AX'_k \mid \dots) \\
&= \det(\dots \mid AX_k \mid \dots) + \lambda \det(\dots \mid AX'_k \mid \dots).
\end{align*}

\textit{Alternation.} If $X_i = X_j$ with $i \neq j$, then $AX_i = AX_j$, so that $f(X) = \det(AX) = 0$.

By the uniqueness of Theorem \ref{teo:existencia_unicidad}, $f(X) = f(I_n) \cdot \det(X)$, and since $f(I_n) = \det(A \cdot I_n) = \det(A)$, taking $X = B$ one obtains $\det(A \cdot B) = \det(A) \cdot \det(B)$.
\end{proof}

\begin{proposicion}{Invertibility Criterion}{inversibilidad}
Let $A \in \mathbb{K}^{n \times n}$. The matrix $A$ is invertible if and only if its determinant is nonzero:
\[
A \text{ is invertible} \iff \det(A) \neq 0
\]
\end{proposicion}

\begin{proof}
Both implications are proved.

\textit{($\Rightarrow$).} If $A$ is invertible, there is $B$ with $A \cdot B = I_n$. Applying the determinant and the multiplicative property (Proposition \ref{prop:producto}), $\det(A) \cdot \det(B) = \det(I_n) = 1$, so that $\det(A) \neq 0$.

\textit{($\Leftarrow$).} If $\det(A) \neq 0$, its columns are linearly independent (Proposition \ref{prop:determinantes}), hence $A$ is invertible.
\end{proof}

\subsection{Adjugate of a Matrix}

\begin{definicion}{Adjugate}{adjunta}
Let $A=(a_{ij}) \in \mathbb{K}^{n \times n}$. The \textbf{adjugate} of $A$, denoted $adj(A)$, is the matrix $adj(A) \in \mathbb{K}^{n \times n}$ defined by:
\[
(adj(A))_{ij} = (-1)^{i+j} \det(A(j|i))
\]
\end{definicion}

\begin{proposicion}{Property of the Adjugate}{adjunta}
Let $A \in \mathbb{K}^{n \times n}$. Then:
\[ A \cdot adj(A) = \det(A) \cdot I_n \]
Hence, if $\det(A) \neq 0$:
\[ A^{-1} = \frac{1}{\det(A)} \cdot adj(A) \]
\end{proposicion}

\begin{proof}
Let $A = (a_{ij})$. The $(k,l)$ coefficient of the product is
\[
(A \cdot adj(A))_{kl} = \sum_{i=1}^{n} a_{ki} \cdot (adj(A))_{il} = \sum_{i=1}^{n} a_{ki} (-1)^{i+l} \det(A(l|i)).
\]

\textit{Case $k = l$.} The sum is the expansion of $\det(A)$ along row $l$, hence $(A \cdot adj(A))_{ll} = \det(A)$.

\textit{Case $k \neq l$.} The sum is the expansion along row $l$ of the determinant of the matrix $A^{(kl)}$ given by
\[
(A^{(kl)})_{ij} = \begin{cases} a_{ij} & \text{if } i \neq l, \\ a_{kj} & \text{if } i = l, \end{cases}
\]
which has rows $k$ and $l$ equal, so that its determinant is $0$.

Thus $A \cdot adj(A) = \det(A) \cdot I_n$, and if $\det(A) \neq 0$ one solves for $A^{-1} = \frac{1}{\det(A)} \cdot adj(A)$.
\end{proof}

\subsection{Cramer's Rule}

\begin{proposicion}{Cramer's Rule}{cramer}
Let $A=(a_{ij}) \in \mathbb{K}^{n \times n}$ be an invertible matrix, and let $b \in \mathbb{K}^{n \times 1}$.
Then the (unique) solution of the linear system $A \cdot x = b$ is given by:
\[
x_i = \frac{\det(A_i)}{\det(A)} \quad \text{for } i=1, \dots, n
\]
where $\det(A_i)$ represents the determinant of the matrix obtained by replacing the $i$-th column of $A$ by the vector $b$.
\end{proposicion}

\begin{proof}
Multiplying $A \cdot x = b$ by $adj(A)$ and using that $adj(A) \cdot A = A \cdot adj(A) = \det(A) \cdot I_n$ (Proposition \ref{prop:adjunta}),
\[
\det(A) \cdot x = adj(A) \cdot b.
\]
Taking coordinate $i$ and the definition of the adjugate (Definition \ref{def:adjunta}),
\[
\det(A) \cdot x_i = \sum_{j=1}^{n} (adj(A))_{ij} \cdot b_j = \sum_{j=1}^{n} (-1)^{i+j} b_j \det(A(j|i)),
\]
which is the expansion along column $i$ of the determinant of the matrix $A_i$ obtained by replacing column $i$ of $A$ by $b$. Therefore $x_i = \det(A_i)/\det(A)$.
\end{proof}

\subsection{Another Formula for the Determinant}

Given $A=(a_{ij}) \in \mathbb{K}^{n \times n}$, writing $A$ by its rows $A = (F_1, \dots, F_n)^t$ and using that the determinant is alternating multilinear in the rows, each row is decomposed in the canonical basis $e_1, \dots, e_n$. Expanding the multilinearity yields a sum over all possible column indices:
\[
\det(A) = \sum_{1 \le i_1, \dots, i_n \le n} a_{1 i_1} a_{2 i_2} \dots a_{n i_n} \det \begin{pmatrix} e_{i_1} \\ \vdots \\ e_{i_n} \end{pmatrix}
\]
Since the determinant vanishes if the matrix has two equal rows, in the previous sum only the terms where the indices $i_1, \dots, i_n$ are all distinct survive, that is, a reordering of $\{1, \dots, n\}$.

To formalize this, the following definition is introduced.

\begin{definicion}{Permutations}{permutaciones}
Let $I_n = \{1, 2, \dots, n\} \subseteq \mathbb{N}$. A \textbf{permutation} of $I_n$ is a bijective function $\sigma: I_n \to I_n$.
The set of all permutations of $I_n$ is denoted $S_n$.
\end{definicion}

With this notation, the sum is restricted to the elements of $S_n$:
\[
\det(A) = \sum_{\sigma \in S_n} a_{1 \sigma(1)} \dots a_{n \sigma(n)} \det \begin{pmatrix} e_{\sigma(1)} \\ \vdots \\ e_{\sigma(n)} \end{pmatrix}
\]
The determinant of the matrix formed by the permuted canonical vectors $\begin{pmatrix} e_{\sigma(1)} \\ \vdots \\ e_{\sigma(n)} \end{pmatrix}$ depends only on the permutation $\sigma$ and always equals $1$ or $-1$, since this matrix is obtained by swapping rows of the identity matrix $I_n$.
This allows defining the sign of the permutation.

\begin{definicion}{Sign of a Permutation}{signo_perm}
Given a permutation $\sigma \in S_n$, the \textbf{sign} of $\sigma$, denoted $sg(\sigma)$, is defined as the determinant of the corresponding permutation matrix:
\[
sg(\sigma) = \det \begin{pmatrix} e_{\sigma(1)} \\ \vdots \\ e_{\sigma(n)} \end{pmatrix}
\]
\end{definicion}

Substituting this into the previous expression yields the explicit formula of the determinant (Leibniz formula).

\begin{proposicion}{Permutation Formula}{leibniz}
Let $A = (a_{ij}) \in \mathbb{K}^{n \times n}$. Then:
\[
\det(A) = \sum_{\sigma \in S_n} sg(\sigma) a_{1 \sigma(1)} a_{2 \sigma(2)} \dots a_{n \sigma(n)}
\]
\end{proposicion}

\obs To compute the sign of a permutation $\sigma \in S_n$ it suffices to count how many row swaps must be performed on the identity matrix to obtain the permuted matrix.
If the number of swaps is $r$, the sign of the permutation is $(-1)^r$.

\begin{definicion}{Determinant of an Endomorphism}{det_endo}
Let $f: V \to V$ be an endomorphism of a finite-dimensional vector space. The \textbf{determinant} of $f$, denoted $\det(f)$, is defined as the determinant of the matrix associated with $f$ in any given basis. This value is independent of the chosen basis.
\end{definicion}

\obs Independence of the basis is justified as follows: two matrices $A$ and $B$ associated with $f$ in different bases are similar, $B = P^{-1} A P$ (see Theorem \ref{teo:cambio_base_semejanza}, below), and by the multiplicative property (Proposition \ref{prop:producto}),
\[
\det(B) = \det(P^{-1}) \det(A) \det(P) = \det(A),
\]
since $\det(P^{-1}) \det(P) = \det(I_n) = 1$.

\newpage

\section{Eigenvalues, Eigenvectors and Diagonalization}

\subsection{Definitions}

\begin{definicion}{Diagonal Matrix}{matriz_diagonal}
A square matrix $D \in \mathbb{K}^{n \times n}$ is called \textbf{diagonal} if all its entries off the main diagonal are zero.
Formally, if its entries $d_{ij}$ satisfy:
\[ d_{ij} = 0 \quad \text{for every } i \neq j \]
Visually, it has the following structure:
\[ D = \begin{pmatrix}
\lambda_1 & 0 & \dots & 0 \\
0 & \lambda_2 & \dots & 0 \\
\vdots & \vdots & \ddots & \vdots \\
0 & 0 & \dots & \lambda_n
\end{pmatrix} \]
\end{definicion}

\obs Diagonal matrices are simple to operate with:
\begin{itemize}
    \item their determinant is the product of the diagonal entries;
    \item raising them to a power $k$ is equivalent to raising each diagonal entry to the $k$-th power.
\end{itemize}

\begin{definicion}{Eigenvalue and Eigenvector}{autovalor_autovector}
Let $V$ be a $\mathbb{K}$-vector space and $f: V \to V$ an endomorphism.
A scalar $\alpha \in \mathbb{K}$ is an \textbf{eigenvalue} of $f$ if there is a vector $v \neq 0$ in $V$ such that:
\[ f(v) = \alpha v \]
In that case, $v$ is an \textbf{eigenvector} associated with $\alpha$.
\end{definicion}

\begin{definicion}{Eigenspace}{autoespacio}
Given an eigenvalue $\alpha$, the set of all eigenvectors associated with $\alpha$, together with the zero vector, forms a subspace of $V$, called the \textbf{eigenspace}, usually denoted by $S_\alpha$:
\[ S_\alpha = \{ v \in V \mid f(v) = \alpha v \} = \ker(f - \alpha \operatorname{id}_V), \]
where $\operatorname{id}_V$ is the identity transformation.
\end{definicion}

\obs The following properties of $S_\alpha$ follow from the definition:
\begin{itemize}
    \item $S_\alpha$ is a vector subspace of $V$ (it is the kernel of the linear transformation $f - \alpha \operatorname{id}_V$).
    \item The subspace is invariant under $f$: if $x \in S_\alpha$, then $f(x) = \alpha x \in S_\alpha$, since $S_\alpha$ is a subspace.
\end{itemize}

\subsection{Examples}

\begin{ejemplo}{Rotation in the Real Plane}{rotacion_plano}
Let $V = \mathbb{R}^2$ and $f$ the linear transformation that rotates a vector by 90 degrees:
\[ f(x, y) = (-y, x) \]
Geometrically, the transformed vector is orthogonal to the original one ($f(v) \perp v$). For an eigenvalue $\alpha$ to exist, $f(v)$ would have to be parallel to $v$. Since no nonzero vector can be orthogonal and parallel to itself at the same time, this transformation has no eigenvalues in $\mathbb{R}$.
\end{ejemplo}

\subsection{Computing Eigenvalues}

\begin{proposicion}{Equivalences}{equivalencias_autovalor}
Let $V$ be a vector space of finite dimension $n$. The following are equivalent:
\begin{enumerate}
    \item[i)] $\alpha$ is an eigenvalue of $f$.
    \item[ii)] The endomorphism $(f - \alpha \operatorname{id}_V)$ is not injective.
    \item[iii)] $\det(f - \alpha \operatorname{id}_V) = 0$.
\end{enumerate}
\end{proposicion}

\begin{proof}
\textit{(i) $\Leftrightarrow$ (ii).} One has $f(v) = \alpha v \iff (f - \alpha \operatorname{id}_V)(v) = 0$, and the existence of $v \neq 0$ with this property is equivalent to the kernel of $f - \alpha \operatorname{id}_V$ not being zero, that is, to it not being injective.

\textit{(ii) $\Leftrightarrow$ (iii).} In finite dimension, an endomorphism is injective if and only if it is an isomorphism (Corollary \ref{cor:cor_endomorfismos}), and it is an isomorphism if and only if its determinant is nonzero (Proposition \ref{prop:inversibilidad}). By contraposition, $f - \alpha \operatorname{id}_V$ is not injective $\iff \det(f - \alpha \operatorname{id}_V) = 0$.
\end{proof}

\obs Item (ii) gives a way of finding eigenvectors given an eigenvalue.

\begin{definicion}{Characteristic Polynomial}{polinomio_caracteristico}
The \textbf{characteristic polynomial} of $f$, denoted by $P(\alpha)$, is defined as:
\[ P(\alpha) = \det(f - \alpha \operatorname{id}_V) \]
\end{definicion}

\obs The problem of finding eigenvalues reduces to computing the roots of a polynomial.

\subsection{Similarity of Matrices}

\begin{definicion}{Similar Matrices}{matrices_semejantes}
Let $A, B \in \mathbb{K}^{n \times n}$ be two square matrices. $A$ is said to be \textbf{similar} to $B$ if there is an invertible matrix $P \in \mathbb{K}^{n \times n}$ such that:
\[ B = P^{-1} A P \]
\end{definicion}

\obs Similarity is an equivalence relation.

\begin{teorema}{Change of Basis and Similarity}{cambio_base_semejanza}
Let $V$ be a finite-dimensional vector space and $\mathcal{B}, \mathcal{B}'$ two bases of $V$.
If $A = [f]_\mathcal{B}$ and $B = [f]_{\mathcal{B}'}$ are the matrices associated with the same endomorphism $f: V \to V$ in those bases, then $A$ and $B$ are similar.
\end{teorema}

\begin{proof}
Let $P$ be the change-of-basis matrix from $\mathcal{B}'$ to $\mathcal{B}$ (its columns are the coordinates of the vectors of $\mathcal{B}'$ in the basis $\mathcal{B}$), which satisfies $[v]_\mathcal{B} = P [v]_{\mathcal{B}'}$ for every $v \in V$ and, being invertible, $[v]_{\mathcal{B}'} = P^{-1} [v]_\mathcal{B}$. By definition of the associated matrix, $[f(v)]_\mathcal{B} = A [v]_\mathcal{B}$ and $[f(v)]_{\mathcal{B}'} = B [v]_{\mathcal{B}'}$. Then
\[
[f(v)]_{\mathcal{B}'} = P^{-1} [f(v)]_\mathcal{B} = P^{-1} A [v]_\mathcal{B} = (P^{-1} A P) [v]_{\mathcal{B}'},
\]
and since this holds for every $v$, comparing with $[f(v)]_{\mathcal{B}'} = B [v]_{\mathcal{B}'}$ one obtains $B = P^{-1} A P$. Hence $A$ and $B$ are similar.
\end{proof}

\obs Since $A$ and $B$ represent the same endomorphism $f$, they share all intrinsic properties: they have the same determinant, the same trace and the same characteristic polynomial.

\subsection{Diagonalization}

\begin{definicion}{Diagonalizable Matrix}{matriz_diagonalizable}
A matrix $A \in \mathbb{K}^{n \times n}$ is called \textbf{diagonalizable} if it is similar to a diagonal matrix, that is, if there are an invertible matrix $P$ and a diagonal matrix $D$ such that:
\[ D = P^{-1} A P \]
\end{definicion}

\begin{definicion}{Diagonalizable Endomorphism}{endo_diagonalizable}
Let $V$ be a finite-dimensional vector space. An endomorphism $f: V \to V$ is called \textbf{diagonalizable} if there is a basis $\mathcal{B}$ of $V$ such that the associated matrix $[f]_\mathcal{B}$ is a diagonal matrix.
\end{definicion}

\obs If $f$ is diagonalizable, there is a coordinate system (the basis $\mathcal{B}$) in which the action of the linear transformation reduces to a scaling along each coordinate axis.

\begin{proposicion}{Equivalence of Definitions}{equiv_diag}
An endomorphism $f$ is diagonalizable if and only if its associated matrix $[f]_\mathcal{B}$ (in any basis $\mathcal{B}$) is a diagonalizable matrix.
\end{proposicion}

\begin{proof}
Both implications are proved.

\textit{($\Rightarrow$).} If $f$ is diagonalizable, there is a basis $\mathcal{B}'$ with $[f]_{\mathcal{B}'} = D$ diagonal. For any other basis $\mathcal{B}$, with $A = [f]_\mathcal{B}$, by Theorem \ref{teo:cambio_base_semejanza} $A$ and $D$ are similar, that is, $D = P^{-1} A P$ for some invertible $P$; hence $A$ is diagonalizable.

\textit{($\Leftarrow$).} If $A = [f]_\mathcal{B}$ is diagonalizable, there are $D$ diagonal and $P$ invertible with $D = P^{-1} A P$. The columns of $P$, which are linearly independent, are the coordinates of a basis $\mathcal{B}'$ in terms of $\mathcal{B}$, and by the change-of-basis formula $[f]_{\mathcal{B}'} = P^{-1} A P = D$; hence $f$ is diagonalizable.
\end{proof}

\begin{teorema}{Diagonalization Condition}{condicion_diag}
Let $f: V \to V$ be an endomorphism of a finite-dimensional vector space.
$f$ is diagonalizable if and only if there is a basis of $V$ formed by eigenvectors of $f$.
\end{teorema}

\begin{proof}
If $\mathcal{B} = \{v_1, \dots, v_n\}$, column $j$ of $[f]_\mathcal{B}$ contains the coordinates of $f(v_j)$ in that basis.

\textit{($\Rightarrow$).} If $f$ is diagonalizable, there is a basis $\mathcal{B} = \{v_1, \dots, v_n\}$ with $[f]_\mathcal{B} = D$ diagonal, so that column $j$ has $\alpha_j$ in position $j$ and zeros elsewhere; that is, $f(v_j) = \alpha_j v_j$, so that each $v_j$ is an eigenvector and $\mathcal{B}$ is a basis of eigenvectors.

\textit{($\Leftarrow$).} If $\mathcal{B} = \{v_1, \dots, v_n\}$ is a basis with $f(v_j) = \alpha_j v_j$, the coordinates of $f(v_j)$ in $\mathcal{B}$ are $(0, \dots, \alpha_j, \dots, 0)^T$, and the associated matrix is diagonal with $D_{jj} = \alpha_j$; hence $f$ is diagonalizable.
\end{proof}

\subsection{Operators on Inner Product Spaces}

In the following definitions $(V, \langle \cdot, \cdot \rangle)$ is assumed to be a finite-dimensional inner product space.

\begin{definicion}{Self-Adjoint Operator}{operador_autoadjunto}
Let $f: V \to V$ be an endomorphism. $f$ is said to be a \textbf{self-adjoint operator} if for every pair of vectors $u, v \in V$ the equality
\[ \langle f(u), v \rangle = \langle u, f(v) \rangle \]
holds.
\end{definicion}

\begin{definicion}{Orthogonal Operator}{operador_ortogonal}
Let $f: V \to V$ be an endomorphism. $f$ is said to be an \textbf{orthogonal operator} if it preserves the inner product, that is, for all $u, v \in V$:
\[ \langle f(u), f(v) \rangle = \langle u, v \rangle. \]
\end{definicion}

\obs \textit{Geometric Intuition:}

\begin{enumerate}
    \item[i)] Orthogonal Transformations (They preserve distances):
    The condition $\langle f(v), f(w) \rangle = \langle v, w \rangle$ implies that lengths and angles are preserved. Taking $w = v$:
    \[ \|f(v)\|^2 = \langle f(v), f(v) \rangle = \langle v, v \rangle = \|v\|^2 \]
    This means that the transformation neither stretches nor shrinks vectors; it works as a rigid motion (such as a rotation or a reflection) that keeps the geometry of the object intact.

    \item[ii)] Self-Adjoint Transformations (Symmetry):
    The condition \[ \langle f(u), v \rangle = \langle u, f(v) \rangle \] reflects a symmetry in the interaction between the vectors.
    Geometrically, these operators scale the space along perpendicular axes (their eigenvectors), stretching or compressing.
\end{enumerate}

\subsection{Matrix Representation of Self-Adjoint Operators}

\begin{proposicion}{Hermitian Matrix and Self-Adjoint Operator}{matriz_hermitica_iff}
Let $(V, \langle \cdot, \cdot \rangle)$ be a finite-dimensional inner product space and $\mathcal{B} = \{e_1, \dots, e_n\}$ an orthonormal basis.
Let $f: V \to V$ be an endomorphism and $A = [f]_\mathcal{B}$ its associated matrix.
Then:
\[ f \text{ is self-adjoint} \iff A \text{ is Hermitian, that is, } A = A^*, \]
where $A^*$ denotes the conjugate transpose of $A$ (in the real case, $A = A^T$: a symmetric matrix).
\end{proposicion}

\begin{proof}
Since $\mathcal{B}$ is orthonormal, $a_{ij} = \langle f(e_j), e_i \rangle$.

\textit{($\Rightarrow$).} If $f$ is self-adjoint, by the conjugate symmetry of the inner product,
\[ a_{ij} = \langle f(e_j), e_i \rangle = \langle e_j, f(e_i) \rangle = \overline{\langle f(e_i), e_j \rangle} = \overline{a_{ji}}, \]
hence $A = A^*$.

\textit{($\Leftarrow$).} If $A$ is Hermitian ($a_{ij} = \overline{a_{ji}}$), let $u = \sum_i x_i e_i$ and $v = \sum_j y_j e_j$. By linearity,
\[ \langle f(u), v \rangle = \sum_{i,j} x_i \overline{y_j} \langle f(e_i), e_j \rangle = \sum_{i,j} x_i \overline{y_j}\, a_{ji}, \]
\[ \langle u, f(v) \rangle = \sum_{i,j} x_i \overline{y_j} \langle e_i, f(e_j) \rangle = \sum_{i,j} x_i \overline{y_j}\, \overline{a_{ij}}. \]
Since $a_{ji} = \overline{a_{ij}}$, both sums coincide, so that $\langle f(u), v \rangle = \langle u, f(v) \rangle$.
\end{proof}

\begin{proposicion}{Unitary Matrix and Orthogonal Operator}{matriz_ortogonal_iff}
Let $(V, \langle \cdot, \cdot \rangle)$ be a finite-dimensional inner product space and $\mathcal{B} = \{e_1, \dots, e_n\}$ an orthonormal basis.
Let $f: V \to V$ be an endomorphism and $A = [f]_\mathcal{B}$ its associated matrix.
Then:
\[ f \text{ is orthogonal} \iff A \text{ is unitary (} A^{-1} = A^* \text{)} \]
\end{proposicion}

\begin{proof}
Let $C_1, \dots, C_n$ be the columns of $A$; by definition of the associated matrix, $C_j$ are the coordinates of $f(e_j)$ in $\mathcal{B}$. Since $\mathcal{B}$ is orthonormal, $\langle C_j, C_i \rangle_{\mathbb{K}^n} = \langle f(e_j), f(e_i) \rangle_V$, and since row $i$ of $A^*$ is the conjugate transpose of $C_i$,
\[ (A^* A)_{ij} = \langle C_j, C_i \rangle_{\mathbb{K}^n} = \langle f(e_j), f(e_i) \rangle_V. \]

\textit{($\Rightarrow$).} If $f$ is orthogonal, $\langle f(e_j), f(e_i) \rangle = \langle e_j, e_i \rangle = \delta_{ij}$, so that $(A^* A)_{ij} = \delta_{ij}$, that is, $A^* A = I$ and $A^{-1} = A^*$.

\textit{($\Leftarrow$).} If $A^* A = I$, then $\langle f(e_j), f(e_i) \rangle = (A^* A)_{ij} = \delta_{ij}$, so that $\{f(e_1), \dots, f(e_n)\}$ is an orthonormal basis. For $u = \sum_k x_k e_k$ and $v = \sum_k y_k e_k$,
\[ \langle f(u), f(v) \rangle = \sum_{k,l} x_k \overline{y_l} \langle f(e_k), f(e_l) \rangle = \sum_k x_k \overline{y_k} = \langle u, v \rangle, \]
so that $f$ is orthogonal.
\end{proof}

\subsection{Spectral Properties of Self-Adjoint Operators}

\begin{teorema}{Real Eigenvalues}{autovalores_reales}
Let $(V, \langle \cdot, \cdot \rangle)$ be an inner product space and $f$ a self-adjoint endomorphism.
Then all the eigenvalues of $f$ are real numbers.
\end{teorema}

\begin{proof}
Let $\lambda$ be an eigenvalue of $f$ and $v \neq 0$ an associated eigenvector, $f(v) = \lambda v$. By linearity in the first variable, $\langle f(v), v \rangle = \lambda \|v\|^2$; and since $f$ is self-adjoint, $\langle f(v), v \rangle = \langle v, f(v) \rangle = \overline{\lambda} \|v\|^2$. Equating, $\lambda \|v\|^2 = \overline{\lambda} \|v\|^2$, and since $\|v\| > 0$, one has $\lambda = \overline{\lambda}$, that is, $\lambda \in \mathbb{R}$.
\end{proof}

\begin{lema}\label{lema_invariancia}
Let $f$ be a self-adjoint endomorphism on $V$ and let $S \subseteq V$ be a nonzero subspace invariant under $f$. Then:
\begin{enumerate}
    \item[i)] $S^\perp$ is also invariant under $f$.
    \item[ii)] $S$ contains at least one eigenvector of $f$.
\end{enumerate}
\end{lema}

\begin{proof}
\textit{(i).} Let $y \in S^\perp$; it is proved that $f(y) \in S^\perp$, that is, that $f(y)$ is orthogonal to every $x \in S$. Since $f$ is self-adjoint, $\langle x, f(y) \rangle = \langle f(x), y \rangle$, and since $S$ is invariant, $f(x) \in S$, so that $\langle f(x), y \rangle = 0$ because $y \in S^\perp$. Since this holds for every $x \in S$, $f(y) \perp S$, that is, $f(y) \in S^\perp$.

\textit{(ii).} Consider the restriction $f|_S: S \to S$, well defined because $S$ is invariant, and self-adjoint because $f$ is. By the Fundamental Theorem of Algebra, the characteristic polynomial of $f|_S$ has a root $\lambda$ (considering the associated matrix as a complex matrix), which by Theorem \ref{teo:autovalores_reales} is real. Then there is $v \in S$, $v \neq 0$, with $f|_S(v) = \lambda v$, and since $f|_S(v) = f(v)$, $v$ is an eigenvector of $f$ in $S$.
\end{proof}

\begin{proposicion}{Orthogonality of Eigenspaces}{ortogonalidad_autoespacios}
Let $f$ be a self-adjoint endomorphism. If $\alpha$ and $\beta$ are two distinct eigenvalues of $f$ ($\alpha \neq \beta$), then their associated eigenspaces are orthogonal to each other.
\end{proposicion}

\begin{proof}
Let $u \in S_\alpha$ and $v \in S_\beta$, with $f(u) = \alpha u$ and $f(v) = \beta v$. Compute $\langle f(u), v \rangle$ in two ways: on the one hand, $\langle f(u), v \rangle = \alpha \langle u, v \rangle$; on the other, since $f$ is self-adjoint and its eigenvalues are real ($\overline{\beta} = \beta$), $\langle f(u), v \rangle = \langle u, f(v) \rangle = \beta \langle u, v \rangle$. Equating, $(\alpha - \beta) \langle u, v \rangle = 0$, and since $\alpha \neq \beta$, $\langle u, v \rangle = 0$. Hence $S_\alpha \perp S_\beta$.
\end{proof}

\obs This result is useful in practice: to orthogonally diagonalize a symmetric matrix it suffices to find a basis of each eigenspace separately, since the eigenspaces of distinct eigenvalues are already orthogonal to each other; it only remains to orthonormalize (with Gram-Schmidt) within each eigenspace.

\obs The previous property ($S_\alpha \perp S_\beta$) does not hold in general.
    
    For an arbitrary endomorphism, if $\alpha \neq \beta$, the only thing that can be guaranteed is that the eigenvectors are linearly independent (they form a basis, if there are enough of them), but they may fail to form a 90-degree angle.
    
    The condition of being self-adjoint is what forces these stretching axes to be perpendicular to each other.

\subsection{Spectral Theorem}

\begin{teorema}{Spectral Decomposition}{descomposicion_espectral}
Let $f$ be a self-adjoint endomorphism on a finite-dimensional space $V$.
Then $V$ is the orthogonal direct sum of the eigenspaces of $f$:
\[ V = S_{\lambda_1} \oplus^\perp S_{\lambda_2} \oplus^\perp \dots \oplus^\perp S_{\lambda_k} \]
\end{teorema}

\begin{proof}
The strategy is to extract orthogonal eigenspaces successively. By the Fundamental Theorem of Algebra and since $f$ is self-adjoint, there is a real eigenvalue $\lambda_1$; let $W_1 = S_{\lambda_1} = \ker(f - \lambda_1 I)$, invariant under $f$. By Lemma \ref{lema_invariancia}, $W_1^\perp$ is also invariant. If $W_1^\perp = \{0\}$, then $V = W_1$ and the decomposition consists of a single eigenspace.

If $W_1^\perp \neq \{0\}$, the restriction $f|_{W_1^\perp}$ is self-adjoint, so that it has an eigenvalue $\lambda_2$ with an eigenvector in $W_1^\perp$. Necessarily $\lambda_2 \neq \lambda_1$: if $\lambda_2 = \lambda_1$, the eigenvector would be in $S_{\lambda_1} \cap S_{\lambda_1}^\perp = \{0\}$. Let $W_2 = S_{\lambda_2}$; since $\lambda_1 \neq \lambda_2$, $W_1 \perp W_2$. The sum $W_1 \oplus W_2$ is invariant and, by Lemma \ref{lema_invariancia}, so is its orthogonal complement, so that the process is repeated on it.

Since $V$ is finite-dimensional, the process ends in finitely many steps and one obtains
\[ V = S_{\lambda_1} \oplus^\perp S_{\lambda_2} \oplus^\perp \dots \oplus^\perp S_{\lambda_k}. \]
\end{proof}

\obs The spectral theorem, in its general form, is stated in Hilbert spaces.

\subsection{Corollary: Orthogonal Diagonalization}

\begin{corolario}{Diagonalization by an Orthogonal Matrix}{diag_ortogonal}
Let $A \in \mathbb{R}^{n \times n}$ be a symmetric matrix.
Then there are an orthogonal matrix $P$ and a diagonal matrix $D$ such that:
\[ P^T A P = D \]
\end{corolario}

\begin{proof}
By Theorem \ref{teo:descomposicion_espectral}, $V = S_{\lambda_1} \oplus^\perp \dots \oplus^\perp S_{\lambda_k}$. Take an orthonormal basis $\mathcal{B}_i$ of each $S_{\lambda_i}$ and let
\[ \mathcal{B} = \mathcal{B}_1 \cup \dots \cup \mathcal{B}_k = \{ v_1, \dots, v_n \}, \]
which is an orthonormal basis of $V$. The matrix $P$ with columns $v_1, \dots, v_n$ is orthogonal, hence $P^{-1} = P^T$. The $(i,j)$ entry of $M = P^T A P$ is
\[ m_{ij} = \langle v_i, A v_j \rangle, \]
interpreting the matrix product as the inner product of rows and columns. Since $v_j$ is an eigenvector, $A v_j = \lambda_j v_j$, so that
\[ m_{ij} = \lambda_j \langle v_i, v_j \rangle = \lambda_j \delta_{ij}. \]
Therefore $M$ is diagonal with $m_{ii} = \lambda_i$, that is,
\[ P^T A P = \begin{pmatrix} \lambda_1 & 0 & \dots & 0 \\ 0 & \lambda_2 & \dots & 0 \\ \vdots & \vdots & \ddots & \vdots \\ 0 & 0 & \dots & \lambda_n \end{pmatrix} = D. \]
\end{proof}

\obs The complex analogue also holds: if $A \in \mathbb{C}^{n \times n}$ is Hermitian, there is a unitary matrix $U$ such that $U^* A U$ is diagonal (with real entries); the proof is the same, replacing transposes by conjugate transposes.

\subsection{Cayley-Hamilton Theorem}

\begin{teorema}{Cayley-Hamilton}{cayley_hamilton}
Let $A \in \mathbb{K}^{n \times n}$ be a diagonalizable matrix.
Let $P(\lambda) = \det(A - \lambda I)$ be its characteristic polynomial.
Then the matrix satisfies its own characteristic equation:
\[ P(A) = 0 \]
Where $0$ represents the zero matrix of size $n \times n$.
\end{teorema}

\begin{proof}
Since $A$ is diagonalizable, $A = C D C^{-1}$ with $D$ diagonal whose entries are the eigenvalues $\lambda_1, \dots, \lambda_n$. If $P(t) = a_n t^n + \dots + a_0$, using $(C D C^{-1})^k = C D^k C^{-1}$,
\[ P(A) = a_n C D^n C^{-1} + \dots + a_0 C I C^{-1} = C\, P(D)\, C^{-1}. \]
Since $D$ is diagonal, $P(D)$ is diagonal with entries $P(\lambda_i)$, and since the $\lambda_i$ are the roots of $P$, $P(\lambda_i) = 0$ for every $i$; hence $P(D) = 0$ and
\[ P(A) = C \cdot 0 \cdot C^{-1} = 0. \]
\end{proof}

\obs The Cayley-Hamilton Theorem holds for any square matrix, whether diagonalizable or not. The proof in the general case requires more advanced tools, such as the Jordan canonical form or continuity arguments, which are not developed here.

\newpage

\section{Application: Optimization of Quadratic Forms}

\subsection{Quadratic Forms}

Every homogeneous polynomial of degree 2 in $m$ variables ($x_1, \dots, x_m$) can be written compactly using matrix notation.
Let $f(x) = \sum_{i,j=1}^m a_{ij} x_i x_j$. Define the symmetric coefficient matrix $A$ and the vector $x = (x_1, \dots, x_m)^T$ such that:
\[ f(x) = x^T A x = \langle x, Ax \rangle \]
This function $f: \mathbb{R}^m \to \mathbb{R}$ is called the \textbf{quadratic form} associated with $A$.

\subsection{The Optimization Problem}

Consider the following constrained extremum problem:
\begin{quote}
    Find the maximum and minimum values of the function $f(x) = x^T A x$ subject to the constraint that $x$ belongs to the unit sphere $S^{m-1}$.
\end{quote}
Formally:
\[ \text{Optimize } f(x) = \langle x, Ax \rangle \quad \text{subject to } \quad g(x) = \|x\|^2 = 1 \]

\subsection{Solution by Lagrange Multipliers}

To find the critical points, the method of Lagrange multipliers is used: one seeks $x$ and a scalar $\lambda$ such that the gradients of the objective function and of the constraint are parallel:
\[ \nabla f(x) = \lambda \nabla g(x) \]

The gradients are:
\begin{enumerate}
    \item $\nabla f(x) = \nabla (x^T A x) = 2Ax$ (here it is used that $A$ is symmetric),
    \item $\nabla g(x) = \nabla (x^T x) = 2x$.
\end{enumerate}

Substituting in the Lagrange equation:
\[ 2Ax = \lambda (2x) \implies Ax = \lambda x \]

\obs The necessary condition for a point $x$ to be an extremum on the sphere is that $x$ be an eigenvector of the matrix $A$, and the Lagrange multiplier $\lambda$ be its associated eigenvalue.

\subsection{Analysis of the Extreme Values}

The objective function $f$ is evaluated at these critical points. Let $x$ be a unit eigenvector ($x^T x = 1$) with associated eigenvalue $\lambda$.
\[ f(x) = x^T A x = x^T (\lambda x) = \lambda (x^T x) = \lambda \cdot 1 = \lambda \]

This leads to the following theorem:

\begin{teorema}{Extrema on the Sphere (Rayleigh Quotient Theorem)}{extremos_esfera}
Let $A \in \mathbb{R}^{m \times m}$ be a symmetric matrix with real eigenvalues ordered $\lambda_1 \geq \lambda_2 \geq \dots \geq \lambda_m$.
Then, on the unit sphere $\|x\|=1$:
\begin{enumerate}
    \item The absolute maximum value of the quadratic form is the largest eigenvalue, $\lambda_{\max} = \lambda_1$, and it is attained at the associated eigenvector $v_1$.
    \item The absolute minimum value is the smallest eigenvalue, $\lambda_{\min} = \lambda_m$, and it is attained at the associated eigenvector $v_m$.
\end{enumerate}
\end{teorema}

\begin{proof}
The unit sphere is compact and $f$ is continuous, so that the maximum and the minimum are attained. By the previous analysis (Lagrange multipliers), every extremum is attained at a unit eigenvector $x$ of $A$, and there $f(x) = \lambda$, its associated eigenvalue. The maximum is then the largest of the eigenvalues, $\lambda_1$, attained at $v_1$, and the minimum is the smallest, $\lambda_m$, attained at $v_m$.
\end{proof}

\begin{corolario}{Orthogonality of the Extrema}{ortogonalidad_extremos}
If the maximum and minimum eigenvalues are distinct, then the points where the maximum and the minimum are attained are orthogonal to each other.
\end{corolario}

\begin{proof}
The maximum is attained at $v_1$ (eigenvector of $\lambda_1$) and the minimum at $v_m$ (eigenvector of $\lambda_m$). Since $A$ is symmetric, by Proposition \ref{prop:ortogonalidad_autoespacios} the eigenvectors associated with distinct eigenvalues are orthogonal, so that $v_1 \perp v_m$.
\end{proof}

\newpage

\section{Singular Values: Generalization of Diagonalization}

\subsection{Motivation}

The previous section dealt with diagonalization ($A = P D P^{-1}$), which has two limitations:
\begin{enumerate}
    \item It applies only to square matrices ($A \in \mathbb{K}^{n \times n}$).
    \item Not every square matrix is diagonalizable (it depends on having a complete basis of eigenvectors).
\end{enumerate}

If $A \in \mathbb{R}^{m \times n}$ is rectangular and represents a linear transformation between spaces of different dimension ($f: \mathbb{R}^n \to \mathbb{R}^m$), classical diagonalization does not apply, because the input and output vectors live in different spaces. Singular values solve this: instead of a single basis that diagonalizes the transformation, two orthonormal bases are sought, one for the domain and another for the codomain.

\subsection{Definitions and Main Theorem}

To find these bases, the properties of self-adjoint operators already proved are used. Given any matrix $A \in \mathbb{R}^{m \times n}$, the matrix $A^T A \in \mathbb{R}^{n \times n}$ is always symmetric (and hence self-adjoint) and positive semidefinite: $v^T (A^T A) v = \|Av\|^2 \ge 0$ for every $v$; in particular, its eigenvalues are nonnegative (if $A^T A v = \lambda v$ with $v \neq 0$, then $\lambda \|v\|^2 = \|Av\|^2 \ge 0$).

\begin{definicion}{Singular Values}{valores_singulares}
Let $A \in \mathbb{R}^{m \times n}$. Consider the symmetric matrix $A^T A$. By the Spectral Theorem, its eigenvalues $\lambda_1, \lambda_2, \dots, \lambda_n$ are real and nonnegative ($\lambda_i \ge 0$).

The \textbf{singular values} of $A$, conventionally denoted by $\sigma_i$, are defined as the positive square roots of the eigenvalues of $A^T A$:
\[ \sigma_i = \sqrt{\lambda_i} \]
They are usually ordered from largest to smallest: $\sigma_1 \ge \sigma_2 \ge \dots \ge \sigma_r > 0$, where $r$ is the number of nonzero singular values.
\end{definicion}

\begin{teorema}{Singular Value Decomposition (SVD)}{svd}
Every matrix $A \in \mathbb{R}^{m \times n}$ can be factorized in the form:
\[ A = U \Sigma V^T \]
Where:
\begin{itemize}
    \item $U \in \mathbb{R}^{m \times m}$ is an orthogonal matrix ($U^T U = I$). Its columns are called \textbf{left singular vectors}.
    \item $V \in \mathbb{R}^{n \times n}$ is an orthogonal matrix ($V^T V = I$). Its columns are called \textbf{right singular vectors}.
    \item $\Sigma \in \mathbb{R}^{m \times n}$ is a rectangular diagonal matrix, whose entries on the main diagonal are the singular values $\sigma_i$ in decreasing order, and whose remaining entries are zero.
\end{itemize}
\end{teorema}

\begin{proof}
The proof is constructive: $V$, $\Sigma$ and $U$ are constructed.

Consider $A^T A \in \mathbb{R}^{n \times n}$, symmetric and positive semidefinite. By Corollary \ref{cor:diag_ortogonal} there is an orthonormal basis $\{v_1, \dots, v_n\}$ of $\mathbb{R}^n$ formed by eigenvectors of $A^T A$, with eigenvalues $\lambda_i \ge 0$. With them ordered from largest to smallest and $\sigma_i = \sqrt{\lambda_i}$, let $r$ be the rank of $A$, which coincides with the number of nonzero eigenvalues: $\ker(A^T A) = \ker(A)$ (if $A^T A v = 0$, then $\|Av\|^2 = v^T A^T A v = 0$, and the other inclusion is direct), so that
\[ \sigma_1 \ge \dots \ge \sigma_r > 0 \quad \text{and} \quad \sigma_{r+1} = \dots = \sigma_n = 0. \]
The matrix $V = (v_1 \mid \dots \mid v_n)$ is orthogonal because it has orthonormal columns.

For $i = 1, \dots, r$ define $u_i = \frac{1}{\sigma_i} A v_i \in \mathbb{R}^m$. This set is orthonormal: using $A^T A v_j = \sigma_j^2 v_j$,
\[ \langle u_i, u_j \rangle = \frac{1}{\sigma_i \sigma_j} v_i^T A^T A v_j = \frac{\sigma_j^2}{\sigma_i \sigma_j}\, v_i^T v_j = \delta_{ij}. \]
Extending $\{u_1, \dots, u_r\}$ to a basis of $\mathbb{R}^m$ and orthonormalizing with Gram-Schmidt, it is completed to an orthonormal basis $\{u_1, \dots, u_m\}$ of $\mathbb{R}^m$, and $U = (u_1 \mid \dots \mid u_m)$ is orthogonal. Let $\Sigma \in \mathbb{R}^{m \times n}$ be the matrix with the $\sigma_i$ on the main diagonal and zeros elsewhere.

Since $V$ is orthogonal, $A = U \Sigma V^T$ is equivalent to $A V = U \Sigma$, which is verified column by column: column $i$ of $AV$ is $A v_i$ and that of $U \Sigma$ is $\sigma_i u_i$.
\begin{itemize}
    \item For $i \le r$, from $u_i = \frac{1}{\sigma_i} A v_i$ one obtains $A v_i = \sigma_i u_i$.
    \item For $i > r$, $\sigma_i = 0$ and $\lambda_i = 0$, so that $\|A v_i\|^2 = v_i^T A^T A v_i = \lambda_i \|v_i\|^2 = 0$, that is, $A v_i = 0 = \sigma_i u_i$.
\end{itemize}
Hence $A V = U \Sigma$ and, multiplying by $V^T$ on the right, $A = U \Sigma V^T$.
\end{proof}

\obs \textit{Geometric Intuition:}

The SVD theorem shows that any matrix linear transformation $A$ can be decomposed into three sequential geometric steps:
\begin{enumerate}
    \item $V^T$: A rotation (or reflection) in the input space $\mathbb{R}^n$.
    \item $\Sigma$: A scaling along the coordinate axes (which moreover changes the dimension from $n$ to $m$). The singular values $\sigma_i$ indicate exactly how much the space is ``stretched'' in each of those orthogonal directions.
    \item $U$: A new rotation (or reflection) in the output space $\mathbb{R}^m$.
\end{enumerate}
This generalizes diagonalization, where the rotation there and back was done with the same matrix ($P$ and $P^{-1}$), restricted to a single space.

\subsection{SVD as a Generalization of Eigenvalues}

\begin{proposicion}{Relation between SVD and Eigenvectors}{relacion_svd_autovectores}
Let $A = U \Sigma V^T$ be the Singular Value Decomposition of $A$. Then:
\begin{enumerate}
    \item[i)] The columns of $V$ are the orthonormal eigenvectors of $A^T A$.
    \item[ii)] The columns of $U$ are the orthonormal eigenvectors of $A A^T$.
\end{enumerate}
\end{proposicion}

\begin{proof}
\textit{(i).} Substituting the SVD,
\[ A^T A = (U \Sigma V^T)^T (U \Sigma V^T) = V \Sigma^T U^T U \Sigma V^T = V (\Sigma^T \Sigma) V^T, \]
since $U^T U = I_m$. The matrix $\Sigma^T \Sigma \in \mathbb{R}^{n \times n}$ is diagonal with entries $\sigma_i^2 = \lambda_i$ (Definition \ref{def:valores_singulares}). Multiplying by $V$ on the right and using $V^T V = I_n$, $(A^T A) V = V (\Sigma^T \Sigma)$, and column by column $(A^T A) v_j = \sigma_j^2 v_j = \lambda_j v_j$; hence each $v_j$ is an eigenvector of $A^T A$.

\textit{(ii).} Analogously,
\[ A A^T = (U \Sigma V^T)(U \Sigma V^T)^T = U \Sigma V^T V \Sigma^T U^T = U (\Sigma \Sigma^T) U^T, \]
with $\Sigma \Sigma^T$ diagonal with entries $\sigma_i^2$. The same reasoning shows that the columns of $U$ are eigenvectors of $A A^T$.
\end{proof}

\newpage

\section{Principal Component Analysis}

\subsection{The Problem of Dimensionality}

Consider a large data set: for example, a database with $m = 1000$ users, each with $n = 50$ different variables (age, height, weight, income and connection time, among others).

Processing a $1000 \times 50$ matrix is costly, and 50 dimensions do not admit direct visualization. Moreover, many of these variables are usually correlated (for example, the number of clicks and the time on the page tend to grow together), which indicates redundant information.

\textbf{Principal Component Analysis} (PCA) seeks to reduce those 50 variables to a few artificial variables that summarize almost all the original information. This compresses the data, speeds up the algorithms and allows plotting them.

\subsection{The Geometric Intuition}

Consider a simpler case, with only two variables, height and weight of several people. In a 2D plane, the data form a cloud of points elongated along a diagonal, since greater height usually corresponds to greater weight.

To summarize the position of each person with a single number (reducing from 2D to 1D), projecting onto the $X$ axis (ignoring weight) flattens the points and loses much of the information about overall size; the same happens when projecting onto the $Y$ axis.

The idea of PCA is to rotate the axes: to draw a new line that crosses the cloud of points along its maximum extent.
\begin{itemize}
    \item This new axis is called the \textbf{first principal component} (PC1); along it the data are as spread out as possible. In statistics, spread is measured by the variance, which is interpreted as information.
    \item The second axis (PC2) is drawn perpendicular to the first one and captures the residual variance.
\end{itemize}

For example, if PC1 (interpretable as the overall size of the person) contains 95\% of the variance, PC2 can be discarded, going from 2 to 1 dimension while losing only 5\% of the information.

\subsection{Formalization: Variance and Covariance}

To translate this intuition into mathematics, the data matrix $X \in \mathbb{R}^{m \times n}$ is assumed to be centered (the mean has been subtracted from each column, so that the origin of coordinates is at the center of the cloud of points).

A direction is sought, given by a unit vector $w \in \mathbb{R}^n$ ($\|w\| = 1$), onto which to project the data so that the variance is maximal.

\begin{definicion}{Covariance Matrix}{matriz_covarianza}
Given the centered data matrix $X \in \mathbb{R}^{m \times n}$, the sample \textbf{covariance matrix} $C \in \mathbb{R}^{n \times n}$ is defined as:
\[ C = \frac{1}{m-1} X^T X \]
The entry $C_{ij}$ indicates how variable $i$ and variable $j$ vary together. The main diagonal $C_{ii}$ contains the individual variance of each variable.
\end{definicion}

\obs It is the natural generalization to higher dimensions of the concept of variance of a scalar random variable.

\obs The matrix $C$ is, by construction, symmetric and positive semidefinite, which allows using the theory of self-adjoint operators already developed.

The projection of the $m$ data points onto the vector $w$ is the product $Xw$. The variance of these projected data is a quadratic form:
\[ \text{Var}(Xw) = \frac{1}{m-1} (Xw)^T (Xw) = \frac{1}{m-1} w^T X^T X w = w^T C w \]

\subsection{PCA as Optimization of a Quadratic Form}

The goal of finding the axis of largest spread is now a precise mathematical problem: maximize the quadratic form $f(w) = w^T C w$ subject to the vector being unit ($\|w\|^2 = 1$).

This is exactly the optimization problem on the unit sphere that was solved in the previous section.

\begin{proposicion}{Principal Components}{pca_autovectores}
The orthogonal directions that successively maximize the variance of the data (the principal components) are exactly the eigenvectors of the covariance matrix $C$. The amount of information (variance) captured by each component is exactly its associated eigenvalue.
\end{proposicion}

\begin{proof}
By the Rayleigh Quotient Theorem \ref{teo:extremos_esfera}, the maximum value of the quadratic form $w^T C w$ on the unit sphere is attained at the eigenvector $v_1$ associated with the largest eigenvalue $\lambda_1$ of $C$, which is the first principal component. The second component is obtained by maximizing the variance over the directions perpendicular to $v_1$; by Corollary \ref{cor:ortogonalidad_extremos}, that maximum is the second largest eigenvalue $\lambda_2$, attained at the orthogonal eigenvector $v_2$. The process is repeated until the $n$ dimensions are exhausted.
\end{proof}

\subsection{PCA via SVD}

The direct construction of PCA goes through forming the matrix $C = \frac{1}{m-1} X^T X$ and finding its eigenvectors. In modern computational practice, however, this is rarely done that way.

Computing $X^T X$ when $X$ is large is slow and can amplify rounding errors. The Singular Value Decomposition offers a numerically stable alternative.

\begin{teorema}{Equivalence between SVD and PCA}{equivalencia_svd_pca}
Let $X \in \mathbb{R}^{m \times n}$ be a centered data matrix and let $X = U \Sigma V^T$ be its Singular Value Decomposition. Then:
\begin{enumerate}
    \item[i)] The principal components of $X$ are exactly the columns of the orthogonal matrix $V$ (the right singular vectors).
    \item[ii)] The variances of the data along each component (the eigenvalues $\lambda_i$ of $C$) are given by:
    \[ \lambda_i = \frac{\sigma_i^2}{m-1} \]
    where $\sigma_i$ are the singular values of $X$ (the diagonal of $\Sigma$).
\end{enumerate}
\end{teorema}

\begin{proof}
Substituting the SVD $X = U \Sigma V^T$ into the covariance matrix and using $U^T U = I_m$,
\[ C = \frac{1}{m-1} (U \Sigma V^T)^T (U \Sigma V^T) = V \left( \frac{\Sigma^T \Sigma}{m-1} \right) V^T. \]
Since $V$ is orthogonal, $V^T = V^{-1}$, so that
% Correction: it said Σ² inside the formula, but Σ is rectangular and cannot be squared;
% the correct diagonal matrix is Σ^T Σ, as in the previous line.
\[ C = V \left( \frac{\Sigma^T \Sigma}{m-1} \right) V^{-1} \]
is the diagonalization of $C$. Therefore, the columns of $V$ are the eigenvectors of $C$ (the principal components) and the eigenvalues are $\lambda_i = \frac{\sigma_i^2}{m-1}$.
\end{proof}

\obs To apply PCA it suffices to pass the matrix $X$ to an SVD algorithm: the columns of $V$ give the rotation of the axes and the singular values of $\Sigma$ indicate how much information each axis retains, which allows discarding the least relevant ones and compressing the data.

\subsection{Properties of the Transformation}

Principal Component Analysis reduces dimensions and, moreover, transforms the statistical nature of the data. Projecting the matrix $X$ onto the basis of eigenvectors $V$ yields a new set of transformed data $Y = XV$.

This linear transformation has two notable properties:

\begin{proposicion}{Decorrelation of the Data}{decorrelacion}
The new variables obtained through PCA (the principal components) are completely decorrelated from each other. That is, the covariance matrix of the transformed data $Y$ is a diagonal matrix.
\end{proposicion}

\begin{proof}
The covariance matrix of $Y = XV$ is
\[ C_Y = \frac{1}{m-1} Y^T Y = \frac{1}{m-1} (XV)^T (XV) = V^T \left( \frac{1}{m-1} X^T X \right) V = V^T C_X V. \]
Since the columns of $V$ are the eigenvectors of $C_X$, the matrix $V$ diagonalizes $C_X$, so that $C_Y = V^T C_X V = D$ is diagonal with the eigenvalues $\lambda_i$. Being diagonal, the cross covariances are zero and the data are decorrelated.
\end{proof}

\obs The basis of eigenvectors provided by PCA (or SVD) is optimal for the reconstruction of the information: to approximate a data set of dimension $n$ by a subspace of lower dimension $k$ (with $k < n$), the Eckart-Young-Mirsky theorem guarantees that projecting onto the first $k$ principal components minimizes the mean squared reconstruction error. No other orthogonal basis retains more variance (information) in so few dimensions.

\obs Since this mathematical tool is universal, it receives different names depending on the discipline that applies it:
\begin{itemize}
    \item In Statistics and data science, it is known as Principal Component Analysis (PCA).
    \item In Signal and Image Processing, it is known as the Karhunen-Loève Transform (KLT). It is used extensively for image compression: adjacent pixels of an image are highly correlated. The KLT decorrelates these pixels, concentrating most of the visual energy of the image in a few singular coefficients, which allows discarding the rest without perceptible loss of quality.
\end{itemize}

\newpage

\section{Matrices in the Computer}

\subsection{Representation in Memory}

The memory of a computer is a linear array: the cells are identified with consecutive addresses $0, 1, 2, \dots$ To store a matrix, which is a two-dimensional object, one must choose the order in which its entries are laid out on that line.

\begin{definicion}{Row-Major and Column-Major Orders}{row_column_major}
Let $A \in \mathbb{K}^{m \times n}$.
\begin{itemize}
    \item In \textbf{row-major} order (by rows), the entries are stored row by row: the entry $a_{ij}$ (with indices starting at $1$) occupies the position
    \[
    (i-1) \, n + (j-1)
    \]
    within the memory block (counting from $0$).
    \item In \textbf{column-major} order (by columns), they are stored column by column: $a_{ij}$ occupies the position $(j-1) \, m + (i-1)$.
\end{itemize}
\end{definicion}

\obs The choice is a convention of each language: C and its derivatives use row-major; Fortran, MATLAB and R use column-major. The distinction matters in practice: traversing the matrix in the same order in which it is stored produces accesses to consecutive addresses, which the hardware serves much faster than scattered accesses (locality of reference). The same algorithm can change its real running time by a large factor just because of the order of its loops, even though its number of arithmetic operations is identical.

A \textbf{block partition} of $A \in \mathbb{K}^{m \times n}$ is obtained by cutting its rows into consecutive stretches of sizes $m_1, \dots, m_q$ (with $m_1 + \dots + m_q = m$) and its columns into consecutive stretches of sizes $n_1, \dots, n_r$ (with $n_1 + \dots + n_r = n$): the matrix is divided into $q \cdot r$ \textbf{blocks} $A_{IJ} \in \mathbb{K}^{m_I \times n_J}$, where $A_{IJ}$ is the submatrix formed by the rows of stretch $I$ and the columns of stretch $J$. There are two ways of materializing a partition:
\begin{enumerate}
    \item \textit{By copying:} memory is reserved for each block and its entries are copied. It costs time proportional to the number of copied entries, $\Theta(mn)$.
    \item \textit{By index computation:} a block is described by three pieces of data (the address of its upper-left corner, its dimensions and the width of the original matrix, the \emph{stride}, which separates two consecutive rows of the block in memory), without touching any entry. Partitioning costs $\Theta(1)$, and every modification of the block updates the original matrix, since they occupy the same memory.
\end{enumerate}
In what follows, partitions are assumed to be done by index computation, in time $\Theta(1)$.

The usefulness of blocks comes from the fact that the matrix product respects them:

\begin{proposicion}{Block Product}{producto_bloques}
Let $A \in \mathbb{K}^{m \times n}$ and $B \in \mathbb{K}^{n \times p}$, partitioned in a \emph{conformal} way: the partition of the columns of $A$ and that of the rows of $B$ are the same, $n = n_1 + \dots + n_r$. Let $C = AB$, partitioned according to the rows of $A$ ($m = m_1 + \dots + m_q$) and the columns of $B$ ($p = p_1 + \dots + p_s$). Then, for all $I \in \{1, \dots, q\}$ and $J \in \{1, \dots, s\}$,
\[
C_{IJ} = \sum_{K=1}^{r} A_{IK} \, B_{KJ},
\]
where each product is well defined: $A_{IK} \in \mathbb{K}^{m_I \times n_K}$ and $B_{KJ} \in \mathbb{K}^{n_K \times p_J}$. That is, the blocks are multiplied with the same ``row by column'' rule as the entries.
\end{proposicion}

\begin{proof}
Let $i$ be a row of stretch $I$ and $j$ a column of stretch $J$, and let $t_K = n_1 + \dots + n_K$ (with $t_0 = 0$), so that stretch $K$ of the common partition consists of the indices $t_{K-1} + 1, \dots, t_K$. Splitting the sum in the definition of the product according to those stretches,
\[
c_{ij} = \sum_{k=1}^{n} a_{ik} \, b_{kj} = \sum_{K=1}^{r} \; \sum_{k = t_{K-1}+1}^{t_K} a_{ik} \, b_{kj}.
\]
For each $K$, the inner sum runs over the columns of $A$ of stretch $K$ against the rows of $B$ of the same stretch: it is the ``row by column'' rule applied to the row of $A_{IK}$ corresponding to $i$ and to the column of $B_{KJ}$ corresponding to $j$, that is, the entry of $A_{IK} B_{KJ}$ at the relative position that $(i, j)$ occupies within block $(I, J)$. Summing over $K$ yields the corresponding entry of $\sum_{K} A_{IK} B_{KJ}$, as claimed.
\end{proof}

The case used by divide and conquer algorithms is the square one with all stretches equal: $A, B \in \mathbb{K}^{n \times n}$ with $n$ even (a power of $2$ if partitioned recursively) and $q = r = s = 2$ with stretches of size $n/2$,
\[
A = \begin{pmatrix} A_{11} & A_{12} \\ A_{21} & A_{22} \end{pmatrix},
\qquad
\begin{aligned}
C_{11} &= A_{11} B_{11} + A_{12} B_{21}, & C_{12} &= A_{11} B_{12} + A_{12} B_{22}, \\
C_{21} &= A_{21} B_{11} + A_{22} B_{21}, & C_{22} &= A_{21} B_{12} + A_{22} B_{22}.
\end{aligned}
\]

\subsection{The Cost of Multiplication}

Computing $C = AB$ directly from the definition requires, for each of the $n^2$ entries $c_{ij}$, a sum of $n$ products: in total $\Theta(n^3)$ arithmetic operations, organized in three nested loops over $i$, $j$ and $k$.

Proposition \ref{prop:producto_bloques} suggests a divide and conquer algorithm: compute the four blocks of $C$ through $8$ multiplications of $n/2 \times n/2$ matrices, plus the block additions. With partitions by index computation, its time satisfies the recurrence
\[
T(n) = 8 \, T(n/2) + \Theta(1),
\]
whose solution is $\Theta(n^3)$ (the recurrence is solved with the master method; see the \emph{Foundations} notes of the Algorithms corpus): dividing into blocks, by itself, does not improve the asymptotic cost, because the number of subproblems ($8$) exactly compensates for the reduction in size.

\textbf{Strassen's algorithm} obtains, through a combination of additions and subtractions of blocks, the four blocks of $C$ with only $7$ multiplications of size $n/2$, at the price of a constant number of block operations of cost $\Theta(n^2)$. Applied recursively, its time satisfies
\[
T(n) = 7 \, T(n/2) + \Theta(n^2),
\]
whose solution is $\Theta(n^{\lg 7}) = O(n^{2.81})$: asymptotically better than $\Theta(n^3)$. The seven formulas and their verification appear in the \emph{Foundations} notes of the Algorithms corpus; for these notes their existence and cost suffice. In practice it only pays off for large matrices, because of the hidden constants and a less regular use of memory (see the remark on locality in the previous subsection).

\obs The optimal exponent of matrix multiplication, denoted $\omega$, is an open problem: it is known that $2 \le \omega$ and, through successive refinements of Strassen's idea, that $\omega \le 2.372$ (results cited without proof). It is conjectured that $\omega = 2$.

\subsection{Linear Systems and LU Factorization}

The central computational problem of linear algebra is solving $Ax = b$ with $A \in \mathbb{K}^{n \times n}$ invertible. Cramer's Rule (Proposition \ref{prop:cramer}) solves it in theory, but requires computing $n+1$ determinants; the efficient route is to factorize $A$ into triangular matrices, because triangular systems are solved by reading the equations in order: the first involves a single unknown, the second two (one of them already known), and so on.

A lower triangular matrix is \textbf{unit triangular} if its diagonal consists of ones.

\begin{proposicion}{Forward and Back Substitution}{sustituciones}
Let $L \in \mathbb{K}^{n \times n}$ be unit lower triangular and $U \in \mathbb{K}^{n \times n}$ upper triangular with nonzero diagonal.
\begin{enumerate}
    \item The system $Ly = b$ has a unique solution, computable in the order $i = 1, \dots, n$ (\textbf{forward substitution}):
    \[
    y_i = b_i - \sum_{j=1}^{i-1} l_{ij} \, y_j.
    \]
    \item The system $Ux = y$ has a unique solution, computable in the order $i = n, \dots, 1$ (\textbf{back substitution}):
    \[
    x_i = \frac{1}{u_{ii}} \Big( y_i - \sum_{j=i+1}^{n} u_{ij} \, x_j \Big).
    \]
\end{enumerate}
Each substitution performs $\Theta(n^2)$ operations.
\end{proposicion}

\begin{proof}
Equation $i$ of $Ly = b$ is $\sum_{j \le i} l_{ij} y_j = b_i$ with $l_{ii} = 1$; solving for $y_i$ gives the formula of the statement, which only uses the $y_j$ with $j < i$, already computed. Equation $i$ of $Ux = y$ is $\sum_{j \ge i} u_{ij} x_j = y_i$ with $u_{ii} \neq 0$; solving for $x_i$ gives the formula, which only uses the $x_j$ with $j > i$. In both cases each value is uniquely determined, and equation $i$ costs $O(n)$ operations: $\Theta(n^2)$ in total.
\end{proof}

\begin{definicion}{LU Factorization}{def_lu}
An \textbf{LU factorization} of $A \in \mathbb{K}^{n \times n}$ is a writing $A = LU$ with $L$ unit lower triangular and $U$ upper triangular.
\end{definicion}

The factorization is the matrix formulation of Gaussian elimination, and it is convenient to construct it by examining \emph{one} elimination step. Suppose $a_{11} \neq 0$ and partition
\[
A = \begin{pmatrix} a_{11} & w^T \\ v & A' \end{pmatrix},
\]
with $v \in \mathbb{K}^{n-1}$ the rest of the first column, $w^T$ the rest of the first row and $A' \in \mathbb{K}^{(n-1) \times (n-1)}$ the rest of the matrix. Eliminating the first column means subtracting, from each row $i \ge 2$, the first row multiplied by the \textbf{multiplier} $v_i / a_{11}$; with that, entry $(i,1)$ becomes $v_i - \frac{v_i}{a_{11}} \, a_{11} = 0$, and $\frac{v_i}{a_{11}} \, w^T$ is subtracted from the rest of row $i$. Stacking those corrections for all the rows, what is subtracted from the block $A'$ is the outer product $\frac{v \, w^T}{a_{11}}$ (the matrix whose row $i$ is $\frac{v_i}{a_{11}} w^T$). The result of the elimination is then
\[
\begin{pmatrix} a_{11} & w^T \\ 0 & S \end{pmatrix},
\qquad S = A' - \frac{v \, w^T}{a_{11}} \in \mathbb{K}^{(n-1) \times (n-1)},
\]
where $S$ (the matrix of the system that remains on the unknowns $2, \dots, n$ once the first one has been eliminated) is called the \textbf{Schur complement} of $A$ with respect to the \textbf{pivot} $a_{11}$.

The elimination step is undone by multiplying by the unit triangular matrix that \emph{adds back} what was subtracted:
\[
\begin{pmatrix} 1 & 0 \\ v / a_{11} & I_{n-1} \end{pmatrix}
\begin{pmatrix} a_{11} & w^T \\ 0 & S \end{pmatrix}
= \begin{pmatrix} a_{11} & w^T \\[2pt] \frac{v}{a_{11}} \, a_{11} + 0 & \frac{v}{a_{11}} \, w^T + S \end{pmatrix}
= \begin{pmatrix} a_{11} & w^T \\ v & A' \end{pmatrix} = A,
\]
where the lower right block recovers $A'$ because $\frac{v w^T}{a_{11}} + S = A'$ by the definition of $S$. This is already a factorization of $A$ with a unit lower triangular first factor; the second factor is not yet upper triangular, but the only block that prevents it is $S$, of size $n - 1$: the same process is applied to it. If recursively $S = L' U'$, with $L'$ unit lower triangular and $U'$ upper triangular, then
\[
\begin{pmatrix} a_{11} & w^T \\ 0 & S \end{pmatrix}
= \begin{pmatrix} 1 & 0 \\ 0 & L' \end{pmatrix}
\begin{pmatrix} a_{11} & w^T \\ 0 & U' \end{pmatrix}
\]
(the block product gives $(a_{11} \mid w^T)$ on top and $(0 \mid L'U') = (0 \mid S)$ below), and the two unit triangular factors combine into one:
\[
A = \begin{pmatrix} 1 & 0 \\ v/a_{11} & I_{n-1} \end{pmatrix}
\begin{pmatrix} 1 & 0 \\ 0 & L' \end{pmatrix}
\begin{pmatrix} a_{11} & w^T \\ 0 & U' \end{pmatrix}
= \underbrace{\begin{pmatrix} 1 & 0 \\ v/a_{11} & L' \end{pmatrix}}_{L}
\underbrace{\begin{pmatrix} a_{11} & w^T \\ 0 & U' \end{pmatrix}}_{U}.
\]
The structure of the result can be read directly: $U$ accumulates, row by row, the pivots with their rows as they remain after each elimination, and $L$ stores, column by column, the multipliers used to produce them ($v / a_{11}$ in the first column; inside $L'$, those of the following steps). Gaussian elimination and LU factorization are, thus, the same computation. The cost is dominated by the computation of $S$, which costs $\Theta((n-1)^2)$ operations; the total time satisfies $T(n) = T(n-1) + \Theta(n^2) = \Theta(n^3)$.

\begin{teorema}{Uniqueness of the LU Factorization}{unicidad_lu}
If $A$ is invertible, its LU factorization, when it exists, is unique.
\end{teorema}

\begin{proof}
Let $A = L_1 U_1 = L_2 U_2$ be two factorizations. Since $\det(A) = \det(L_1) \det(U_1) \neq 0$ and $\det(L_1) = 1$ (Propositions \ref{prop:producto}, \ref{prop:transpuesta} and \ref{prop:triangular}), the matrices $U_1$ and $U_2$ are invertible, and the unit triangular $L_1, L_2$ always are (determinant $1$). Multiplying the equality $L_1 U_1 = L_2 U_2$ on the left by $L_2^{-1}$ and on the right by $U_1^{-1}$,
\[
L_2^{-1} L_1 = U_2 \, U_1^{-1}.
\]
The left-hand side is unit lower triangular and the right-hand side is upper triangular (the product of two triangular matrices of the same type is triangular of that type, and the inverse of an invertible unit triangular or triangular matrix is of the same type; both claims are verified directly by induction). A matrix that is at the same time unit lower triangular and upper triangular is the identity, so that $L_2^{-1} L_1 = I_n = U_2 U_1^{-1}$, that is, $L_1 = L_2$ and $U_1 = U_2$.
\end{proof}

Not every invertible matrix admits an LU factorization: the pivot may vanish. For example,
\[
A = \begin{pmatrix} 0 & 1 \\ 1 & 1 \end{pmatrix}
\]
is invertible, but if $A = LU$, entry $(1,1)$ of the product would be $u_{11} = 0$, and then entry $(2,1)$ would be $l_{21} u_{11} = 0 \neq 1$. The solution is to permute rows before eliminating; for numerical accuracy, in practice the entry of largest absolute value in the column is chosen as the pivot, even when the natural pivot is not zero (dividing by small pivots amplifies rounding errors).

\begin{definicion}{LUP Factorization}{def_lup}
An \textbf{LUP factorization} of $A \in \mathbb{K}^{n \times n}$ is a triple $(L, U, P)$ with $L$ unit lower triangular, $U$ upper triangular and $P$ a permutation matrix (the identity with its rows reordered), such that
\[
P A = L U.
\]
\end{definicion}

Existence is proved with the same elimination as before, ensured by swaps: if the natural pivot is zero, some nonzero entry of the first column, which exists because $A$ is invertible, is moved up to the first position (by swapping two rows), and elimination proceeds normally. The only delicate point is the bookkeeping of the multipliers: the deeper levels of the recursion permute rows of the Schur complement $S$, and the multipliers already stored, which travel with those rows, must be permuted in the same way; in the proof that follows, that is exactly the role of the block $P'v/a$ inside $L$.

\begin{teorema}{Existence of the LUP Factorization}{existencia_lup}
Every invertible matrix $A \in \mathbb{K}^{n \times n}$ admits an LUP factorization.
\end{teorema}

\begin{proof}
By induction on $n$. For $n = 1$, $A = (a)$ with $a \neq 0$, and it suffices to take $P = L = (1)$, $U = (a)$.

Let $n > 1$. Since $A$ is invertible, its first column is not zero (if it were, the columns would not be linearly independent, against Proposition \ref{prop:prop_matrices_invertibles}); let $a_{k1} \neq 0$ and let $P_1$ be the permutation that swaps rows $1$ and $k$. Then
\[
P_1 A = \begin{pmatrix} a & w^T \\ v & A' \end{pmatrix} \quad \text{with } a = a_{k1} \neq 0,
\qquad
P_1 A = \begin{pmatrix} 1 & 0 \\ v/a & I_{n-1} \end{pmatrix} \begin{pmatrix} a & w^T \\ 0 & S \end{pmatrix},
\]
with $S = A' - v w^T / a$ the Schur complement. Taking determinants in the second equality and expanding the second factor along its first column (Corollary \ref{cor:calculo_det}), $\det(P_1 A) = a \cdot \det(S)$; since $\det(P_1 A) = \pm \det(A) \neq 0$, the matrix $S \in \mathbb{K}^{(n-1) \times (n-1)}$ is invertible. By the inductive hypothesis there is a factorization $P' S = L' U'$. Define
\[
P = \begin{pmatrix} 1 & 0 \\ 0 & P' \end{pmatrix} P_1, \qquad
L = \begin{pmatrix} 1 & 0 \\ P'v/a & L' \end{pmatrix}, \qquad
U = \begin{pmatrix} a & w^T \\ 0 & U' \end{pmatrix}.
\]
The matrix $P$ is a permutation matrix (a product of permutations), $L$ is unit lower triangular and $U$ is upper triangular. Multiplying by blocks,
\[
P A = \begin{pmatrix} 1 & 0 \\ 0 & P' \end{pmatrix} \begin{pmatrix} a & w^T \\ v & A' \end{pmatrix}
= \begin{pmatrix} a & w^T \\ P'v & P'A' \end{pmatrix},
\qquad
L U = \begin{pmatrix} a & w^T \\ P'v & \frac{P'v \, w^T}{a} + L'U' \end{pmatrix},
\]
and both coincide: by the inductive hypothesis, $L'U' = P'S = P' \big( A' - v w^T/a \big)$, so that $\frac{P'v \, w^T}{a} + L'U' = P'A'$.
\end{proof}

With the factorization $PA = LU$, solving $Ax = b$ reduces to two substitutions. Since $P$ is invertible (it only reorders the equations), $Ax = b$ is equivalent to $PAx = Pb$, that is, to $L(Ux) = Pb$; calling $y = Ux$, first $Ly = Pb$ is solved by forward substitution and then $Ux = y$ by back substitution, all in time $\Theta(n^2)$ (Proposition \ref{prop:sustituciones}). Factorizing costs $\Theta(n^3)$, but it is done only once: solving the system afterwards for many different vectors $b$ costs only $\Theta(n^2)$ per vector.

\begin{ejemplo}{Solving a System by LU}{ej_lu}
Let
\[
A = \begin{pmatrix} 1 & 2 & 0 \\ 3 & 4 & 4 \\ 5 & 6 & 3 \end{pmatrix}, \qquad b = \begin{pmatrix} 3 \\ 7 \\ 8 \end{pmatrix}.
\]
\textit{Factorization.} The first pivot is $a_{11} = 1$, with $v = (3, 5)^T$ and $w^T = (2, 0)$:
\[
S = \begin{pmatrix} 4 & 4 \\ 6 & 3 \end{pmatrix} - \begin{pmatrix} 3 \\ 5 \end{pmatrix} \begin{pmatrix} 2 & 0 \end{pmatrix}
= \begin{pmatrix} 4 - 6 & 4 \\ 6 - 10 & 3 \end{pmatrix} = \begin{pmatrix} -2 & 4 \\ -4 & 3 \end{pmatrix}.
\]
The second pivot is $-2$, with multiplier $\frac{-4}{-2} = 2$ and Schur complement $3 - 2 \cdot 4 = -5$. Gathering the multipliers in $L$ and the pivots with their rows in $U$:
\[
L = \begin{pmatrix} 1 & 0 & 0 \\ 3 & 1 & 0 \\ 5 & 2 & 1 \end{pmatrix}, \qquad
U = \begin{pmatrix} 1 & 2 & 0 \\ 0 & -2 & 4 \\ 0 & 0 & -5 \end{pmatrix}
\]
(the first column of $L$ stores the multipliers $3$ and $5$ of the first pivot, and the second, the multiplier $2$ of the second; no permutation was needed, $P = I_3$, and one checks directly that $LU = A$).

\textit{Forward substitution} ($Ly = b$): $y_1 = 3$; $y_2 = 7 - 3 \cdot 3 = -2$; $y_3 = 8 - 5 \cdot 3 - 2 \cdot (-2) = -3$.

\textit{Back substitution} ($Ux = y$): $x_3 = \frac{-3}{-5} = \frac{3}{5}$; $x_2 = \frac{1}{-2} \big( {-2} - 4 \cdot \frac{3}{5} \big) = \frac{11}{5}$; $x_1 = 3 - 2 \cdot \frac{11}{5} = -\frac{7}{5}$.

\textit{Verification:} $A x = \big( {-\frac{7}{5}} + \frac{22}{5}, \ {-\frac{21}{5}} + \frac{44}{5} + \frac{12}{5}, \ {-\frac{35}{5}} + \frac{66}{5} + \frac{9}{5} \big)^T = (3, 7, 8)^T = b$.
\end{ejemplo}

\obs Two by-products of the LUP factorization:
\begin{itemize}
    \item \textit{Determinant:} from $\det(P) \det(A) = \det(L) \det(U)$ one solves $\det(A) = (-1)^s \prod_{i=1}^n u_{ii}$, where $s$ is the number of row swaps performed by $P$. It is the practical way of computing determinants, in $\Theta(n^3)$, instead of the $n!$ terms of the permutation formula (Proposition \ref{prop:leibniz}).
    \item \textit{Inverse:} the columns of $A^{-1}$ are the solutions of the $n$ systems $A x_i = e_i$; with the factorization already computed, they cost $\Theta(n^2)$ each: $\Theta(n^3)$ in total. In practice $A^{-1}$ is almost never needed explicitly: solving the system is less costly and numerically more stable.
\end{itemize}

\obs The factorization can be computed \emph{in situ}: since $L$ is unit triangular (its diagonal of ones need not be stored) and $U$ is upper triangular, both fit exactly in the array of $A$, overwriting the multipliers below the diagonal and $U$ from the diagonal upward. It is an example of how the representation in memory (first subsection) shapes the design of the algorithm.

\subsection{QR Factorization and Least Squares}

The QR factorization is the matrix version of the Gram-Schmidt process (Theorem \ref{teo:gram_schmidt}). Applied to the columns $a_1, \dots, a_n$ of $A$, Gram-Schmidt produces orthonormal vectors $q_1, \dots, q_n$ that preserve the spans:
\[
\operatorname{span}\{a_1, \dots, a_k\} = \operatorname{span}\{q_1, \dots, q_k\} \quad \text{for every } k.
\]
That preservation is the central point of what follows: it says that each $a_k$ is a linear combination of the first $k$ orthonormal vectors (it does not need the later ones), so that the coefficient matrix that expresses the columns of $A$ in terms of $q_1, \dots, q_n$ is \emph{upper triangular}. That observation, written as a product of matrices, is precisely the QR factorization.

\begin{definicion}{QR Factorization}{def_qr}
Let $A \in \mathbb{R}^{m \times n}$ have linearly independent columns ($m \ge n$). A \textbf{QR factorization} of $A$ is a writing $A = QR$ with $Q \in \mathbb{R}^{m \times n}$ having orthonormal columns and $R \in \mathbb{R}^{n \times n}$ upper triangular with positive diagonal.
\end{definicion}

\obs Every matrix $Q$ with orthonormal columns satisfies two identities that will be used repeatedly. First: $Q^T Q = I_n$ \emph{even if $Q$ is not square}, since the $(i, j)$ entry of $Q^T Q$ is the product of row $i$ of $Q^T$ (that is, $q_i$) against column $j$ of $Q$: $\langle q_i, q_j \rangle = \delta_{ij}$. Second: when $m > n$, $Q Q^T \in \mathbb{R}^{m \times m}$ is \emph{not} the identity, but the matrix of the orthogonal projection onto the column space of $Q$: for $b \in \mathbb{R}^m$,
\[
Q Q^T b = \sum_{i=1}^{n} \langle b, q_i \rangle \, q_i = P_{\operatorname{Col}(Q)}(b),
\]
by the formula of the projection in an orthonormal basis (Corollary \ref{cor:caso_ortonormal}).

\begin{teorema}{Existence and Uniqueness of the QR}{existencia_qr}
Every matrix $A \in \mathbb{R}^{m \times n}$ with linearly independent columns admits a unique QR factorization.
\end{teorema}

\begin{proof}
\textit{Existence.} Let $a_1, \dots, a_n$ be the columns of $A$, linearly independent by hypothesis, and apply Gram-Schmidt, normalizing:
\[
v_k = a_k - \sum_{j=1}^{k-1} \frac{\langle a_k, v_j \rangle}{\|v_j\|^2} \, v_j, \qquad q_k = \frac{v_k}{\|v_k\|}
\]
(the $v_k$ are nonzero by linear independence, as in Theorem \ref{teo:gram_schmidt}, so that the normalization is well defined). Each term of the sum is the projection of $a_k$ onto the direction of $v_j$, which written with the normalized vector takes the form
\[
\frac{\langle a_k, v_j \rangle}{\|v_j\|^2} \, v_j
= \Big\langle a_k, \tfrac{v_j}{\|v_j\|} \Big\rangle \, \tfrac{v_j}{\|v_j\|}
= \langle a_k, q_j \rangle \, q_j.
\]
Solving for $a_k$ in the definition of $v_k$ and using $v_k = \|v_k\| \, q_k$,
\[
a_k = \underbrace{\langle a_k, q_1 \rangle}_{=: \, r_{1k}} \, q_1 + \dots + \underbrace{\langle a_k, q_{k-1} \rangle}_{=: \, r_{k-1,k}} \, q_{k-1} + \underbrace{\|v_k\|}_{=: \, r_{kk}} \, q_k \, :
\]
each column of $A$ is a linear combination of $q_1, \dots, q_k$ \emph{only}. Define then $R \in \mathbb{R}^{n \times n}$ with the indicated entries $r_{jk}$ and $r_{jk} = 0$ for $j > k$: it is upper triangular, with diagonal $r_{kk} = \|v_k\| > 0$. The equality above says that $a_k = Q \, r_k$, where $r_k$ is column $k$ of $R$ (the product $Q r_k$ is the linear combination of the columns of $Q$ with the coefficients of $r_k$ (Proposition \ref{prop:prop_img_cols}), and the coefficients from the $(k+1)$-th on are zero). Stacking the $n$ equalities column by column, $A = QR$.

\textit{Uniqueness.} The triangular structure forces the columns of $Q$ and of $R$ one by one. Let $A = QR$ be \emph{any} factorization under the conditions of the statement; it is proved by induction on $k$ that the column $q_k$ and the coefficients $r_{1k}, \dots, r_{kk}$ are completely determined by $A$, so that any two factorizations coincide. The starting point is the column-by-column reading of $A = QR$, the same as in the existence part:
\[
a_k = r_{1k} \, q_1 + \dots + r_{k-1,k} \, q_{k-1} + r_{kk} \, q_k
\]
(the coefficients with index greater than $k$ are zero because $R$ is upper triangular).

\textit{Case $k = 1$:} the equality is $a_1 = r_{11} \, q_1$. Taking norms, $\|a_1\| = r_{11} \|q_1\| = r_{11}$ (since $\|q_1\| = 1$ and $r_{11} > 0$): necessarily $r_{11} = \|a_1\|$ and $q_1 = a_1 / \|a_1\|$.

\textit{Inductive step:} suppose that $q_1, \dots, q_{k-1}$ and the first $k-1$ columns of $R$ have already been determined. Taking in the equality above the inner product against $q_i$ with $i < k$, the orthonormality of the columns of $Q$ cancels all the terms but one:
\[
\langle a_k, q_i \rangle = \sum_{j=1}^{k} r_{jk} \, \langle q_j, q_i \rangle = r_{ik},
\]
so that the coefficients $r_{ik}$ with $i < k$ are determined ($q_i$ is, by the inductive hypothesis). With those coefficients known, only one unknown remains in the equality, and it is solved for:
\[
r_{kk} \, q_k = a_k - \sum_{j=1}^{k-1} r_{jk} \, q_j =: z_k,
\]
where $z_k$ is an already determined vector, and nonzero, since $\|z_k\| = r_{kk} \|q_k\| = r_{kk} > 0$. Taking norms, $r_{kk} = \|z_k\|$, and dividing, $q_k = z_k / \|z_k\|$: the whole column $k$, of $Q$ and of $R$, is determined. The induction concludes.
\end{proof}

\obs In the proof of uniqueness, the solved vector $z_k$ is exactly the $v_k$ of Gram-Schmidt, and the formulas obtained ($r_{ik} = \langle a_k, q_i \rangle$, $r_{kk} = \|z_k\|$) coincide with those of the existence part: every QR factorization necessarily reproduces the Gram-Schmidt process.

\obs The columns of $Q$ span the same as those of $A$: each $a_k$ is a combination of $q_1, \dots, q_k$ (existence) and, solving with $R^{-1}$, each $q_k$ is a combination of $a_1, \dots, a_k$. Therefore $\operatorname{Col}(Q) = \operatorname{Col}(A)$, and the previous remark says that $Q Q^T$ is the orthogonal projection onto $\operatorname{Col}(A)$: the QR factorization gathers, in matrix form, an orthonormal basis of the column space and the coefficient matrix $R$ that expresses the original columns in that basis.

\begin{ejemplo}{A QR Factorization}{ej_qr}
Let
\[
A = \begin{pmatrix} 1 & 1 \\ 1 & 0 \\ 0 & 1 \end{pmatrix},
\]
with columns $a_1 = (1,1,0)^T$ and $a_2 = (1,0,1)^T$. Gram-Schmidt gives $v_1 = a_1$, with $\|v_1\| = \sqrt{2}$, and
\[
v_2 = a_2 - \frac{\langle a_2, v_1 \rangle}{\|v_1\|^2} v_1 = (1,0,1)^T - \frac{1}{2}(1,1,0)^T = \Big( \frac{1}{2}, -\frac{1}{2}, 1 \Big)^T,
\qquad \|v_2\| = \frac{\sqrt{6}}{2}.
\]
With $r_{12} = \langle a_2, q_1 \rangle = \frac{1}{\sqrt{2}}$, the factorization is
\[
Q = \begin{pmatrix} \frac{1}{\sqrt{2}} & \frac{1}{\sqrt{6}} \\ \frac{1}{\sqrt{2}} & -\frac{1}{\sqrt{6}} \\ 0 & \frac{2}{\sqrt{6}} \end{pmatrix}, \qquad
R = \begin{pmatrix} \sqrt{2} & \frac{1}{\sqrt{2}} \\ 0 & \frac{\sqrt{6}}{2} \end{pmatrix},
\]
and one checks column by column that $QR = A$: the first is $\sqrt{2} \, q_1 = a_1$ and the second is $\frac{1}{\sqrt{2}} q_1 + \frac{\sqrt{6}}{2} q_2 = \big( \frac{1}{2}, \frac{1}{2}, 0 \big)^T + \big( \frac{1}{2}, -\frac{1}{2}, 1 \big)^T = a_2$.
\end{ejemplo}

When there are more equations than unknowns ($m > n$), for example when fitting a model with few parameters to many noisy data points, the system $Ax = b$ is in general inconsistent: the vectors of the form $Ax$ run only over the subspace $\operatorname{Col}(A)$, of dimension at most $n$, and the measured $b$ does not fall in it. The best that can be done is to minimize the norm of the \textbf{residual} $b - Ax$, and the theory of orthogonal projection has already solved that problem: the point of $\operatorname{Col}(A)$ closest to $b$ is unique and is the projection $P_{\operatorname{Col}(A)}(b)$ (Theorems \ref{teo:mejor_aprox} and \ref{teo:descomposicion_ortogonal}), characterized by the residual being \emph{perpendicular} to the subspace. Imposing that perpendicularity against each column of $A$ produces a system of $n$ linear equations: the \emph{normal equations}.

\begin{teorema}{Normal Equations}{ecuaciones_normales}
Let $A \in \mathbb{R}^{m \times n}$ and $b \in \mathbb{R}^m$. A vector $\hat{x} \in \mathbb{R}^n$ minimizes $\|Ax - b\|$ if and only if
\[
A^T A \, \hat{x} = A^T b.
\]
If the columns of $A$ are linearly independent, the minimizer is unique.
\end{teorema}

\begin{proof}
Minimizing $\|Ax - b\|$ over $x \in \mathbb{R}^n$ is minimizing $\|y - b\|$ over $y \in \operatorname{Col}(A)$, and by Theorem \ref{teo:mejor_aprox} the minimum is attained exactly at $y = P_{\operatorname{Col}(A)}(b)$. By the characterization of the projection (Theorem \ref{teo:descomposicion_ortogonal}), $A\hat{x} = P_{\operatorname{Col}(A)}(b)$ if and only if the residual $b - A\hat{x}$ belongs to $\operatorname{Col}(A)^\perp$; since the columns $a_1, \dots, a_n$ generate $\operatorname{Col}(A)$, that is equivalent to the $n$ conditions $\langle a_i, \, b - A\hat{x} \rangle = 0$. Row $i$ of $A^T$ is $a_i^T$, so that the $n$ conditions, stacked, are exactly
\[
A^T (b - A\hat{x}) = 0 \iff A^T A \, \hat{x} = A^T b.
\]

If the columns of $A$ are linearly independent, $\ker(A^T A) = \ker(A) = \{0\}$ (as in the proof of Theorem \ref{teo:svd}), so that $A^T A$ is invertible and $\hat{x}$ is unique.
\end{proof}

\obs With the factorization $A = QR$ the problem is solved without forming $A^T A$, through a geometric derivation. Since $\operatorname{Col}(A) = \operatorname{Col}(Q)$, the projection onto $\operatorname{Col}(A)$ is $Q Q^T$ (previous remarks), and the condition for the minimizer is rewritten
\[
A \hat{x} = P_{\operatorname{Col}(A)}(b)
\iff Q R \, \hat{x} = Q Q^T b
\iff R \, \hat{x} = Q^T b,
\]
where the last equivalence is obtained by multiplying on the left by $Q^T$ (with $Q^T Q = I_n$) in one direction, and by $Q$ in the other. The system $R \hat{x} = Q^T b$ is triangular and is solved by back substitution. The two factors have a direct interpretation: $Q^T b$ contains the coordinates of the projection of $b$ in the orthonormal basis $q_1, \dots, q_n$, and solving with $R$ translates those coordinates into coefficients on the original columns of $A$. The same system is obtained from the normal equations: $A^T A = R^T Q^T Q R = R^T R$ and $A^T b = R^T Q^T b$, and cancelling $R^T$, which is invertible, one gets $R \hat{x} = Q^T b$. Avoiding the passage through $A^T A$ matters numerically, because that product amplifies rounding errors: it is the same reason why PCA is computed via the SVD and not by diagonalizing the covariance matrix.

\begin{ejemplo}{Least Squares}{ej_minimos_cuadrados}
With the matrix $A$ of Example \ref{ej:ej_qr} and $b = (1, 1, 1)^T$, the system $Ax = b$ has no solution. The normal equations give
\[
A^T A = \begin{pmatrix} 2 & 1 \\ 1 & 2 \end{pmatrix}, \qquad A^T b = \begin{pmatrix} 2 \\ 2 \end{pmatrix}
\qquad \Longrightarrow \qquad \hat{x} = \Big( \frac{2}{3}, \frac{2}{3} \Big)^T.
\]
By the QR route: $Q^T b = \big( \sqrt{2}, \frac{2}{\sqrt{6}} \big)^T$, and back substitution in $R \hat{x} = Q^T b$ gives $\hat{x}_2 = \frac{2/\sqrt{6}}{\sqrt{6}/2} = \frac{2}{3}$ and $\hat{x}_1 = \frac{1}{\sqrt{2}} \big( \sqrt{2} - \frac{1}{\sqrt{2}} \cdot \frac{2}{3} \big) = \frac{2}{3}$: the same result. The residual is
\[
b - A\hat{x} = (1,1,1)^T - \Big( \frac{4}{3}, \frac{2}{3}, \frac{2}{3} \Big)^T = \Big( {-\frac{1}{3}}, \frac{1}{3}, \frac{1}{3} \Big)^T,
\]
orthogonal to both columns of $A$ ($-\frac{1}{3} + \frac{1}{3} = 0$ in both cases), as the projection requires.
\end{ejemplo}

\obs The classical use case is \emph{polynomial fitting}: given $m$ points $(t_i, b_i)$ and a degree $d$ with $d + 1 \le m$, finding the polynomial $p(t) = x_0 + x_1 t + \dots + x_d t^d$ that minimizes $\sum_i (p(t_i) - b_i)^2$ is exactly the least squares problem with the Vandermonde matrix $a_{ij} = t_i^{\, j-1}$, of size $m \times (d+1)$.

\subsection{The Power Method}

The sections on diagonalization, SVD and PCA assume that eigenvalues and eigenvectors can be computed, but they do not say how. The direct route (finding the roots of the characteristic polynomial) is not numerically practicable: for $n \ge 5$ there is no closed formula for the roots (Abel-Ruffini theorem, cited without proof) and, worse, the roots of a polynomial are extremely sensitive to perturbations of its coefficients. The real methods are iterative, and the simplest one (the conceptual core of the others) is the \textbf{power method}: multiply repeatedly by $A$ and normalize. Given an initial vector $x^{(0)}$ of norm $1$, iterate
\[
x^{(t+1)} = \frac{A \, x^{(t)}}{\| A \, x^{(t)} \|},
\]
and the eigenvalue is estimated with the \textbf{Rayleigh quotient}
\[
\lambda^{(t)} = \langle x^{(t)}, A \, x^{(t)} \rangle,
\]
the value of the quadratic form of $A$ on the unit sphere (Theorem \ref{teo:extremos_esfera}).

\begin{teorema}{Convergence of the Power Method}{potencia}
Let $A \in \mathbb{R}^{n \times n}$ be symmetric with eigenvalues $\lambda_1 > |\lambda_2| \ge \dots \ge |\lambda_n|$ and an orthonormal basis of eigenvectors $v_1, \dots, v_n$ (Corollary \ref{cor:diag_ortogonal}). If $c_1 = \langle x^{(0)}, v_1 \rangle \neq 0$, then the angle $\theta_t$ between $x^{(t)}$ and $v_1$ satisfies
\[
\tan \theta_t \le \Big( \frac{|\lambda_2|}{\lambda_1} \Big)^{\! t} \tan \theta_0,
\]
so that $x^{(t)}$ converges to $v_1$ or to $-v_1$ (according to the sign of $c_1$) with geometric ratio $|\lambda_2| / \lambda_1$, and $\lambda^{(t)} \to \lambda_1$ with ratio $(|\lambda_2| / \lambda_1)^2$.
\end{teorema}

\begin{proof}
Since each normalization divides by a positive scalar, $x^{(t)}$ is the vector $y^{(t)} = A^t x^{(0)}$ normalized. Write $x^{(0)} = \sum_i c_i v_i$ in the orthonormal basis of eigenvectors; since $A^t v_i = \lambda_i^t v_i$,
\[
y^{(t)} = \sum_{i=1}^{n} c_i \, \lambda_i^t \, v_i.
\]
The component of $y^{(t)}$ along $v_1$ is $c_1 \lambda_1^t \neq 0$, and the component orthogonal to $v_1$ has norm $\big( \sum_{i \ge 2} c_i^2 \lambda_i^{2t} \big)^{1/2}$. The tangent of the angle between $y^{(t)}$ and $v_1$ is the quotient of both:
\[
\tan \theta_t = \frac{\big( \sum_{i \ge 2} c_i^2 \, \lambda_i^{2t} \big)^{1/2}}{|c_1| \, \lambda_1^t}
\le \Big( \frac{|\lambda_2|}{\lambda_1} \Big)^{\! t} \, \frac{\big( \sum_{i \ge 2} c_i^2 \big)^{1/2}}{|c_1|}
= \Big( \frac{|\lambda_2|}{\lambda_1} \Big)^{\! t} \tan \theta_0,
\]
using $|\lambda_i|^{2t} \le |\lambda_2|^{2t}$ for $i \ge 2$. Since $|\lambda_2| / \lambda_1 < 1$, the angle tends to $0$; and the component of $x^{(t)}$ along $v_1$, which is $c_1 \lambda_1^t / \|y^{(t)}\|$, keeps the sign of $c_1$, so that the limit is $\operatorname{sign}(c_1) \, v_1$.

For the Rayleigh quotient, write $x^{(t)} = \sum_i w_i v_i$, with $\sum_i w_i^2 = 1$ and $w_1^2 = \cos^2 \theta_t$. By orthonormality, $\lambda^{(t)} = \sum_i \lambda_i w_i^2$, and then
\[
\lambda_1 - \lambda^{(t)} = \sum_{i \ge 2} (\lambda_1 - \lambda_i) \, w_i^2 \le (\lambda_1 - \lambda_n) \sin^2 \theta_t,
\]
which decays with the \emph{square} of the ratio, $(|\lambda_2| / \lambda_1)^{2t}$, since $\sin \theta_t \le \tan \theta_t$.
\end{proof}

\begin{ejemplo}{The Power Method on a $2 \times 2$ Matrix}{ej_potencia}
Let $A = \begin{pmatrix} 2 & 1 \\ 1 & 2 \end{pmatrix}$, the matrix $A^T A$ of Example \ref{ej:ej_minimos_cuadrados}, with eigenvalues $\lambda_1 = 3$ and $\lambda_2 = 1$ and eigenvectors $(1, 1)^T$ and $(1, -1)^T$. With $x^{(0)} = (1, 0)^T = \frac{1}{2} \big[ (1,1)^T + (1,-1)^T \big]$,
\[
A^t x^{(0)} = \frac{1}{2} \Big( 3^t \, (1, 1)^T + (1, -1)^T \Big) = \Big( \frac{3^t + 1}{2}, \ \frac{3^t - 1}{2} \Big)^{T}:
\]
unnormalized, the iterates are $(1,0)$, $(2,1)$, $(5,4)$, $(14,13)$, $(41,40)$, \dots, whose direction tends to that of $(1,1)$ with ratio $|\lambda_2| / \lambda_1 = \frac{1}{3}$, as the theorem predicts. The Rayleigh quotients (for example, for $(2,1)$: $\frac{(2,1) \cdot (5,4)}{2^2 + 1^2} = \frac{14}{5}$) are
\[
\lambda^{(0)} = 2, \qquad \lambda^{(1)} = \frac{14}{5} = 2.8, \qquad \lambda^{(2)} = \frac{122}{41} \approx 2.9756, \qquad \lambda^{(3)} = \frac{1094}{365} \approx 2.9973,
\]
and the errors $1, \ 0.2, \ 0.024, \ 0.0027$ are indeed reduced by a factor close to $(|\lambda_2| / \lambda_1)^2 = \frac{1}{9}$ per step.
\end{ejemplo}

\obs For the following components of a symmetric matrix, \emph{deflation} is used: once the pair $(\lambda_1, v_1)$ is found, one moves to the matrix
\[
A' = A - \lambda_1 \, v_1 v_1^T.
\]
Expanding the orthogonal diagonalization $A = P D P^T$ column by column one obtains $A = \sum_i \lambda_i \, v_i v_i^T$, so that $A'$ has the same eigenvectors as $A$ with the eigenvalue $\lambda_1$ replaced by $0$: the power method on $A'$ produces $(\lambda_2, v_2)$, and so on. Applied to the covariance matrix, this scheme computes the first $k$ principal components of PCA (Section 9) without diagonalizing the full matrix.

\obs Three practical notes. First: a random $x^{(0)}$ satisfies $c_1 \neq 0$ with probability $1$. Second: each iteration costs a matrix-vector product, $\Theta(n^2)$, or proportional to the nonzero entries if the matrix is sparse; the celebrated example is PageRank, in essence a power method on the link matrix of the web. Third: numerical libraries (LAPACK) do not use this method for the full problem, but the shifted QR algorithm, a more elaborate iteration that obtains all the eigenvalues at once and that is cited without proof; the power method remains the choice when only the first components of a huge sparse matrix are needed.

\newpage

\section{Machine Learning Algorithms}

\subsection{Clustering: the k-means Problem}

Given $n$ unlabeled points of $\mathbb{R}^d$, one wants to group them into $k$ groups of points close to each other. As a measure of closeness, the \textbf{squared distance}
\[
\delta(x, y) = \|x - y\|^2 = \sum_{a=1}^{d} (x_a - y_a)^2,
\]
the square of the Euclidean distance, is used. It is not a metric (it does not satisfy the triangle inequality), but it is the conventional choice because of its algebraic convenience, as the theorems of this subsection show.

\begin{definicion}{The k-means Problem}{kmeans}
Let $S \subseteq \mathbb{R}^d$ be a set of $n$ points and let $k$ be a positive integer. A \textbf{k-clustering} of $S$ is a partition $\langle S^{(1)}, \dots, S^{(k)} \rangle$ of $S$ together with a sequence of \textbf{centers} $C = \langle c^{(1)}, \dots, c^{(k)} \rangle$ in $\mathbb{R}^d$. Its cost is the sum of the squared distances from each point to the center of its group,
\[
f(S, C) = \sum_{\ell=1}^{k} \sum_{x \in S^{(\ell)}} \delta(x, c^{(\ell)}),
\]
and the \textbf{k-means problem} consists in finding the k-clustering of minimum cost.
\end{definicion}

The following two theorems solve separately the two halves of the problem: which is the best center for a given group, and which is the best grouping for given centers.

\begin{teorema}{The Centroid as the Optimal Center}{centroide}
Let $S^{(\ell)}$ be a nonempty cluster. The unique center $c \in \mathbb{R}^d$ that minimizes $\sum_{x \in S^{(\ell)}} \delta(x, c)$ is its \textbf{centroid} (or mean)
\[
c^{(\ell)} = \frac{1}{|S^{(\ell)}|} \sum_{x \in S^{(\ell)}} x.
\]
\end{teorema}

\begin{proof}
Let $q = |S^{(\ell)}|$. The sum splits by coordinates:
\[
\sum_{x \in S^{(\ell)}} \delta(x, c) = \sum_{a=1}^{d} \sum_{x \in S^{(\ell)}} (x_a - c_a)^2
= \sum_{a=1}^{d} \Bigg( \sum_{x \in S^{(\ell)}} x_a^2 - 2 \Big( \sum_{x \in S^{(\ell)}} x_a \Big) c_a + q \, c_a^2 \Bigg).
\]
For each $a$, the term in parentheses is a quadratic in $c_a$ with leading coefficient $q > 0$; completing the square, with $\mu_a = \frac{1}{q} \sum_{x \in S^{(\ell)}} x_a$,
\[
\sum_{x \in S^{(\ell)}} x_a^2 - 2 q \mu_a c_a + q c_a^2 = q (c_a - \mu_a)^2 + \Big( \sum_{x \in S^{(\ell)}} x_a^2 - q \mu_a^2 \Big),
\]
which is minimized, uniquely, at $c_a = \mu_a$. Since this holds for each coordinate separately, the unique minimizer is $c = (\mu_1, \dots, \mu_d)$, the centroid.
\end{proof}

\begin{teorema}{Assignment to the Nearest Center}{asignacion_cercana}
Given $S$ and a sequence of centers $\langle c^{(1)}, \dots, c^{(k)} \rangle$, a k-clustering minimizes $f(S, C)$ if and only if it assigns each point $x \in S$ to a cluster $S^{(\ell)}$ whose center minimizes $\delta(x, c^{(\ell)})$.
\end{teorema}

\begin{proof}
Each point $x \in S$ contributes to the sum $f(S, C)$ exactly one term, $\delta(x, c^{(\ell)})$ for the $\ell$ of its cluster, and the contributions are independent of each other: the sum is minimal if and only if each contribution is, that is, if each point is assigned to a nearest center (in squared or Euclidean distance, which induce the same order).
\end{proof}

\textbf{Lloyd's procedure} alternates the two partial optimizations given by the theorems:
\begin{enumerate}
    \item \textit{Initialize:} choose $k$ centers (for example, $k$ points of $S$ at random).
    \item \textit{Assign:} form the clustering by assigning each point of $S$ to a nearest center (Theorem \ref{teo:asignacion_cercana}); ties are broken in a fixed way, for example by the smallest index.
    \item \textit{Recompute:} replace each center by the centroid of its cluster (Theorem \ref{teo:centroide}).
    \item Repeat steps 2 and 3 until an iteration does not change the clustering.
\end{enumerate}
Each iteration costs $O(dkn)$ operations ($nk$ distances of $d$ coordinates, plus the centroids); with $T$ iterations, the total is $O(T d k n)$.

\begin{proposicion}{Monotonicity and Termination of Lloyd}{lloyd_termina}
In Lloyd's procedure, the cost $f(S, C)$ never increases, and the procedure terminates after a finite number of iterations.
\end{proposicion}

\begin{proof}
Monotonicity is a direct consequence of the two theorems: the assignment step does not increase $f$, because with the centers fixed the assignment to the nearest one is optimal among all possible assignments (Theorem \ref{teo:asignacion_cercana}), and neither does the recomputation step, because with the clustering fixed the centroid is the optimal center of each group (Theorem \ref{teo:centroide}).

For termination, the two cases of the recomputation step are examined. If it modifies some center, then $f$ \emph{decreases strictly}: the centroid is the \emph{unique} minimizer (Theorem \ref{teo:centroide}), so that the replaced center, different from it, had strictly larger cost. If it modifies none, the centers remained fixed, the next assignment (deterministic by the fixed tie-breaking rule) reproduces the previous clustering, and the procedure stops.

Consequently, every iteration that is not the last one strictly decreases $f$. Moreover, after a recomputation the value of $f$ is determined by the partition (the centers are its centroids), and the partitions of $S$ into at most $k$ groups are finitely many: the values that $f$ can take form a finite set, and a strictly decreasing sequence inside a finite set is necessarily finite. The number of iterations is, therefore, finite.
\end{proof}

\obs The procedure ends at a \emph{local} minimum: no reassignment or recomputation improves it. It does not guarantee the global optimum (the result depends on the initialization, and in practice it is run several times with different initializations), and finding the global optimum is computationally hard in general.

\obs Lloyd's procedure illustrates the scheme common to a large part of machine learning: define a parameterized hypothesis space (here, the sequences of $k$ centers), define a function that measures how badly each hypothesis fits the data (here, $f(S, C)$) and minimize it, at least locally, with an optimization procedure. The gradient descent of the next subsection, and the Principal Component Analysis of Section 9, follow the same scheme.

\begin{ejemplo}{Lloyd on the Line}{ej_lloyd}
Let $S = \{0, 2, 5, 6\} \subset \mathbb{R}$ (that is, $d = 1$) with $k = 2$, and take as initial centers $c^{(1)} = 0$ and $c^{(2)} = 2$.

\textit{Iteration 1.} Assignment: $0 \to c^{(1)}$; $2, 5, 6 \to c^{(2)}$ (for example, $\delta(5, 0) = 25 > 9 = \delta(5, 2)$). Recomputation: $c^{(1)} = 0$, $c^{(2)} = \frac{2+5+6}{3} = \frac{13}{3}$. Cost:
\[
f = 0 + \Big( 2 - \tfrac{13}{3} \Big)^2 + \Big( 5 - \tfrac{13}{3} \Big)^2 + \Big( 6 - \tfrac{13}{3} \Big)^2 = \tfrac{49}{9} + \tfrac{4}{9} + \tfrac{25}{9} = \tfrac{26}{3}.
\]

\textit{Iteration 2.} Assignment with $c^{(1)} = 0$, $c^{(2)} = \frac{13}{3}$: the point $2$ \emph{changes cluster}, since $\delta(2, 0) = 4 < \frac{49}{9} = \delta \big( 2, \frac{13}{3} \big)$; this leaves $S^{(1)} = \{0, 2\}$ and $S^{(2)} = \{5, 6\}$. Recomputation: $c^{(1)} = 1$, $c^{(2)} = \frac{11}{2}$. Cost: $f = 1 + 1 + \frac{1}{4} + \frac{1}{4} = \frac{5}{2}$.

\textit{Iteration 3.} The assignment with $c^{(1)} = 1$, $c^{(2)} = \frac{11}{2}$ changes nothing ($\delta(2,1) = 1 < \frac{49}{4}$, $\delta(5, \frac{11}{2}) = \frac{1}{4} < 16$): the procedure ends with cost $\frac{5}{2}$, which in this case is the global optimum (the other partitions into two intervals, $\{0\} \cup \{2,5,6\}$ and $\{0,2,5\} \cup \{6\}$, cost $\frac{26}{3}$ and $\frac{114}{9}$).
\end{ejemplo}

\obs Image compression by \emph{vector quantization} is a direct application: each pixel is a point of $\mathbb{R}^3$ (its RGB intensities), and running Lloyd with $k = 256$ produces a palette of $256$ representative colors; replacing each pixel by its nearest center, the image goes from $24$ to $8$ bits per pixel.

\subsection{Gradient Descent}

Proposition \ref{prop:prop_max_descenso} identified the direction in which a differentiable function decreases fastest: $-\nabla f(x)$. \textbf{Gradient descent} turns it into a minimization algorithm: from an initial point $x^{(0)} \in \mathbb{R}^n$, each iteration moves in that direction, with a length regulated by a fixed parameter $\eta > 0$ called the \textbf{step size} (or \emph{learning rate}):
\[
x^{(t+1)} = x^{(t)} - \eta \, \nabla f(x^{(t)}), \qquad t = 0, 1, \dots, T - 1.
\]
After $T$ iterations, the algorithm returns the \emph{average} $\bar{x} = \frac{1}{T} \big( x^{(0)} + \dots + x^{(T-1)} \big)$; the average is returned, and not the last iterate, because it is the variant for which the convergence proof that follows obtains a guarantee. Without hypotheses on $f$, a local method of this kind can only aspire to a local minimum; the hypothesis that turns local into global is convexity:

\begin{definicion}{Convex Function}{convexidad}
A function $f: \mathbb{R}^n \to \mathbb{R}$ is \textbf{convex} if for all $x, y \in \mathbb{R}^n$ and every $\lambda \in [0, 1]$,
\[
f(\lambda x + (1 - \lambda) y) \le \lambda f(x) + (1 - \lambda) f(y).
\]
\end{definicion}

\obs For a convex function, every local minimum is global: if $x$ is a local minimum and there were $y$ with $f(y) < f(x)$, the points $\lambda y + (1-\lambda) x$ with $\lambda \to 0^+$ would be arbitrarily close to $x$ with value at most $\lambda f(y) + (1-\lambda) f(x) < f(x)$, against local minimality.

\begin{lema}\label{lem_tangente}
Let $f: \mathbb{R}^n \to \mathbb{R}$ be convex and differentiable. For all $x, y \in \mathbb{R}^n$,
\[
f(y) \ge f(x) + \langle \nabla f(x), y - x \rangle,
\]
that is, the function lies above its tangent planes.
\end{lema}

\begin{proof}
By convexity, for $\lambda \in (0, 1]$,
\[
f(x + \lambda (y - x)) = f((1-\lambda) x + \lambda y) \le (1 - \lambda) f(x) + \lambda f(y),
\]
whence
\[
\frac{f(x + \lambda (y - x)) - f(x)}{\lambda} \le f(y) - f(x).
\]
As $\lambda \to 0^+$, the left-hand side tends to the directional derivative of $f$ at $x$ along $y - x$, which by differentiability is $\langle \nabla f(x), y - x \rangle$ (Definition \ref{def:def_diferencial} and Theorem \ref{teo:teo_gradiente_riesz}). Rearranging yields the inequality.
\end{proof}

\begin{lema}\label{lem_jensen}
Let $f: \mathbb{R}^n \to \mathbb{R}$ be convex. For all $x^{(0)}, \dots, x^{(T-1)} \in \mathbb{R}^n$,
\[
f \Bigg( \frac{x^{(0)} + \dots + x^{(T-1)}}{T} \Bigg) \le \frac{f(x^{(0)}) + \dots + f(x^{(T-1)})}{T}.
\]
\end{lema}

\begin{proof}
By induction on $T$; the case $T = 1$ is an equality. For the step, the average of $T$ points is written as a convex combination of the last point and the average of the first $T - 1$:
\[
\frac{x^{(0)} + \dots + x^{(T-1)}}{T} = \frac{1}{T} \, x^{(T-1)} + \frac{T-1}{T} \cdot \frac{x^{(0)} + \dots + x^{(T-2)}}{T-1},
\]
and applying Definition \ref{def:convexidad} with $\lambda = \frac{1}{T}$ and then the inductive hypothesis yields the bound.
\end{proof}

\begin{teorema}{Convergence of Gradient Descent}{convergencia_gd}
Let $f: \mathbb{R}^n \to \mathbb{R}$ be convex and differentiable with minimizer $x^*$. Let $R = \|x^{(0)} - x^*\|$ and $L$ an upper bound of $\|\nabla f(x^{(t)})\|$ over the iterates $t = 0, \dots, T-1$. If the step size is $\eta = \frac{R}{L \sqrt{T}}$, then the average $\bar{x}$ returned by gradient descent satisfies
\[
f(\bar{x}) - f(x^*) \le \frac{R L}{\sqrt{T}}.
\]
\end{teorema}

\begin{proof}
How the distance to the minimizer evolves is studied. Expanding the squared norm of $x^{(t+1)} - x^* = (x^{(t)} - x^*) - \eta \nabla f(x^{(t)})$,
\[
\|x^{(t+1)} - x^*\|^2 = \|x^{(t)} - x^*\|^2 - 2 \eta \, \langle \nabla f(x^{(t)}), \, x^{(t)} - x^* \rangle + \eta^2 \, \|\nabla f(x^{(t)})\|^2.
\]
Solving that equality for the inner product,
\[
\langle \nabla f(x^{(t)}), \, x^{(t)} - x^* \rangle
= \frac{\|x^{(t)} - x^*\|^2 - \|x^{(t+1)} - x^*\|^2}{2 \eta} + \frac{\eta}{2} \, \|\nabla f(x^{(t)})\|^2.
\]
On the other hand, Lemma \ref{lem_tangente} with $x = x^{(t)}$ and $y = x^*$ gives $f(x^*) \ge f(x^{(t)}) + \langle \nabla f(x^{(t)}), x^* - x^{(t)} \rangle$, that is, $f(x^{(t)}) - f(x^*) \le \langle \nabla f(x^{(t)}), x^{(t)} - x^* \rangle$. Combining both and bounding $\|\nabla f(x^{(t)})\| \le L$,
\[
f(x^{(t)}) - f(x^*) \le \frac{\|x^{(t)} - x^*\|^2 - \|x^{(t+1)} - x^*\|^2}{2 \eta} + \frac{\eta L^2}{2}.
\]
Summing over $t = 0, \dots, T-1$, the distances in the first fraction cancel in a chain (telescoping), only the first minus the last remains, and the term $-\|x^{(T)} - x^*\|^2 \le 0$ is discarded:
\[
\sum_{t=0}^{T-1} \big( f(x^{(t)}) - f(x^*) \big) \le \frac{\|x^{(0)} - x^*\|^2}{2 \eta} + \frac{T \eta L^2}{2} = \frac{R^2}{2 \eta} + \frac{T \eta L^2}{2}.
\]
Dividing by $T$ and applying Lemma \ref{lem_jensen} to the left-hand side,
\[
f(\bar{x}) - f(x^*) \le \frac{1}{T} \sum_{t=0}^{T-1} \big( f(x^{(t)}) - f(x^*) \big) \le \frac{R^2}{2 \eta T} + \frac{\eta L^2}{2},
\]
and with $\eta = \frac{R}{L \sqrt{T}}$ the two summands equal $\frac{RL}{2\sqrt{T}}$ each.
\end{proof}

\obs Three remarks on the bound. First: to achieve error $\varepsilon$, $T = (RL/\varepsilon)^2$ iterations suffice, each with the cost of evaluating a gradient: the error decays as $O(1/\sqrt{T})$. Second: the optimal step size balances the two summands of the proof; a step that is too large makes the term $\eta L^2 / 2$ dominate (the algorithm overshoots the minimum at each iteration), and one that is too small, the term $R^2 / (2\eta T)$ (the progress per iteration is insufficient). Third: the proof does not need $f$ to decrease at each step; with additional hypotheses (Lipschitz gradient, strong convexity) better rates are obtained, $O(1/T)$ and geometric, which are cited without proof.

\begin{ejemplo}{Gradient Descent and Least Squares}{ej_gd}
The squared error function $f(x) = \frac{1}{2} \|Ax - b\|^2$ is convex: given $x, y$ and $\lambda \in [0,1]$, since $A(\lambda x + (1-\lambda) y) - b = \lambda (Ax - b) + (1-\lambda)(Ay - b)$, the triangle inequality and the convexity of $t \mapsto t^2$ on $[0, \infty)$ give
\[
f(\lambda x + (1-\lambda) y) \le \tfrac{1}{2} \big( \lambda \|Ax - b\| + (1-\lambda) \|Ay - b\| \big)^2 \le \lambda f(x) + (1-\lambda) f(y).
\]
Its gradient is $\nabla f(x) = A^T A x - A^T b$ (as in the section on quadratic forms, using that $A^T A$ is symmetric), so that $\nabla f(x) = 0$ are exactly the normal equations (Theorem \ref{teo:ecuaciones_normales}): gradient descent converges to the least squares solution, and it is the practical alternative to QR when $A$ is too large to factorize.

With the matrix $A$ and the vector $b$ of Example \ref{ej:ej_minimos_cuadrados} ($\hat{x} = (\frac{2}{3}, \frac{2}{3})^T$), starting from $x^{(0)} = (0,0)^T$ with step size $\eta = \frac{1}{4}$:
\begin{gather*}
\nabla f(x) = \begin{pmatrix} 2 & 1 \\ 1 & 2 \end{pmatrix} x - \begin{pmatrix} 2 \\ 2 \end{pmatrix};
\\
x^{(1)} = \tfrac{1}{4} \begin{pmatrix} 2 \\ 2 \end{pmatrix} = \begin{pmatrix} 1/2 \\ 1/2 \end{pmatrix}, \qquad
x^{(2)} = \begin{pmatrix} 5/8 \\ 5/8 \end{pmatrix}, \qquad
x^{(3)} = \begin{pmatrix} 21/32 \\ 21/32 \end{pmatrix}, \; \dots
\end{gather*}
(for example, $\nabla f(x^{(1)}) = (\frac{3}{2}, \frac{3}{2})^T - (2,2)^T = -(\frac{1}{2}, \frac{1}{2})^T$ and $x^{(2)} = x^{(1)} + \frac{1}{4} (\frac{1}{2}, \frac{1}{2})^T$). The values of $f$ decrease toward $f(\hat{x}) = \frac{1}{6}$:
\[
f(x^{(0)}) = \tfrac{3}{2}, \qquad f(x^{(1)}) = \tfrac{1}{4}, \qquad f(x^{(2)}) = \tfrac{11}{64}, \qquad f(x^{(3)}) = \tfrac{171}{1024} \approx 0.167.
\]
The convergence here is much faster than the general guarantee of Theorem \ref{teo:convergencia_gd}: the iterates live on the line with direction $(1,1)^T$, an eigenvector of $A^T A$ with eigenvalue $3$, and the error contracts by the factor $|1 - \eta \cdot 3| = \frac{1}{4}$ per step: $x^{(t)} = \frac{2}{3} - \frac{2}{3} \cdot 4^{-t}$ in each coordinate. It is a preview of the role of eigenvalues in the speed of gradient descent.
\end{ejemplo}

\obs In modern machine learning, $f$ is usually the error of a model over millions of data points and $x$ its parameters; evaluating the full gradient is costly, and it is replaced by the gradient over a random sample of data (\emph{stochastic gradient descent}, SGD). Moreover, the error functions of deep models are not convex, so that the guarantees of this section do not apply literally; even so, gradient descent and its variants constitute the standard training procedure.

\end{document}
